\documentclass[11pt]{article}

\usepackage[margin=0.9in]{geometry}
\usepackage[T1]{fontenc}
\usepackage{lmodern}
\usepackage{microtype}
\usepackage{amsmath,amssymb,amsthm}
\usepackage{booktabs}
\usepackage{array}
\usepackage{longtable}
\usepackage{algorithm}
\usepackage{algpseudocode}
\usepackage[dvipsnames]{xcolor}
\usepackage{graphicx}
% hyperref is loaded BEFORE tikz on purpose.  hyperref pulls in url, which
% defines \path; tikz defines \path too, and whichever loads last wins.  TikZ
% needs its own (\draw and \node are built on it), so we capture url's verbatim
% behaviour as \filepath first and then let tikz take \path back.
\usepackage{hyperref}
\DeclareUrlCommand\filepath{\urlstyle{tt}}
\usepackage{tikz}
\usetikzlibrary{arrows.meta,positioning,calc}
\usepackage{enumitem}
\usepackage{fancyhdr}
\usepackage[font=small,labelfont=bf,skip=6pt]{caption}

\definecolor{fishgreen}{HTML}{0B5D46}
\definecolor{fishlight}{HTML}{EAF4EF}
\definecolor{fishgray}{HTML}{4B5954}
\hypersetup{colorlinks=true, linkcolor=fishgreen, citecolor=fishgreen, urlcolor=fishgreen}
\setlist{nosep,leftmargin=*}
\setlength{\parindent}{1.2em}
\setlength{\tabcolsep}{4.5pt}
\setlength{\parskip}{0pt plus 1pt}
\pagestyle{fancy}
\fancyhf{}
\lhead{\small FishBot v0.5}
\rhead{\small Live-ask gating, feasible declaration, and a retracted freeze claim}
\cfoot{\thepage}

\newtheorem{theorem}{Theorem}
\newtheorem{proposition}[theorem]{Proposition}
\newtheorem{corollary}[theorem]{Corollary}
\theoremstyle{definition}
\newtheorem{definition}[theorem]{Definition}
\newtheorem{observation}[theorem]{Observation}

% Consistent names for the policies compared.  v0.4-Fast is the policy shipped
% and evaluated by the v0.4 study; v0.4-Block is the same fitted policy with the
% exact reference belief substituted.  v0.5 is the policy of this report.
% These names appear inside captions and section titles, which are moving
% arguments, so the commands must be robust; and they end up in PDF bookmarks,
% so \texorpdfstring supplies a plain spelling for the bookmark text.
\DeclareRobustCommand{\vfive}{\texorpdfstring{\mbox{FishBot v0.5}}{FishBot v0.5}}
\DeclareRobustCommand{\vfour}{\texorpdfstring{\mbox{FishBot v0.4}}{FishBot v0.4}}
\DeclareRobustCommand{\vthree}{\texorpdfstring{\mbox{FishBot v0.3}}{FishBot v0.3}}
\DeclareRobustCommand{\vtwo}{\texorpdfstring{\mbox{FishBot v0.2}}{FishBot v0.2}}
\DeclareRobustCommand{\vfast}{\texorpdfstring{\mbox{v0.4-Fast}}{v0.4-Fast}}
\DeclareRobustCommand{\vblock}{\texorpdfstring{\mbox{v0.4-Block}}{v0.4-Block}}

\title{\textbf{FishBot v0.5: A Self-Inflicted Deadlock, Its Measurement,\\
and the Repair of the Ask Policy That Caused It}\\[5pt]
\large Live-ask gating, capacity-feasible declaration, and the retraction\\
of an information-freeze claim in six-player Canadian Fish}
\author{Dylan Nguyen\\\small FishLab Research Project}
\date{%
  \small Technical report --- \today\\[2pt]
  \small Repository: \filepath{https://github.com/dylann29/fish-optimization}\\[1pt]
  \small The diagnosis of \S\ref{sec:diagnosis} and the corrections of
  \S\ref{sec:corrections} were produced at commit \texttt{fe21e19}; the \vfive{}
  policy of \S\ref{sec:mechanisms} is built on top of that commit. Diagnosis
  artifacts are in \filepath{research/v05/results/}.%
}

%% --- begin numbers_v05 ---
%% numbers_v05.tex -- every number quoted in the v0.5 manuscript.
%%
%% PLACEHOLDER.  This file is hand-written for now so that fishbot_v05.tex
%% compiles before the v0.5 mode of engine/build_tables.py exists.  Each value
%% below is the figure carried by the artifact named in the comment above it;
%% the generator will replace this file wholesale with \newcommand definitions
%% read straight from those artifacts.  \providecommand is used so that a
%% partially generated file can be \input first without a clash.
%%
%% Discipline: no number appears as a literal anywhere in sections_v05/.  If a
%% section needs a figure, it gets a macro here and the macro cites its source.

%% ---------------------------------------------------------------------------
%% research/v05/results/P0-v04-pathology.md -- v0.4 mirror, 600 games, seed 31
%% ---------------------------------------------------------------------------
\providecommand{\numMirrorGames}{300}
\providecommand{\numMirrorEvents}{143.6}
\providecommand{\numMirrorHit}{34.2}
\providecommand{\numDeadAsk}{39.0}
\providecommand{\numRepeatAsk}{40.0}
\providecommand{\numDeadRunLongest}{286}
\providecommand{\numDeadRunGames}{34.3}
\providecommand{\numStarved}{0.20}
\providecommand{\numOwnLocked}{16.5}
\providecommand{\numDeclWrong}{10.4}
\providecommand{\numPostHorizonDecl}{384}
\providecommand{\numPostHorizonWrong}{58.6}
\providecommand{\numForcedWrongRate}{100}

%% research/v05/results/P0-v04-pathology.md -- v0.4 vs v0.3, 600 games, seed 90210
\providecommand{\numVthreeEvents}{101.7}
\providecommand{\numVthreeHit}{54.9}
\providecommand{\numVthreeDeadAsk}{2.82}
\providecommand{\numVthreeRepeat}{5.29}
\providecommand{\numVthreeDeadRun}{5}
\providecommand{\numVthreeDeclWrong}{8.96}
\providecommand{\numVthreeDeadRunGames}{0}
\providecommand{\numVthreeStarved}{0.03}
\providecommand{\numVthreeOwnLocked}{11.7}

%% research/v05/results/P0-v04-pathology.md -- exact-belief substitution
\providecommand{\numBlockDeadAsk}{28}
\providecommand{\numBlockDeadRun}{280}

%% ---------------------------------------------------------------------------
%% research/v05/results/P0b-forcing-dilemma.md -- the forcing dilemma
%% ---------------------------------------------------------------------------
\providecommand{\numPatientWin}{56.2}
\providecommand{\numPatientWinCI}{54.93--57.47}
\providecommand{\numPatientDeclAcc}{98.40}
\providecommand{\numForcingDeclAcc}{85.66}
\providecommand{\numPatientCap}{22.5}
\providecommand{\numPatientDeadAsk}{51.2}
\providecommand{\numPatientDeadRun}{379}
\providecommand{\numPatientDeclWrong}{2.1}
\providecommand{\numPatientN}{1{,}500}

%% ---------------------------------------------------------------------------
%% research/v05/results/P1-deadlock-forensics.md and its two verifications
%% (P1-verify-deadlock-lock-census.md, P1-verify-e11-information.md)
%% ---------------------------------------------------------------------------
\providecommand{\numScanGames}{120}
\providecommand{\numLongGames}{14}
\providecommand{\numCycleLen}{245.1}
\providecommand{\numCycleTriples}{2.14}
\providecommand{\numPureCycle}{12}
\providecommand{\numNoLockPooled}{41}
\providecommand{\numNoLockPooledN}{62}
\providecommand{\numNoLockPct}{66}
\providecommand{\numRunAsksLocked}{13.0}
\providecommand{\numCycleInLock}{3.2}
\providecommand{\numCertN}{2{,}880}
\providecommand{\numCertRate}{68.5}
\providecommand{\numCertMean}{0.096}
\providecommand{\numCertMax}{0.664}
\providecommand{\numEntropyRate}{59.9}
\providecommand{\numEntropyMean}{0.325}
\providecommand{\numOutsideSupport}{20{,}898}
\providecommand{\numLadderRate}{82.9}
\providecommand{\numLadderMedian}{4}
\providecommand{\numTeamGain}{0.780}
\providecommand{\numOppGain}{0.021}
\providecommand{\numLeakRatio}{37}
\providecommand{\numDonationCost}{0.128}
\providecommand{\numDonationDead}{0.002}
\providecommand{\numInfoFreeStates}{42}
\providecommand{\numInfoRank}{17.6}
\providecommand{\numNoRepeatWin}{42.2}
\providecommand{\numTraceLive}{51}
\providecommand{\numTraceLegal}{54}
\providecommand{\numTraceRepeats}{142}
\providecommand{\numDonationNil}{99.6}
\providecommand{\numOwnProbPairs}{36}
\providecommand{\numOwnProbMean}{0.068}
\providecommand{\numOwnProbMax}{0.342}
\providecommand{\numProvableThresh}{0.999}
\providecommand{\numNoRepeatBaseEvents}{140.6}
\providecommand{\numNoRepeatEvents}{101.1}
\providecommand{\numNoRepeatBaseRun}{284}
\providecommand{\numNoRepeatRun}{4}
\providecommand{\numNoRepeatGames}{2{,}000}
%% P1-verify-e11-information.md sections 5-6 -- the original-seed figures that the
%% pooled re-measurement moved, quoted in sec:limitations
\providecommand{\numCertRateOrig}{70.6}
\providecommand{\numCertMeanOrig}{0.126}
\providecommand{\numLadderOrigN}{21}
\providecommand{\numMapRatioOrig}{14.3}
\providecommand{\numMapRatioPooled}{3.3}

%% ---------------------------------------------------------------------------
%% The compiled v0.4 ask weights (engine/src/v04.hpp:70--92) and the static
%% score evaluation in research/v05/lit/v05-refresh.md 2.1
%% ---------------------------------------------------------------------------
\providecommand{\numHitWeight}{11.506}
\providecommand{\numOwnershipPay}{+8.76}
\providecommand{\numOwnSetWeight}{1.688}
\providecommand{\numTeamControlWeight}{2.133}
\providecommand{\numLockWeight}{4.071}
\providecommand{\numCompletionWeight}{1.428}
\providecommand{\numRepeatWeight}{1.270}
\providecommand{\numEntropyWeight}{-2.6534}
\providecommand{\numDeadScore}{7.44}
\providecommand{\numLiveScore}{5.09}

%% research/v05/DESIGN.md 0.2 -- voluntariness of the dead asks, and the
%% dead-ask-filter prototype that isolates the cause
\providecommand{\numVoluntaryDead}{99.3}
\providecommand{\numFilterBaseDeadRun}{289}
\providecommand{\numFilterBaseGames}{32}
\providecommand{\numFilterBaseDeclWrong}{10.9}
\providecommand{\numFilterDeadRun}{1}
\providecommand{\numFilterDeclWrong}{1.9}
\providecommand{\numFilterWin}{49.75}
\providecommand{\numFilterBasePostHorizon}{474}

%% research/v05/DESIGN.md 0.3 -- the time-varying holding cost, measured and rejected
\providecommand{\numHoldEventsBefore}{146}
\providecommand{\numHoldEventsAfter}{121}
\providecommand{\numHoldRunGamesBefore}{32}
\providecommand{\numHoldRunGamesAfter}{18.5}
\providecommand{\numHoldPninetyBefore}{313}
\providecommand{\numHoldPninetyAfter}{311}
\providecommand{\numHoldLongestBefore}{289}
\providecommand{\numHoldLongestAfter}{289}
\providecommand{\numHoldDeclBefore}{10.9}
\providecommand{\numHoldDeclAfter}{13.4}
\providecommand{\numHoldLoss}{3.9}

%% ---------------------------------------------------------------------------
%% research/v05/results/P2-forced-endgame.md and P8-coordination.md
%% ---------------------------------------------------------------------------
%% NOTE on the sample size.  P2-forced-endgame.md quotes 562 forced declarations
%% over 12,000 mirror games at --rotations=6.  P2-verify-forced-endgame.md section 5
%% and P2-verify-forced-ladder.md section 4 establish that in a MIRROR match the six
%% rotations are three deal shifts x two team orientations, and both orientations
%% run identical policies, so each game is played twice: the 562 events are 281
%% distinct declarations over 6,000 distinct games (2,000 deals x three seat
%% shifts; \numForcedGames counts games, not deals).  The rate is unaffected
%% (281/281 = 100%); the sample size is halved here.
\providecommand{\numForcedGames}{6{,}000}
\providecommand{\numForcedN}{281}
\providecommand{\numForcedMissRate}{92.2}
\providecommand{\numFeasibleRight}{40.6}
\providecommand{\numRungZeroPairs}{210}
\providecommand{\numLadderBottomRung}{0.10}

%% ---------------------------------------------------------------------------
%% research/v05/results/P4-policy-review.md -- the forcing horizon
%% ---------------------------------------------------------------------------
\providecommand{\numStageTwoCost}{8.13}
\providecommand{\numPreHorizonWrong}{1.83}
\providecommand{\numHorizonWrong}{66.54}
\providecommand{\numHorizonReach}{23.0}
\providecommand{\numEfourteenCost}{2.67}
\providecommand{\numEfourteenRange}{1--5}
\providecommand{\numLowConfDecl}{130}
\providecommand{\numLowConfGames}{200}
\providecommand{\numLowConfThresh}{0.1}
\providecommand{\numHorizonGames}{46}

%% ---------------------------------------------------------------------------
%% research/v05/results/P3-deception.md -- the policy prior under deception
%% ---------------------------------------------------------------------------
\providecommand{\numAgnosticCost}{4.4}
\providecommand{\numAgnosticCI}{1.6, 7.2}
\providecommand{\numPriorDoubleCost}{5.3}
\providecommand{\numPhiWin}{49.00}
\providecommand{\numPhiCI}{45.8, 52.2}
\providecommand{\numPhiPaired}{+0.3}
\providecommand{\numFeintTrue}{28.2}
\providecommand{\numFeintBelief}{37.8}
\providecommand{\numFeintBias}{0.096}
\providecommand{\numBeliefErrHonest}{0.2255}
\providecommand{\numVsFeint}{54.0}
\providecommand{\numVsSilent}{56.1}
\providecommand{\numVsWithholder}{67.8}
\providecommand{\numFeintCalib}{0.07}
\providecommand{\numBeliefErrTol}{0.004}

%% ---------------------------------------------------------------------------
%% research/v05/results/P5-human-strategy.md and P5-verify-target-channel.md
%% ---------------------------------------------------------------------------
\providecommand{\numAskDecisions}{154{,}318}
\providecommand{\numTargetTwo}{46.6}
\providecommand{\numTargetBits}{0.639}
\providecommand{\numTargetSpread}{0.081}
\providecommand{\numTargetTwoVthree}{66.2}
\providecommand{\numTargetBitsVthree}{0.919}
\providecommand{\numProvableOwn}{0.34}
\providecommand{\numRouteTwo}{45.5}

%% ---------------------------------------------------------------------------
%% research/v05/results/P6-declaration.md and P8-coordination.md -- channels
%% ---------------------------------------------------------------------------
\providecommand{\numArbCost}{0.37}
\providecommand{\numArbGames}{30{,}000}
\providecommand{\numTransferRate}{0.148}
\providecommand{\numTransferOracle}{800}
\providecommand{\numTransferOracleN}{886}
\providecommand{\numTransferValue}{-0.0007}
\providecommand{\numTransferSE}{0.0024}
\providecommand{\numLadderLeak}{2.19}
\providecommand{\numTransferAskShare}{0.92}
\providecommand{\numTransferGames}{6{,}000}

%% ---------------------------------------------------------------------------
%% research/v05/results/P7-valuefn.md -- the value model and the fitting panel
%% ---------------------------------------------------------------------------
\providecommand{\numValueRSq}{0.2327}
\providecommand{\numValueRSqScore}{0.2354}
\providecommand{\numValueRSqTree}{0.388}
\providecommand{\numValueRSqTail}{0.168}
\providecommand{\numDeadFeatureDecisions}{31{,}788}
\providecommand{\numAskSideValue}{+0.006}
\providecommand{\numAskSideCI}{-0.016, +0.029}
\providecommand{\numDeclSideValue}{+0.025}
\providecommand{\numDeclSideCI}{+0.004, +0.047}
\providecommand{\numSoftMinBeta}{10}
\providecommand{\numSoftMinRatio}{1.9}
\providecommand{\numDeadFeatureSwing}{9}
\providecommand{\numPanelFloor}{0.664}
\providecommand{\numTailEvents}{160}

%% ---------------------------------------------------------------------------
%% v0.5 single-mechanism measurements (200 deals x 6 rotations, seed 90210),
%% and the v0.4 figures they are compared against.  These are regenerated by
%% the v0.5 mode of engine/build_tables.py once sec:results is written.
%% ---------------------------------------------------------------------------
\providecommand{\numMechDeals}{200}
\providecommand{\numMechSeed}{90210}
\providecommand{\numControlWin}{50.17}
\providecommand{\numMoneAlone}{49.50}
\providecommand{\numMonePAlone}{50.75}
\providecommand{\numMtwoAlone}{51.08}
\providecommand{\numMeightAlone}{55.17}
\providecommand{\numGuardAlone}{50.17}

%% ---------------------------------------------------------------------------
%% research/v05/runs/v04vector-in-v05.txt -- the v0.5 mechanisms on v0.4's frozen
%% 34-coordinate vector.  TWO configurations, which must not be pooled:
%%   numVfourParam*  the DESIGN-TIME assembly (ownership-by-p scaling and the
%%                   repetition guard both ON, both later rejected), 200 deals x
%%                   6 rotations at seed 90210 -- the bank sec:mech-built's
%%                   pre-refit figure came from;
%%   numShipParam*   the assembly that actually SHIPS, on the same bank
%%                   (numShipParamSameBank), on two larger banks, and on E2's
%%                   mirror pathology instrument.
%% numRejectedCost is the gap between the two on the shared bank.
%% Every one of these is \renewcommand'ed from that artifact.
%% ---------------------------------------------------------------------------
\providecommand{\numVfourParamWin}{41.8}
\providecommand{\numVfourParamCI}{39.25--44.50}
\providecommand{\numVfourParamDecl}{4.21}
\providecommand{\numVfourDecl}{4.75}
\providecommand{\numShipParamSameBank}{49.08}
\providecommand{\numShipParamSameCI}{46.67--51.58}
\providecommand{\numRejectedCost}{7.3}
\providecommand{\numShipParamWin}{50.22}
\providecommand{\numShipParamCI}{48.40--52.03}
\providecommand{\numShipParamRange}{50.13--50.30}
\providecommand{\numShipParamSeeds}{2}
\providecommand{\numShipParamGames}{6{,}000}
\providecommand{\numShipParamDecl}{4.49}
\providecommand{\numShipParamVfourDecl}{4.47}
\providecommand{\numShipParamDeadAsk}{0.02}
\providecommand{\numShipParamDeadRun}{1}
\providecommand{\numShipParamDeadRunGames}{0}
\providecommand{\numShipParamRepeat}{2.97}
\providecommand{\numShipParamDeclWrong}{1.89}
\providecommand{\numShipParamPostHorizon}{0}
\providecommand{\numShipParamForcedN}{4}

%% ---------------------------------------------------------------------------
%% Carried over from the v0.4 study (paper/numbers.tex) for the corrections
%% ---------------------------------------------------------------------------
\providecommand{\numVfourVsVthree}{75.07}
\providecommand{\numVfourLimitRate}{0.000\%}

%% ---------------------------------------------------------------------------
%% v0.5 EVALUATION BATTERY.  Defaults for the macros the results, limitations
%% and conclusion sections quote; every one of them is \renewcommand'ed by
%% paper/numbers_v05_generated.tex from the artifact named above it, so these
%% values only reach the PDF if the battery has not been run.
%% ---------------------------------------------------------------------------

%% research/v05/results/E1-verify.txt -- the rules audit on the frozen config
\providecommand{\numVerifyViolations}{0}
\providecommand{\numVerifyChecks}{23{,}594{,}580}

%% research/v05/results/E2-pathology.txt -- 600 mirror games per arm, seed 31
\providecommand{\numFourDeadAsk}{39.04}
\providecommand{\numFourDeadRunGames}{34.33}
\providecommand{\numFourDeadRunLongest}{286}
\providecommand{\numFourDeclWrong}{10.44}
\providecommand{\numFourPostHorizonDecl}{768}
\providecommand{\numFourRepeatAsk}{40.03}
\providecommand{\numFiveRepeatAsk}{2.63}
\providecommand{\numFiveDeadAsk}{0}
\providecommand{\numFiveDeadRunGames}{0}
\providecommand{\numFiveDeadRunLongest}{0}
\providecommand{\numFiveDeclWrong}{2.07}
\providecommand{\numFivePostHorizonDecl}{0}

%% research/v05/results/E3-headtohead.jsonl -- v0.5 vs v0.4 at five seed banks
\providecommand{\numHeadWin}{50.79}
\providecommand{\numHeadRange}{49.50--52.33}
\providecommand{\numHeadSeeds}{5}
\providecommand{\numHeadDeals}{1{,}500}
\providecommand{\numHeadGames}{9{,}000}
\providecommand{\numHeadGames}{9{,}000}

%% research/v05/results/E4-perstyle.jsonl -- the nine-style panel, 300 deals in
%% six rotations per cell, seed 515253 (engine/experiments_v05.sh).  Per-style
%% cells carry the "Panel" prefix so that they cannot collide with the v0.4
%% study's published figures quoted in sec:corrections.
\providecommand{\numStyleDeals}{300}
\providecommand{\numStyleSeed}{515253}
\providecommand{\numStyleCIVfour}{48.83--53.39}
\providecommand{\numPanelVfiveVfour}{51.11}
\providecommand{\numDeltaVfour}{+1.11}
\providecommand{\numDeltaVthree}{-1.00}
\providecommand{\numDeltaVtwo}{-1.78}
\providecommand{\numDeltaLockout}{+1.61}
\providecommand{\numDeltaDiversifier}{+0.89}
\providecommand{\numDeltaHunter}{+0.06}
\providecommand{\numDeltaBluffer}{-0.06}
\providecommand{\numDeltaRandom}{0.00}
\providecommand{\numVfiveMean}{83.60}
\providecommand{\numVfourMean}{83.57}
\providecommand{\numVfiveWorst}{51.11}
\providecommand{\numVfourWorst}{50.00}
\providecommand{\numVfiveRegret}{1.78}
\providecommand{\numVfourRegret}{1.61}

%% research/v05/results/E8-forced-endgame.txt -- 4,000 deals in six rotations
%% per arm, seed 909090, reported per DECLARING team
\providecommand{\numVfiveForcedAcc}{24.35}
\providecommand{\numVfourForcedAcc}{0.14}
\providecommand{\numForcedRateRatio}{6.4}
\providecommand{\numForcedEightGames}{24{,}000}

%% research/v05/results/E10-deception.md -- 400 deals in six rotations per cell
%% at two independent seed banks
\providecommand{\numDecSeeds}{2}
\providecommand{\numDecGames}{2{,}400}
\providecommand{\numDecSilentDelta}{+2.3}
\providecommand{\numDecFeintFive}{51.52}
\providecommand{\numDecFeintFour}{53.71}
\providecommand{\numDecFeintDelta}{-2.2}
\providecommand{\numDecWithholderFive}{72.53}
\providecommand{\numDecWithholderFour}{65.35}
\providecommand{\numDecWithholderDelta}{+7.2}
\providecommand{\numDecWorst}{50.96}
\providecommand{\numDecWorstOpp}{feint}
\providecommand{\numThetaSweep}{0.20, 0.26, 0.35, 0.445, 0.60}
\providecommand{\numThetaSweepN}{5}

%% research/v05/runs/v05-fitted.txt and engine/src/v04.hpp -- the policy-prior
%% weight before and after the v0.5 fit
\providecommand{\numPriorThetaFour}{0.264}
\providecommand{\numPriorThetaFive}{0.445}

%% research/v05/runs/fit-round1.jsonl, research/v04/runs/ -- the fit itself
\providecommand{\numFitRounds}{1}
\providecommand{\numFitGens}{40}
\providecommand{\numFitWorst}{57.00}
\providecommand{\numVfourFitRounds}{5}

%% research/v05/results/P5-human-strategy.md -- the discarded target dimension
\providecommand{\numTargetBitsGame}{41}

%% ---------------------------------------------------------------------------
%% sec:fitting.  The flat parameter vector, the freeze/parse round-trip
%% assertion, the CEM schedule, the objective, and the winner's-curse guard.
%% ---------------------------------------------------------------------------

%% engine/src/factory.hpp, engine/freeze_config_v05.py -- the vector's shape and
%% the offset the v0.4 pipeline hard-coded
\providecommand{\numFitDims}{34}
\providecommand{\numFitAskFeats}{20}
\providecommand{\numFitKnobs}{14}
\providecommand{\numAliasOffset}{18}

%% research/v05/runs/roundtrip-assert.txt -- the mechanical assertion that the
%% parser and the freezer agree, its negative control at the v0.4 offset, and
%% the bake-precision check
\providecommand{\numRoundTrip}{50.0000}
\providecommand{\numRoundTripPass}{passes}
\providecommand{\numRoundTripSeeds}{2}
\providecommand{\numRoundTripGames}{6{,}000}
\providecommand{\numAliasWin}{45.87}
\providecommand{\numAliasCI}{44.07--47.70}
\providecommand{\numAliasLoss}{4.13}
\providecommand{\numAliasGames}{3{,}000}
\providecommand{\numQuantWin}{50.43}
\providecommand{\numQuantCI}{49.43--51.40}
\providecommand{\numQuantGames}{3{,}000}

%% docs/FISHBOT_V05.md section 9 -- the fitting command of record
\providecommand{\numFitPop}{24}
\providecommand{\numFitElite}{6}
\providecommand{\numFitSeed}{505101}
\providecommand{\numFitDeals}{200}
\providecommand{\numFitCellGames}{400}
\providecommand{\numFitTotalGames}{2{,}304{,}000}

%% research/v05/runs/fit-round1.jsonl -- the panel, the objective's temperature,
%% and the gradient-weight arithmetic that makes a soft minimum a minimum
\providecommand{\numFitPanel}{v0.4, v0.3, lockout, detective, diversifier, hunter}
\providecommand{\numFitPanelN}{6}
\providecommand{\numFitBeta}{25}
\providecommand{\numFitSpan}{0.415}
\providecommand{\numFitRatio}{32048.3}
\providecommand{\numFitRatioAtTen}{63.43}
\providecommand{\numFitRatioCross}{4.83}

%% research/v04/runs/selected.json -- the v0.4 objective read the same way
\providecommand{\numSoftMinSpan}{0.063}
\providecommand{\numSoftMinPanelN}{4}

%% research/v05/runs/fit-round1.jsonl -- the best generation, and the final
%% common-seed re-evaluation the tuner returns from
\providecommand{\numFitScore}{0.5693}
\providecommand{\numFitBestGen}{24}
\providecommand{\numFitProfile}{57.0 / 75.2 / 80.5 / 79.0 / 94.8 / 98.5}
\providecommand{\numFitGuardWinner}{mean}
\providecommand{\numFitGuardMu}{0.5272}
\providecommand{\numFitGuardBest}{0.4798}
\providecommand{\numFitGuardGap}{0.0474}
\providecommand{\numFitCurse}{0.0895}

%% derived: the best generation's in-panel mirror cell against the held-out
%% multi-seed mean of the frozen configuration
\providecommand{\numFitRegression}{6.21}

%% ===========================================================================
%% sec:results (paper/sections_v05/10-results.tex).  Defaults for the whole
%% results macro set.  Every value below is the figure carried by the artifact
%% named above it and is regenerated by engine/build_tables_v05.py, which
%% \renewcommand's over these; they exist so the manuscript compiles when the
%% battery has not been run.
%% ===========================================================================

% research/v05/results/E1-verify.txt
\providecommand{\numVerifyDeterminism}{PASS}
\providecommand{\numVerifyLimitGames}{0}
\providecommand{\numVerifySetFail}{0}

% engine/experiments_v05.sh
\providecommand{\numAblFiveSeed}{606060}
\providecommand{\numCalibFiveDeals}{400}
\providecommand{\numCalibFiveSeed}{717171}
\providecommand{\numDialectSeed}{828282}
\providecommand{\numForcedEightDeals}{4{,}000}
\providecommand{\numForcedEightSeed}{909090}
\providecommand{\numHeadSeedFive}{424242}
\providecommand{\numHeadSeedFour}{777001}
\providecommand{\numHeadSeedOne}{90210}
\providecommand{\numHeadSeedThree}{515151}
\providecommand{\numHeadSeedTwo}{31337}
\providecommand{\numPathologyCrossSeed}{90210}
\providecommand{\numPathologyDeals}{300}
\providecommand{\numPathologySeed}{31}
\providecommand{\numVerifyGames}{600}

% research/v05/results/E2-pathology.txt
\providecommand{\numCrossAsks}{90.5}
\providecommand{\numCrossCapRate}{0}
\providecommand{\numCrossDeadAsk}{3.07}
\providecommand{\numCrossDeadRunGames}{0}
\providecommand{\numCrossDeadRunLongest}{1}
\providecommand{\numCrossDeadRuns}{1{,}665}
\providecommand{\numCrossDeclN}{5{,}400}
\providecommand{\numCrossDeclWrong}{2.26}
\providecommand{\numCrossEventsMean}{99.8}
\providecommand{\numCrossEventsMedian}{99}
\providecommand{\numCrossEventsPninety}{115}
\providecommand{\numCrossEventsPninetynine}{132}
\providecommand{\numCrossForced}{29}
\providecommand{\numCrossForcedWrong}{89.66}
\providecommand{\numCrossHit}{54.6}
\providecommand{\numCrossOwnLocked}{9.2}
\providecommand{\numCrossPostHorizonDecl}{0}
\providecommand{\numCrossRepeatAsk}{4.98}
\providecommand{\numCrossStarved}{0.0018}
\providecommand{\numFiveAsks}{87.2}
\providecommand{\numFiveCapRate}{0}
\providecommand{\numFiveDeadRuns}{0}
\providecommand{\numFiveDeclN}{5{,}400}
\providecommand{\numFiveEventsMean}{96.6}
\providecommand{\numFiveEventsMedian}{96}
\providecommand{\numFiveEventsPninety}{112}
\providecommand{\numFiveEventsPninetynine}{125}
\providecommand{\numFiveForced}{2}
\providecommand{\numFiveForcedWrong}{0}
\providecommand{\numFiveHit}{55.47}
\providecommand{\numFiveOwnLocked}{9.63}
\providecommand{\numFiveStarved}{0}
\providecommand{\numFourAsks}{134.3}
\providecommand{\numFourCapRate}{0}
\providecommand{\numFourDeadRuns}{2{,}610}
\providecommand{\numFourDeclN}{5{,}400}
\providecommand{\numFourEventsMean}{143.6}
\providecommand{\numFourEventsMedian}{106}
\providecommand{\numFourEventsPninety}{312}
\providecommand{\numFourEventsPninetynine}{321}
\providecommand{\numFourForced}{28}
\providecommand{\numFourForcedWrong}{100}
%% mirror columns are reported per distinct game: `fish pathology --rotations=2`
%% plays both team orientations, and in a mirror those are byte-identical games
\providecommand{\numFourGames}{300}
\providecommand{\numFiveGames}{300}
\providecommand{\numCrossGames}{600}
\providecommand{\numFourHit}{34.25}
\providecommand{\numFourOwnLocked}{16.46}
\providecommand{\numFourPostHorizonWrong}{58.59}
\providecommand{\numFourStarved}{0.2}

% research/v05/results/E3-headtohead.jsonl
\providecommand{\numHeadCIFive}{48.61--53.61}
\providecommand{\numHeadCIFour}{48.67--53.11}
\providecommand{\numHeadCIOne}{47.89--52.33}
\providecommand{\numHeadCIThree}{50.00--54.72}
\providecommand{\numHeadCITwo}{47.22--51.83}
\providecommand{\numHeadWinFive}{51.11}
\providecommand{\numHeadWinFour}{50.89}
\providecommand{\numHeadWinOne}{50.11}
\providecommand{\numHeadWinThree}{52.33}
\providecommand{\numHeadWinTwo}{49.50}

% research/v05/results/E3-headtohead.jsonl (derived)
\providecommand{\numHeadDeclAccFive}{98.15}
\providecommand{\numHeadDeclAccFour}{98.54}
\providecommand{\numHeadDeclRateFive}{4.51}
\providecommand{\numHeadDeclRateFour}{4.46}
\providecommand{\numHeadEvents}{99.8}
\providecommand{\numHeadSetsFive}{4.52}
\providecommand{\numHeadSetsFour}{4.48}

% research/v05/results/E4-perstyle.jsonl
\providecommand{\numPanelVfiveBluffer}{99.89}
\providecommand{\numPanelVfiveDetective}{76.78}
\providecommand{\numPanelVfiveDiversifier}{93.78}
\providecommand{\numPanelVfiveHunter}{97.72}
\providecommand{\numPanelVfiveLockout}{79.56}
\providecommand{\numPanelVfiveRandom}{100.00}
\providecommand{\numPanelVfiveVthree}{72.33}
\providecommand{\numPanelVfiveVtwo}{81.28}
\providecommand{\numPanelVfourBluffer}{99.94}
\providecommand{\numPanelVfourDetective}{77.28}
\providecommand{\numPanelVfourDiversifier}{92.89}
\providecommand{\numPanelVfourHunter}{97.67}
\providecommand{\numPanelVfourLockout}{77.94}
\providecommand{\numPanelVfourRandom}{100.00}
\providecommand{\numPanelVfourVfour}{50.00}
\providecommand{\numPanelVfourVthree}{73.33}
\providecommand{\numPanelVfourVtwo}{83.06}
\providecommand{\numStyleGames}{1{,}800}
\providecommand{\numVfiveStyleN}{9}
\providecommand{\numVfiveWorstOpp}{v0.4}
\providecommand{\numVfourWorstOpp}{v0.4}

% research/v05/results/E4-perstyle.jsonl (derived)
\providecommand{\numDeltaDetective}{-0.50}
\providecommand{\numVfiveRegretOpp}{v0.2}
\providecommand{\numVfourRegretOpp}{lockout}

% research/v05/results/E5-ablations.jsonl
\providecommand{\numAblControl}{50.87}
\providecommand{\numAblControlCI}{48.60--53.20}
\providecommand{\numAblControlDeclAcc}{89.02}
\providecommand{\numAblControlEvents}{141.2}
\providecommand{\numAblControlOppDeclAcc}{88.63}
\providecommand{\numAblFiveDeals}{250}
\providecommand{\numAblFiveGames}{1{,}500}
\providecommand{\numAblFull}{49.07}
\providecommand{\numAblFullCI}{46.60--51.53}
\providecommand{\numAblFullDeclAcc}{98.31}
\providecommand{\numAblFullEvents}{99.7}
\providecommand{\numAblFullHit}{53.61}
\providecommand{\numAblFullOppDeclAcc}{98.49}
\providecommand{\numAblMeight}{56.60}
\providecommand{\numAblMeightCI}{54.27--58.93}
\providecommand{\numAblMeightDeclAcc}{96.60}
\providecommand{\numAblMeightEvents}{141.7}
\providecommand{\numAblMeightOppDeclAcc}{83.55}
\providecommand{\numAblMone}{49.53}
\providecommand{\numAblMoneCI}{47.07--52.00}
\providecommand{\numAblMoneDeclAcc}{98.76}
\providecommand{\numAblMoneEvents}{99.7}
\providecommand{\numAblMoneMeight}{49.53}
\providecommand{\numAblMoneMeightCI}{47.07--52.00}
\providecommand{\numAblMoneMeightDeclAcc}{98.76}
\providecommand{\numAblMoneMeightEvents}{99.7}
\providecommand{\numAblMoneMeightOppDeclAcc}{98.49}
\providecommand{\numAblMoneMtwo}{49.07}
\providecommand{\numAblMoneMtwoCI}{46.60--51.53}
\providecommand{\numAblMoneMtwoDeclAcc}{98.31}
\providecommand{\numAblMoneMtwoEvents}{99.7}
\providecommand{\numAblMoneMtwoOppDeclAcc}{98.49}
\providecommand{\numAblMoneOppDeclAcc}{98.49}
\providecommand{\numAblMtwo}{52.60}
\providecommand{\numAblMtwoCI}{50.27--54.93}
\providecommand{\numAblMtwoDeclAcc}{91.20}
\providecommand{\numAblMtwoEvents}{141.4}
\providecommand{\numAblMtwoMeight}{56.40}
\providecommand{\numAblMtwoMeightCI}{54.07--58.80}
\providecommand{\numAblMtwoMeightDeclAcc}{96.16}
\providecommand{\numAblMtwoMeightEvents}{141.7}
\providecommand{\numAblMtwoMeightOppDeclAcc}{83.55}
\providecommand{\numAblMtwoOppDeclAcc}{86.78}
\providecommand{\numAblOwnershipP}{47.73}
\providecommand{\numAblOwnershipPCI}{45.13--50.33}
\providecommand{\numAblOwnershipPDeclAcc}{97.50}
\providecommand{\numAblOwnershipPEvents}{99.2}
\providecommand{\numAblOwnershipPHit}{54.51}
\providecommand{\numAblOwnershipPOppDeclAcc}{98.72}
\providecommand{\numAblRepeatGuard}{42.93}
\providecommand{\numAblRepeatGuardCI}{40.47--45.33}
\providecommand{\numAblRepeatGuardDeclAcc}{98.37}
\providecommand{\numAblRepeatGuardEvents}{99.6}
\providecommand{\numAblRepeatGuardHit}{52.67}
\providecommand{\numAblRepeatGuardOppDeclAcc}{98.56}

% research/v05/results/E5-ablations.jsonl (derived)
\providecommand{\numAblOwnershipPDelta}{-1.33}
\providecommand{\numAblRepeatGuardDelta}{-6.13}
\providecommand{\numThroughputFixed}{279}
\providecommand{\numThroughputPatho}{164}
\providecommand{\numThroughputRatioFive}{1.7}

% research/v05/results/E6-calibration-v05.txt
\providecommand{\numCalibAskBrier}{0.12187}
\providecommand{\numCalibAskECE}{0.02234}
\providecommand{\numCalibAskLogloss}{0.36529}
\providecommand{\numCalibAskLowN}{3{,}593}
\providecommand{\numCalibAskLowObs}{2.0}
\providecommand{\numCalibAskLowPred}{4.6}
\providecommand{\numCalibAskMidN}{2{,}775}
\providecommand{\numCalibAskMidObs}{51.5}
\providecommand{\numCalibAskMidPred}{45.0}
\providecommand{\numCalibAskN}{36{,}467}
\providecommand{\numCalibAskObs}{53.2}
\providecommand{\numCalibAskPred}{53.0}
\providecommand{\numCalibDeclBrier}{0.01500}
\providecommand{\numCalibDeclECE}{0.01542}
\providecommand{\numCalibDeclLogloss}{0.05373}
\providecommand{\numCalibDeclMidN}{184}
\providecommand{\numCalibDeclMidObs}{89.7}
\providecommand{\numCalibDeclMidPred}{75.7}
\providecommand{\numCalibDeclN}{3{,}609}
\providecommand{\numCalibDeclObs}{98.2}
\providecommand{\numCalibDeclPred}{96.7}
\providecommand{\numCalibDeclTopObs}{99.9}
\providecommand{\numCalibDeclTopPred}{99.8}

% research/v05/results/E6-calibration-v05.txt (derived)
\providecommand{\numCalibDeclTopShare}{86.6}

% research/v05/results/E7-rules.jsonl
\providecommand{\numDialectDeals}{250}
\providecommand{\numDialectDefault}{51.60}
\providecommand{\numDialectDefaultCI}{49.13--54.07}
\providecommand{\numDialectDefaultForcedFive}{42.86}
\providecommand{\numDialectDefaultForcedFour}{0.00}
\providecommand{\numDialectDefaultForcedRate}{0.0047}
\providecommand{\numDialectDefaultHit}{53.91}
\providecommand{\numDialectGames}{1{,}500}
\providecommand{\numDialectLegacy}{52.13}
\providecommand{\numDialectLegacyCI}{49.60--54.60}
\providecommand{\numDialectLegacyForcedFive}{94.41}
\providecommand{\numDialectLegacyForcedFour}{88.20}
\providecommand{\numDialectLegacyForcedRate}{0.2027}
\providecommand{\numDialectNoCardless}{51.87}
\providecommand{\numDialectNoCardlessCI}{49.40--54.33}
\providecommand{\numDialectNoCardlessForcedFive}{42.86}
\providecommand{\numDialectNoCardlessForcedFour}{0.00}
\providecommand{\numDialectNoCardlessForcedRate}{0.0047}
\providecommand{\numDialectNoOOT}{52.47}
\providecommand{\numDialectNoOOTCI}{50.00--54.87}
\providecommand{\numDialectNoOOTDeclFive}{98.47}
\providecommand{\numDialectNoOOTDeclFour}{95.43}
\providecommand{\numDialectNoOOTForcedFive}{98.36}
\providecommand{\numDialectNoOOTForcedFour}{88.14}
\providecommand{\numDialectNoOOTForcedRate}{0.2027}
\providecommand{\numDialectNoOOTHit}{44.81}

% research/v05/results/E7-rules.jsonl (derived)
\providecommand{\numDialectForcedRatio}{43}

% research/v05/results/E8-forced-endgame.txt
\providecommand{\numVfiveDeclAccEight}{98.13}
\providecommand{\numVfiveDeclRateEight}{4.50}
\providecommand{\numVfiveForcedRate}{0.0048}
\providecommand{\numVfourDeclAccEight}{98.51}
\providecommand{\numVfourDeclRateEight}{4.47}
\providecommand{\numVfourForcedRate}{0.0307}

% research/v05/results/E9-throughput.txt
\providecommand{\numThroughputFive}{286}
\providecommand{\numThroughputGamesFive}{600}

% research/v05/results/E10-deception.md
\providecommand{\numDecDeals}{400}
\providecommand{\numDecFeintFiveA}{50.96}
\providecommand{\numDecFeintFiveB}{52.08}
\providecommand{\numDecFeintFourA}{54.13}
\providecommand{\numDecFeintFourB}{53.29}
\providecommand{\numDecSeedA}{31415926}
\providecommand{\numDecSeedB}{8675309}
\providecommand{\numDecSilentFiveA}{80.42}
\providecommand{\numDecSilentFiveB}{83.17}
\providecommand{\numDecSilentFourA}{79.96}
\providecommand{\numDecSilentFourB}{79.00}
\providecommand{\numDecWithholderFiveA}{73.63}
\providecommand{\numDecWithholderFiveB}{71.42}
\providecommand{\numDecWithholderFourA}{66.25}
\providecommand{\numDecWithholderFourB}{64.46}
\providecommand{\numTwelveStyles}{12}

% research/v05/results/P3-verify-deception.md
\providecommand{\numPriorDeleteCI}{2.63, 6.58}
\providecommand{\numPriorDeleteCost}{4.60}
\providecommand{\numPriorCells}{14}
\providecommand{\numPriorCellsWorse}{12}
\providecommand{\numPriorCellsTie}{1}
\providecommand{\numPriorCellsReversed}{1}
\providecommand{\numPriorSignP}{0.003}

% engine/src/v05.hpp
\providecommand{\numGuardReject}{6.0}
\providecommand{\numMonePReject}{5.6}

%% research/v05/results/E3-headtohead.jsonl
\providecommand{\numHeadBankDeals}{300}
\providecommand{\numHeadBankGames}{1{,}800}

%% research/v05/results/E3,E4,E5,E7 (derived)
\providecommand{\numLimitRows}{37}
\providecommand{\numLimitNonzero}{0}

%% research/v05/results/E4-perstyle.jsonl (derived)
\providecommand{\numStyleMeanGap}{0.04}

%% engine/src/main.cpp -- the policy pool `fish verify` sweeps
\providecommand{\numVerifyPool}{9}
\providecommand{\numVerifyPairs}{81}

%% ---- audit pass: denominators, spans and per-bank deltas ----------------
%% research/v05/results/E8-forced-endgame.txt (derived)
\providecommand{\numVfiveForcedN}{115}
\providecommand{\numVfiveForcedRightN}{28}
\providecommand{\numVfourForcedN}{737}
\providecommand{\numVfourForcedRightN}{1}

%% research/v05/results/E3-headtohead.jsonl (derived)
\providecommand{\numHeadDeclErrFour}{1.46}

%% research/v05/results/E3,E4,E5,E7 (derived) -- every independent
%% v0.5-against-v0.4 bank in the report, at the frozen configuration
\providecommand{\numHeadBanks}{8}
\providecommand{\numHeadBankSpan}{49.07--52.33}
\providecommand{\numHeadBankLow}{49.07}

%% research/v05/results/E10-deception.md (derived) -- per-bank deltas
\providecommand{\numDecSilentDeltaA}{+0.46}
\providecommand{\numDecSilentDeltaB}{+4.17}
\providecommand{\numDecFeintDeltaA}{-3.17}
\providecommand{\numDecFeintDeltaB}{-1.21}
\providecommand{\numDecWithholderDeltaA}{+7.38}
\providecommand{\numDecWithholderDeltaB}{+6.96}

%% ---------------------------------------------------------------------------
%% engine/build_tables_v05.py -- provenance of the manuscript's own figures,
%% computed by the generator over paper/sections_v05/ so the paragraph in
%% section 12 that quotes them cannot itself go stale.  Audited by
%% paper/check_provenance.py.
%% ---------------------------------------------------------------------------
\providecommand{\numProvUsed}{0}
\providecommand{\numProvGenerated}{0}
\providecommand{\numProvTranscribed}{0}
%% --- end numbers_v05 ---
% Generated from the experiment artifacts by engine/build_tables_v05.py.  It
% \renewcommand's every macro the artifacts settle, so the \providecommand
% defaults above act only as a compile-before-the-battery-runs fallback and can
% never silently reach the PDF once the battery has been run.
\InputIfFileExists{numbers_v05_generated}{}{}

\begin{document}
\maketitle
\thispagestyle{fancy}

\begin{abstract}
%% --- begin sections_v05/abstract ---
Canadian Fish, also called Literature, is a six-player imperfect-information card
game in which two teams of three score only by declaring an exact allocation of a
six-card half-suit among their own members, and in which nothing in the rules
compels anyone to declare. \vfour{} was evaluated against a panel of weaker
policies and won decisively. Played against a copy of itself it does something
else: over \numMirrorGames{} mirror games, \numDeadAsk\% of its asks name a card
it can \emph{prove} the target does not hold, \numRepeatAsk\% are exact repeats
of an ask the same seat has already made, \numDeadRunGames\% of games contain a
run of at least six consecutive provably dead asks --- the longest run is
\numDeadRunLongest{} --- and \numDeclWrong\% of its declarations are wrong,
rising to \numPostHorizonWrong\% for the declarations its own termination rule
forces. Against \vthree{} on the same instrument the dead-ask rate is
\numVthreeDeadAsk\%. The failure is invisible to the evaluation that shipped it
because it is a strong-opponent phenomenon.
This report diagnoses that failure, corrects the account of it given in the v0.4
study, and specifies the policy that repairs it. The deadlock is not an
information freeze: it is a deterministic two-question cycle, present in
\numPureCycle{} of \numLongGames{} long games at the forensic seed, and across
pooled seed banks no half-suit is locked to any team at any point in the cycle
in \numNoLockPct\% of long games. Where a
half-suit \emph{is} locked, further information is still obtainable ---
\numCertRate\% of legal asks inside such a half-suit strictly raise a teammate's
exact probability of the correct allocation, by a mean of \numCertMean{} and up
to \numCertMax{} --- so the claim in the v0.4 termination study that ownership
monotonicity freezes allocation information is false, and is retracted here with
the measurement that refutes it. The mechanism is arithmetic in the ask score:
four ownership features that are not gated by hit probability pay up to
$\numOwnershipPay$ at hit probability zero, against $\numHitWeight \cdot p$ for
the hit-probability term, so a card the team already owns outscores a genuine
chance at one it does not, and \numVoluntaryDead\% of dead asks are chosen when a
live ask was available. \vfive{} restricts the candidate set to asks with
strictly positive hard-consistent probability, replaces an unconstrained
per-card allocation guess ---
responsible for \emph{all} \numForcedN{} forced-endgame declarations over
\numForcedGames{} distinct mirror games being wrong --- with a capacity-feasible one, and
deletes an event-count rule that cashes a best candidate without inspecting its
probability at a measured cost of \numStageTwoCost{} points against a mirror
opponent. Refitted on a panel that for the first time contains an opponent of its
own strength, and frozen, it takes the mirror dead-ask rate to
\numFiveDeadAsk\%, exact repeats to \numFiveRepeatAsk\%, games containing a dead
run of six or more to \numFiveDeadRunGames\%, declarations at or after the old
horizon to \numFivePostHorizonDecl{}, and declaration error to
\numFiveDeclWrong\%, with the forcing rule switched off and no game reaching the
action cap. None of that shows up as an aggregate win rate. Head to head it scores
\numPanelVfiveVfour\% [\numStyleCIVfour] on the style bank and \numHeadWin\% over
\numHeadSeeds{} further banks of \numHeadGames{} games; over a nine-style panel the two means
are level, \numVfiveMean\% against \numVfourMean\%; and on minimax regret over
that panel \vfour{} is marginally the better of the two, \numVfourRegret{}
against \numVfiveRegret{}. Against three deceptive archetypes held out of the
fit, the differences are larger and run both ways: \numDecWithholderDelta{}
points on the withholding manoeuvre that prompted this work,
\numDecSilentDelta{} on a silent opponent, and \numDecFeintDelta{} against a
feint that manufactures a false ask-legality certificate. Each replicates in
sign at \numDecSeeds{} independent seed banks, and the withholding result
replicates in size as well ($\numDecWithholderDeltaA$ and
$\numDecWithholderDeltaB$); the silent one does not
($\numDecSilentDeltaA$ and $\numDecSilentDeltaB$). Seven further mechanisms are specified and
not yet built. This report claims no equilibrium and no bound on exploitability,
and, unlike the v0.4 study, it ran no exploitability probe at all.
%% --- end sections_v05/abstract ---
\end{abstract}

%% --- begin sections_v05/01-introduction ---
\section{Introduction}
\label{sec:intro}

Canadian Fish, also called Literature, is a six-player card game of imperfect
information: two teams of three, the whole deck dealt out, no communication
outside the public record of play. On turn a player names a card and an opponent;
if the opponent holds it the card transfers and the asker continues, otherwise
the turn passes to the opponent asked. Cards taken this way score nothing. A team
scores only by \emph{declaring} --- naming one half-suit and stating which of its
own members holds each of its six cards --- a correct declaration awarding the
half-suit to the declaring team and an incorrect one to the opponents.
\S\ref{sec:rules} states every consequential rule choice in the dialect studied.
Two of them govern everything in this report: a declaration may be made at any
moment, and nothing in the rules ever obliges a player to make one.

\vfour{} \cite{fishbot04} was built on that structure. It computes an
observer-conditioned posterior over the initial deal, scores asks with a
twenty-feature linear function refined by a one-ply expectimax over a
sixteen-feature value model, and decides declarations by comparing fitted value
estimates. On its published evaluation bank it defeated \vthree{} \cite{fishbot03} in
\numVfourVsVthree\% of games. This report begins from the observation that the
published evaluation could not have detected what \vfour{} does against a strong
opponent, and that the project owner met the resulting behaviour directly at the
table: bots repeating one question at one opponent indefinitely, and then
misdeclaring in bulk when the game was hurried to a close.

\paragraph{The measurement.} \S\ref{sec:diagnosis} reports it. In mirror
self-play, \numDeadAsk\% of \vfour{}'s asks are \emph{provably dead}: the actor
could prove from its own public deduction state that the target does not hold
the named card. \numRepeatAsk\% of asks are exact repeats of an
(actor, card, target) triple the same game has already seen.
\numDeadRunGames\% of games contain a run of at least six consecutive provably
dead asks; the longest such run is \numDeadRunLongest{} asks. Against \vthree{}
on the identical instrument the dead-ask rate is \numVthreeDeadAsk\% and the
longest run is \numVthreeDeadRun{}. The pathology is a strong-opponent
phenomenon, which is precisely why a panel of weaker policies did not expose it,
and why a competent human --- who is closer to the mirror case than \vthree{} is
--- meets it and \vthree{} does not.

\paragraph{Five claims this report makes, and one it retracts.}

\begin{enumerate}
  \item \textbf{The deadlock is an ask-policy fixed point, not an information
        freeze.} The v0.4 termination study explains non-termination by appeal to
        the locked-half-suit theorem, asserting that no further information about
        a locked half-suit's allocation can ever arrive. Both halves of that
        account fail. The observed deadlock is a deterministic two-question
        cycle: in \numPureCycle{} of \numLongGames{} long games the entire dead
        run consists of exactly two distinct (actor, card, target) triples, and
        in \numNoLockPct\% of long games no half-suit is locked to any team at
        any point during the run. Independently of that, the informational claim
        is false on its own terms: \numCertRate\% of legal asks inside a
        half-suit a team does in fact own strictly increase a teammate's exact
        probability of the correct allocation, mean $+\numCertMean$, maximum
        $+\numCertMax$, and \numEntropyRate\% strictly lower its exact marginal
        entropy. \S\ref{sec:corrections} states the retraction.
  \item \textbf{The cause is arithmetic, and it is in the ask score.} Four
        ownership features are not gated by hit probability. Evaluated at the
        compiled weights they pay up to $\numOwnershipPay$ at $p=0$, against
        $\numHitWeight\cdot p$ for the hit-probability term, and the only
        explicit information term in the score, the binary entropy of the ask
        outcome, carries weight $\numEntropyWeight$ --- a negative
        value-of-information term, so a certain miss is rewarded for being
        certain. \numVoluntaryDead\% of the dead asks are voluntary: a live ask
        existed and was passed over.
  \item \textbf{Terminating by force is a strong-opponent tax, and terminating by
        patience does not terminate.} A configuration whose event-count
        escalation never fires beats the shipped one by \numPatientWin\%
        (deal-clustered \numPatientWinCI\%) and declares at \numPatientDeclAcc\%
        accuracy against \numForcingDeclAcc\%, but reaches the action cap in
        \numPatientCap\% of mirror games. The escalation's second stage, which
        cashes a best candidate without inspecting its probability at all, costs
        \numStageTwoCost{} points against a mirror opponent and exactly nothing
        against every weaker style tested. Both settings of the dial are bad; the
        missing capability is an ask rule that terminates the cycle by choosing
        informative asks rather than an event counter that ends the game by
        guessing.
  \item \textbf{The repairs that follow are keyed to measurements, and the ones
        that failed are recorded as failures.} \S\ref{sec:mechanisms} specifies
        ten mechanisms; three are built and seven are designed. It also records
        four plausible fixes that were measured and rejected --- a time-varying
        holding cost, deleting the policy prior, a turn-transfer willingness
        ladder, and confidence-ranked declaration arbitration --- so that they
        are not rebuilt.
  \item \textbf{The repair removes the failure mode, and does not show up as a
        win rate.} The three built mechanisms were refitted, on a panel that for
        the first time in this project's history contains an opponent of the
        agent's own strength (\S\ref{sec:fitting}), and frozen. In the mirror,
        the dead-ask rate goes from \numFourDeadAsk\% to \numFiveDeadAsk\%, exact
        repeats from \numFourRepeatAsk\% to \numFiveRepeatAsk\%, games containing
        a dead run of six or more from \numFourDeadRunGames\% to
        \numFiveDeadRunGames\%, and declaration error from \numFourDeclWrong\% to
        \numFiveDeclWrong\% --- with the event-count escalation switched off.
        Head to head \vfive{} scores
        \numPanelVfiveVfour\% [\numStyleCIVfour] on the style bank and
        \numHeadWin\% over \numHeadSeeds{} further banks; across the nine-style panel of
        \S\ref{sec:results} the means are \numVfiveMean\% against
        \numVfourMean\%; and on minimax regret over that panel \vfour{} is
        marginally the better of the two, \numVfourRegret{} against
        \numVfiveRegret{}. Where the two policies do separate is against
        deception: \numDecWithholderDelta{} points on the withholding manoeuvre
        that prompted this work and \numDecSilentDelta{} on a silent opponent,
        against \numDecFeintDelta{} on a feint that manufactures a false
        ask-legality certificate. Each replicates in sign at \numDecSeeds{}
        independent seed banks; the withholding and feint results replicate in
        size too, the silent one does not ($\numDecSilentDeltaA$ against
        $\numDecSilentDeltaB$).
\end{enumerate}

\paragraph{What this report does not claim.} No result here is an equilibrium
result, and nothing here bounds exploitability from above. The v0.4 study ran an
exploitability probe; this one ran none at all, so \vfive{} makes no robustness
claim that a best response could contradict, and the comparison of the two
policies is not a comparison of their exploitability. Following the project's
standing evaluation preference, no aggregate figure above appears without the
per-opponent breakdown and an explicit worst case beside it: over the twelve
styles tested, and on the per-style banks those figures come from, \vfive{}'s
worst case is \numDecWorst\% against the \numDecWorstOpp{} and \vfour{}'s is
\numVfourWorst\% against a copy of itself, which is the worst case for any
policy in this family. The fit behind those
figures is a single round from a single base seed, against \vfour{}'s
\numVfourFitRounds{} rounds, and the sixteen-coefficient value model inside the
policy was not refitted at all. Seven of the ten mechanisms of
\S\ref{sec:mechanisms} are specified and unbuilt, so nothing here measures the
design this report proposes --- only the three parts of it that exist.
%% --- end sections_v05/01-introduction ---
%% --- begin sections_v05/02-rules ---
\section{The game and the rules dialect studied}
\label{sec:rules}

Canadian Fish (also called Literature) is played by six players in two teams of
three, seated so that the teams alternate around the table. The deck is 54 cards,
partitioned into nine \emph{half-suits} of six cards each: the low band
$\{2,3,4,5,6,7\}$ and the high band
$\{9,10,\mathrm{J},\mathrm{Q},\mathrm{K},\mathrm{A}\}$ of each of the four suits,
together with a ninth half-suit containing the four eights and the two jokers.
The dealer is chosen at random, each player is dealt nine cards, and the player
to the dealer's left leads. A half-suit leaves play when it is claimed; the team
that takes at least five of the nine wins the game.

The dialect below is the one supplied by the project owner as the rules of
record, and the engine's \texttt{Rules} defaults match it on every point checked.
Published descriptions document mutually incompatible variants of several of
these choices \cite{pagat}; what follows is one explicit selection, not a
uniquely canonical rule set. It matches the strategy corpus this work draws on
\cite{develin}, and the contested choices are exposed as engine switches so that
their effect can be measured rather than assumed.

\paragraph{Asking.} On their turn a player asks one opponent for one named card.
The ask is legal only if the named card lies in a half-suit that is still live,
the asker holds at least one \emph{other} card of that half-suit, the asker does
not hold the named card, the target sits on the opposing team, and the target
still holds at least one card. If the target holds the named card they must
surrender it and the asker retains the turn; otherwise the turn passes to the
target. Every ask, its answer, every resulting transfer, and every player's card
count are public. The rules additionally license two public queries that the
engine models implicitly because both quantities are already in the public state:
the content of the previous two asks, and any player's remaining card count.

\paragraph{Declaring.} A claim on a half-suit is a \emph{declaration}: an exact
allocation naming, for each of the six cards, which member of the declaring
player's own team holds it. If every one of the six assignments is correct the
declaring team is awarded the half-suit; \emph{any} inaccuracy at all awards it
to the opposing team. Either way the six cards leave play, the half-suit becomes
inactive, and the turn returns to the player who held it before the declaration.
A declaration is legal at any moment, including during an opponent's turn. The
declarer need hold no card of the half-suit --- indeed need hold no cards at all
--- and may not declare on behalf of the opposing team.

\paragraph{Running out of cards.} A player with no cards leaves the asking cycle:
they can neither ask nor be asked, though they may still be named as the holder
of a card in a teammate's declaration and may still declare. If the turn reaches
a cardless player --- possible only after a declaration --- the turn passes to a
teammate who still holds cards. When an entire team is cardless, the other team
must declare every remaining half-suit with no further asking, and may misdeclare
and thereby hand half-suits to the cardless team.

\subsection{Two places where the engine is more restrictive than the rules}
\label{sec:rules-gap}

Both are coordination channels the rules permit and the engine does not use.
They are stated here because \S\ref{sec:mechanisms} reports that both were
measured and neither is worth building, which is a result about the game and not
only about the implementation.

\begin{enumerate}
  \item \textbf{Turn transfer.} The rules say that when a cardless player must
        pass the turn and both teammates hold cards, the team may openly
        strategise about which teammate receives it, but may share nothing beyond
        their willingness to receive it. The engine instead has the cardless
        player choose unilaterally from its own belief, and solicits no
        willingness bit. A genuine multi-candidate decision arises
        \numTransferRate{} times per game; \S\ref{sec:mech-rejected} reports what
        the channel is worth.
  \item \textbf{Declaration arbitration.} The rules do not say who acts first
        when two seats offer a declaration in the same instant. The engine
        resolves it by lowest seat, deliberately, to avoid comparing private
        confidences across seats. \S\ref{sec:mech-rejected} reports the measured
        cost of that choice.
\end{enumerate}

\subsection{Engine settings, and the switches that change them}
\label{sec:rules-dialect}

\begin{itemize}
  \item \textbf{Deck:} 54 cards, nine half-suits, nine per player
        (\texttt{-{}-sets=8} gives the 48-card, eight-half-suit form).
  \item \textbf{Declaration timing:} any seat may declare at any moment, in or
        out of turn, holding cards or not (\texttt{-{}-no-out-of-turn},
        \texttt{-{}-no-cardless-declare}).
  \item \textbf{Misdeclaration:} the half-suit is awarded to the opposing team;
        published variants differ here \cite{pagat}.
  \item \textbf{Simultaneous declarations:} the lowest seat index that offers a
        declaration acts first --- a modelling choice, not a rule
        (\texttt{-{}-arb=low\textbar high\textbar turn}).
  \item \textbf{Forced endgame:} an eight-level willingness ladder selects the
        declarer when one team is cardless (\texttt{-{}-legacy} gives the
        two-level v0.3 ladder). \S\ref{sec:diag-forced} reports that in the
        shipped policy this ladder is inert.
  \item \textbf{Action cap:} 400 asks, any residue adjudicated neutrally by
        majority physical holding (\texttt{-{}-maxasks}).
\end{itemize}

No prearranged signalling conventions are available to any policy in the
configuration measured here, and team communication is limited to the willingness
bit above and the successor choice made by a cardless turn-holder. The literature
records that pre-agreed conventions are explicitly forbidden in this game
\cite{develin} while also documenting a named convention in real use
\cite{pagat}, and a pair of policies trained together has an implicit
pre-agreement by construction; \S\ref{sec:mech-conventions} states how \vfive{}
handles that tension.

\paragraph{The action cap is a backstop, not a rule.} Two sufficiently patient
policies can stall indefinitely, and this report is largely about a configuration
that does. Its incidence is therefore reported with every experiment rather than
assumed to be zero: it is \numPatientCap\% for the unforced mirror configuration
of \S\ref{sec:diag-dilemma} and zero for the shipped \vfour{} configuration,
which avoids it by forcing declarations that \S\ref{sec:diag-dilemma} shows are
\numPostHorizonWrong\% wrong.
%% --- end sections_v05/02-rules ---
%% --- begin sections_v05/03-related ---
\section{Related work}
\label{sec:related}

The v0.4 study surveyed inference, search, team equilibria and signalling, and
that survey stands. This section covers what the present report additionally
depends on, and it is organised around a gap: a coverage sweep of the eight
reviews assembled for v0.4 finds no occurrence of \emph{war of attrition},
\emph{minimax regret}, \emph{distributionally robust}, \emph{Dynkin},
\emph{restricted Nash}, \emph{data-biased} or \emph{safe opponent exploitation},
and one occurrence of \emph{deadlock}. Those are exactly the two bodies of work
this report needs: the theory of voluntary stopping when both sides may wait, and
the theory of performing well against the worst opponent style rather than the
average one.

\subsection{Fish and Literature}

The published material on this game remains human and informal. McLeod's Pagat
entry \cite{pagat} is the canonical rules source and records the points on which
dialects differ; Develin's chapter \cite{develin} is the most detailed strategy
writing located, and supplies two observations this design uses --- that an ask
teaches the opponents as much as it teaches a teammate, and that pre-agreed
conventions are outside the game. A repeat search on 22 August 2026 --- the arXiv
API for \texttt{"Canadian Fish"}, \texttt{"Literature card game"} and
\texttt{"half-suit"}, returning zero results, plus a web sweep that surfaced only
adjacent games --- reproduced the v0.4 finding. Every novelty statement in this
report therefore takes the form ``we found no prior published work that \dots''
and should be read against that scope. The one agent located, Somani's
\texttt{literature} \cite{somani}, plays the four-player 48-card variant and
publishes no strength numbers; a closed-source mobile application advertises
six-player bots with no stated methodology \cite{playstore}; Literature is not
among the environments packaged with RLCard \cite{rlcard}.

\subsection{Voluntary stopping when waiting is free}
\label{sec:related-stopping}

The declaration decision in Fish is a multi-player, nonzero-sum, undiscounted
optimal-stopping game: each side chooses when to claim, the payoff depends on who
claims first and whether the claim was right, and waiting costs nothing in the
rules. Games of this shape are known to be badly behaved. Christensen,
Lindensj\"o and Neumann \cite{christensen-dynkin} record that for undiscounted
Dynkin games only $\epsilon$-equilibria are established in general, with
counterexamples showing that Nash equilibria need not exist even for
deterministic stationary rewards, and that this survives enlargement to
randomised stopping times. Hamad\`ene, Hassani and Morlais
\cite{hamadene-dynkin} prove existence of an $\epsilon$-Nash equilibrium for the
$N$-player discrete-time case by a constructive algorithm, under a payoff
structure that is the Fish structure: payoffs depend on the set of players who
stop at the first stage at which anyone stops. These are results for general
Markovian reward processes rather than theorems about Fish, so the correspondence
is one of structure and is used here to name the missing ingredient, not to
assert a property of the game.

The classical model of mutual refusal to move is the war of attrition. Hendricks,
Weiss and Wilson \cite{hendricks-attrition} characterise the complete-information
continuous-time case and find a continuum of equilibria including arbitrarily
long delay: delay is an equilibrium phenomenon that a symmetric deterministic
pair will select, and \vfour{}'s mirror match is exactly symmetric and exactly
deterministic. The constructive converse is Mguni \cite{mguni-actioncost}: impose
a strictly positive cost on each action and the stochastic game acquires a unique
value and a Markovian saddle point at strategically chosen stopping times. That
is the cleanest published statement of the natural fix, and \S\ref{sec:diagnosis}
reports that in this game it does not work --- the cost is a mean-shifter that
leaves the tail of the distribution untouched.

There is a reason to be careful with holding costs beyond their measured
ineffectiveness. Admati and Perry \cite{admati-perry} show that delay is used as
a signalling device precisely because it is a more efficient signal of strength
than the alternatives, and that the delay does not vanish as the interval between
moves goes to zero. Applied to Fish this is Farrell--Gibbons two-audience cheap
talk \cite{farrell-gibbons} on the \emph{timing} channel: if how long a seat
waits before declaring depends on its private confidence, three opponents can
read that confidence. Any holding cost must therefore be a deterministic function
of public state, or be deliberately randomised.

We found no substantive literature on self-inflicted non-termination in
game-playing agents. Three distinct query formulations returned only anecdote,
resolved in every case by an arbitrary move cap --- which is what \vfour{} does,
and what \S\ref{sec:diagnosis} measures the cost of. That absence is itself a
result: voluntary-claiming games are an under-studied class and Fish is a clean
instance of one.

\subsection{Information-seeking actions, and what an ask leaks}
\label{sec:related-info}

An ask in Fish is simultaneously a request for material, a certificate to one's
teammates and a disclosure to one's opponents. The v0.4 review formalised the
trade-off as a per-decision secrecy rate in the Wyner \cite{wyner} sense --- the
mutual information the action carries to teammates minus the mutual information
it carries to opponents --- and built none of the four instruments it listed.
\vfour{} contains three information terms in its ask score and all three are
negative; \S\ref{sec:diag-features} gives the arithmetic.

Two results bear directly on the repair. Alvim et al.\ \cite{alvim-leakage} model
information flow as a game in which the defender's utility is a convex function
of the \emph{distribution} over its actions rather than the expected utility of
the action played, and show that a Nash equilibrium nevertheless exists. The
consequence for Fish is sharp: a deterministic argmax ask policy has maximal
leakage by construction, so no reweighting of leak-penalty features can reduce
leakage below the deterministic bound. Reducing it requires randomising over
near-tied asks. The second consequence is a caveat --- because utility is convex
in the mixed strategy, the leakage game is not bilinear and standard regret
matching does not apply unmodified.

Lervold, Peterson and King \cite{lervold-infogain} appear to be the only
published work that places an explicit information-gain incentive inside an
information-set tree search, reporting improvement over baseline and over a
cheating search in a hidden-information multi-action wargame. We were unable to
obtain the full text, so the finding is used here and the formula is not.
Serrino et al.\ \cite{serrino-deeprole} integrate deductive reasoning directly
into vector-form counterfactual regret minimisation for five-player Avalon; it is
the closest published system to Fish's structure, a hidden-role adversarial team
game in which every action is public and its informational content is the whole
game, and it is the existence proof that hard deduction can live inside a solver
rather than beside it.

\subsection{Online opponent modelling, and robustness to being modelled}
\label{sec:related-opponent}

\vfour{} models the other five seats with two global scalars, identical for every
opponent and never updated during a game (\S\ref{sec:diag-prior}). The literature
on doing better within a single short episode is mature and was absent from the
v0.4 corpus. Ganzfried, Wang and Chiswick \cite{ganzfried-multiplayer} give the
only opponent-modelling result located that is explicitly in the more-than-two
player regime. Rosman, Hawasly and Ramamoorthy \cite{rosman-bpr} pose exactly
this problem --- select online, from a library of policies, a response to a novel
task instance, under a short episode --- and a Fish game is
\numMirrorEvents{} public events long. Smith et al.\ \cite{smith-qmixing} learn a
value function against each pure opponent and combine them for any mixture
without retraining, with a classifier that refines the mixture weights during the
episode.

The robustness half matters more here, because the behaviour that prompted this
work is an opponent deliberately feeding a model. Johanson, Zinkevich and Bowling
\cite{johanson-rnash} give the shape of the dial: solve a game in which the
opponent plays a fixed model with probability $p$ and is free otherwise, so that
$p=1$ is a best response and $p=0$ an equilibrium, and intermediate $p$ traces an
exploitation/robustness frontier. Johanson and Bowling \cite{johanson-dbr} set
$p$ per information set as a function of how much data is actually available
there --- the right refinement for Fish, where the per-half-suit evidence about a
given opponent is a handful of observations, and the one that makes a
deliberately silent opponent fall back to the population prior instead of being
misread. Ganzfried and Sandholm \cite{ganzfried-safe} guarantee at least the game
value per period while still exploiting mistakes; M\"uller et al.\
\cite{muller-bestofboth} give the modern bandit-feedback formalisation of the
same idea. Yang et al.\ \cite{yang-bayestomop} handle opponents that are
themselves running Bayesian opponent models, and Milec, Kova\v{r}\'ik and Lis\'y
\cite{milec-abd} adapt to sub-rational opponents beyond a depth limit while
remaining robust against rational play.

\subsection{Worst case rather than mean case}
\label{sec:related-worstcase}

The standing evaluation preference of this project --- report the per-opponent
breakdown and an explicit worst case, never an aggregate alone --- has a matching
equilibrium concept. Aggarwal and How \cite{aggarwal-prmre} define a
probabilistically robust minimax-regret equilibrium for adversarial team games
under asymmetric information, minimising worst-case regret over a
high-confidence subset of the type space rather than over the whole of it, and
solve it inside a robust double-oracle loop. Translated to Fish, the type is the
opponent style, the high-confidence subset is the style set actually built, and
the prescription is to optimise worst-case regret rather than either the mean or
the absolute worst case --- promoting a reporting rule to a training objective.
Two practical notes come with it. Xu et al.\ \cite{xu-mirror} observe that
computing max-regret requires the per-style optimum first, which in this setting
is one arena run per style, so minimax regret is available as a post-hoc
selection criterion over checkpoints and not only as a training objective.
Beukman et al.\ \cite{beukman-remidi} document the failure mode of a naive
worst-case objective: once the bound is attained everywhere, learning stagnates
because the objective stops distinguishing policies --- which is a direct warning
against a raw minimum over a small style set. Vinitsky et al.\
\cite{vinitsky-rap} report that a single adversary produces exploitable policies
and replace it with a sampled population.

\subsection{Team games, and the concepts this report does not compute}

Three teammates share a payoff, hold different information and cannot
communicate, so the target concept is a team-maxmin equilibrium with correlation
\cite{celli-gatti}. Nothing here computes one. Two recent results scope the
ambition. Anagnostides et al.\ \cite{anagnostides-teamzero} show that computing
Nash equilibria in team zero-sum games is PPAD-complete, and remains so at
inverse-polynomial precision and with two-player teams; that concerns the
uncorrelated equilibrium, which is what three independently acting seats without
a shared per-episode seed actually implement. Liu et al.\ \cite{liu-hpsro}
identify a concrete defect in Team-PSRO \cite{mcaleer-teampsro}: its policy
sharing cannot cover the joint policy space in \emph{heterogeneous} team games
where teammates play distinct roles, yielding higher exploitability. Fish seats
are heterogeneous in the relevant sense, since turn flow and distance to the next
opponent differ by seat, so any future population loop for this game should be
the heterogeneous variant.

Two design questions in \S\ref{sec:mechanisms} are the Hanabi-style zero-shot
coordination questions \cite{bard-hanabi} in another game: how much a policy may
rely on its partner running the same blueprint, and how to keep a convention from
degenerating into a private code. Off-belief learning \cite{hu-obl} and
other-play \cite{hu-otherplay} are the reference treatments, and the level of
grounding is the natural explicit knob for the partner-regime distinction of
\S\ref{sec:mech-partner}.

\subsection{Determinization, and why it is not the foundation used here}

Perfect Information Monte Carlo --- sample hidden cards consistently with the
public history, solve the resulting perfect-information game, average
\cite{ginsberg-gib} --- is unsound however many worlds are drawn
\cite{frank-basin-ai}, and is degenerate in Fish for a structural reason as well:
a clairvoyant player never fails an ask, so almost every reachable position
carries the same perfect-information value and the average collapses to choosing
the ask most likely to hit. The v0.4 study gives the collapse in full. It is
worth restating here because the collapse target --- ``choose the ask most likely
to hit'' --- is a policy that would not have deadlocked. \vfour{}'s pathology is
not that it is a hit-maximiser; it is that its ownership features outbid the
hit-probability term at $p=0$ (\S\ref{sec:diag-features}). The policy-aware
reweighting of a rules posterior that the Skat literature uses
\cite{solinas-supervised,rebstock-inference,buro-kermit} is the correct frame for
the prior that \S\ref{sec:diag-prior} examines, and the counting machinery
underneath \cite{lsw,cryan-dyer} is unchanged from v0.4.
%% --- end sections_v05/03-related ---
%% --- begin sections_v05/04-diagnosis ---
\section{Diagnosis: what \vfour{} does against a strong opponent}
\label{sec:diagnosis}

This is the central section of the report. It establishes four things in order:
what the failure is, that it is not the failure the v0.4 study described, what
causes it, and what each contributing defect costs. Every number is produced by
an instrument described in \S\ref{sec:diag-instrument} and carries the artifact
it came from.

\subsection{The instrument}
\label{sec:diag-instrument}

The measurements below come from a diagnostic command added for this study,
\texttt{fish pathology}, together with five probes that re-derive its headline
quantities independently. Three definitions do the work.

\begin{definition}[Provably dead ask]
\label{def:dead}
An ask of card $c$ from target $t$ by actor $a$ is \emph{provably dead} if $a$'s
own public deduction state establishes that $t$ does not hold $c$: either $c$'s
holder is already known and is not $t$, or $t$ has been excluded from $c$'s
candidate set by a certificate. The predicate is a deduction from the actor's own
hand and the public record. No posterior and no policy prior enter it, so a
policy that filters on it cannot acquire a modelling error --- it can only lose
moves whose probability of success is exactly zero.
\end{definition}

\begin{definition}[Dead run]
A maximal sequence of consecutive asks in one game, every one of which is
provably dead in the sense of Definition~\ref{def:dead}.
\end{definition}

\begin{definition}[Starved turn]
A turn at which \emph{every} legal ask is provably dead, so the actor has no live
alternative. Starved turns separate a policy that chooses dead asks from a policy
that is forced into them.
\end{definition}

Two populations are used throughout: \emph{mirror} play, in which both teams run
the same policy, and the \emph{held-out} contrast against \vthree{}, which is the
population on which \vfour{} was evaluated for publication.

\subsection{The baseline, and why the published evaluation could not see it}
\label{sec:diag-baseline}

\begin{table}[t]
\centering
\caption{\vfour{} measured on the same instrument against itself and against
\vthree{}. Mirror: \numMirrorGames{} games, seed 31. Held-out: 600 games, seed
90210. Source: \texttt{research/v05/results/P0-v04-pathology.md}.}
\label{tab:baseline}
\begin{tabular}{lrr}
\toprule
& \textbf{mirror} & \textbf{vs.\ \vthree{}} \\
\midrule
events per game                       & \numMirrorEvents{}   & \numVthreeEvents{} \\
ask hit rate                          & \numMirrorHit\%      & \numVthreeHit\% \\
provably dead asks                    & \numDeadAsk\%        & \numVthreeDeadAsk\% \\
exact repeat (actor, card, target)    & \numRepeatAsk\%      & \numVthreeRepeat\% \\
longest dead run                      & \numDeadRunLongest{} & \numVthreeDeadRun{} \\
games with a dead run $\geq 6$        & \numDeadRunGames\%   & \numVthreeDeadRunGames\% \\
starved turns                         & \numStarved\%        & \numVthreeStarved\% \\
asks inside a half-suit the actor's team already owns & \numOwnLocked\% & \numVthreeOwnLocked\% \\
declarations wrong                    & \numDeclWrong\%      & \numVthreeDeclWrong\% \\
declarations at or after the forcing horizon & \numPostHorizonDecl{} & 0 \\
\quad of those, wrong                 & \numPostHorizonWrong\% & --- \\
forced-endgame declarations wrong     & \numForcedWrongRate\% & \numForcedWrongRate\% \\
\bottomrule
\end{tabular}
\end{table}

Table~\ref{tab:baseline} is the whole problem in one place. On the population
\vfour{} was published against, it is a clean policy: \numVthreeDeadAsk\% dead
asks, no game containing a dead run of six, no game reaching the forcing horizon
at all. Against a copy of itself, two in five of its asks name a card it can
prove the target does not hold, two in five are exact repetitions, and a third of
games contain an extended run of guaranteed misses. The two columns differ by a
factor of fourteen in dead-ask rate. A policy panel of weaker opponents cannot
detect this, and a competent human opponent is closer to the mirror column than
\vthree{} is --- which is how the failure was first reported, from live play, and
not from any experiment in the v0.4 study.

The pathology is not caused by the approximate inference path. Substituting the
exact reference belief reduces dead asks from \numDeadAsk\% to
\numBlockDeadAsk\% and leaves the longest dead run at \numBlockDeadRun{}. The
belief reports the hit probability correctly and the score prefers the move
anyway: this is a policy defect, not an inference defect.

\subsection{The deadlock is a two-question cycle}
\label{sec:diag-cycle}

Of \numScanGames{} mirror games scanned in the forensic probe, \numLongGames{}
exceed 300 events. In those games the mean longest dead run is \numCycleLen{}
asks, and the mean number of distinct (actor, card, target) triples inside that
run is \numCycleTriples{}: in \numPureCycle{} of \numLongGames{} the entire run
consists of exactly two questions, alternating, each asked well over a hundred
times. Two seats on opposite teams ask each other the same guaranteed-miss
question until an event counter stops them.

The reason a deterministic policy can do this is that the position is a fixed
point. Sampling each dumped run at its onset and at $+40$ and $+80$ events, every
observer's exact mean per-card certainty is identical to the last digit at all
three samples, as are the hand counts, the unresolved counts and the score. The
only quantity that changes is a counter: each repeated ask appends a fresh
ask-legality certificate that duplicates one already held, so the stored
certificate list grows without bound while the posterior it induces does not
move. Both policies re-derive the same argmax from the same information state,
and the cycle closes.

\subsection{The deadlock is not an information freeze}
\label{sec:diag-information}

The v0.4 termination study attributes non-termination to the locked-half-suit
theorem, stating that such a half-suit is frozen and that, for the same reason,
no further information about its allocation can ever arrive. That account fails
in two independent ways, and \S\ref{sec:corrections} retracts it.

\paragraph{First: the deadlock is mostly not in a locked half-suit.} A ground-truth
lock census taken at \emph{every} event of the dead run --- not only at its ends
--- finds no half-suit locked to any team at any point in the run in
\numNoLockPooled{} of \numNoLockPooledN{} long games pooled over four seeds
(\numNoLockPct\%). Measuring the stronger and more relevant condition, only
\numRunAsksLocked\% of dead-run asks sit inside a half-suit that is locked at that
moment, and only \numCycleInLock\% of long games have their entire cycle inside a
lock. In the majority of deadlocks the theorem has nothing to bite on: at the
onset of one traced run all nine half-suits are live and genuinely split,
\numTraceLive{} of \numTraceLegal{} legal asks are not provably dead, and the
policy ranks a guaranteed miss first and then repeats it \numTraceRepeats{}
times.

\paragraph{Second: where a half-suit \emph{is} locked, information still arrives.}
This is a statement about the rules before it is a measurement. An ask inside a
half-suit locked to the asker's own team is legal --- the set is active, the
target is a live opponent, the asker holds another card of the set and lacks the
one named --- and the deduction machinery publishes the asker's own exclusion
regardless of whether the ask succeeds. ``The asker does not hold this card'' is
public, permanent information about the allocation of a locked half-suit, and
with six cards split among three teammates it frequently resolves a card
outright.

Pooled over three independent seeds and \numCertN{} legal asks inside a
half-suit the asker's team in fact owns: \numCertRate\% strictly increase a
teammate's exact probability of the correct allocation, by a mean of
$+\numCertMean$ and a maximum of $+\numCertMax$. Re-measured with two statistics
that avoid the disputed estimator entirely --- the combinatorial support size of
the half-suit's cards, and the Shannon entropy of the exact marginals ---
\numEntropyRate\% strictly lower the entropy, by a mean of \numEntropyMean{}
nats. Asks \emph{outside} the locked half-suit shrink its support in
0 of \numOutsideSupport{} cases at any seed, which is the expected structural
result and confirms that the certificate lands where it is aimed. Playing such
asks greedily, a team's best allocation probability reaches certainty in
\numLadderRate\% of attempts in a median of \numLadderMedian{} asks.

Ownership monotonicity therefore does not imply informational monotonicity, and
the two must not be merged. The v0.4 \emph{paper} states this correctly as an
observation; the v0.4 termination artifact does not, and it is that artifact's
claim which \S\ref{sec:corrections} withdraws.

\paragraph{The leak, priced.} These certificates are public, so opponents read
them too. Over the same ladders, the team's allocation probability on the
half-suit it is trying to cash rises by $+\numTeamGain$ while the opponents' best
locked-half-suit allocation probability rises by $+\numOppGain$ --- an exchange
rate of about $\numLeakRatio\!:\!1$ in the sender's favour. The opponents' gain
is indirect, arriving through capacity coupling rather than through the locked
half-suit itself, which they can never take.

\paragraph{What the channel costs, and the qualification that matters.} A
guaranteed miss donates the turn. In \vfour{}'s own value function a donated turn
is worth \numDonationCost{} half-suits of final differential, and empirically the
receiving team's take from a miss inside a dead run is \numDonationDead{}
half-suits, with \numDonationNil\% of such donations yielding the opponents
nothing at all.
The qualification is important and was overstated in the first pass of this
analysis: the lock census that selects these asks is a ground-truth statistic. At
the deadlock states examined, the exact probability an owning team assigns to its
own ownership has mean \numOwnProbMean{} and maximum \numOwnProbMax{}, and not
one of the \numOwnProbPairs{} (observer, half-suit) pairs clears the
\numProvableThresh{} bar at which ownership would be provable. A \vfive{} ask rule therefore \emph{cannot} condition on ``a
half-suit my team provably owns''; the same conclusion is reached independently
in \S\ref{sec:diag-channels}, where the provable version of the condition holds
at \numProvableOwn\% of mirror decisions. The channel is real and it must be
priced against ownership probability, not gated on a certainty nobody has.

\subsection{The cause: ownership features that are not gated by hit probability}
\label{sec:diag-features}

The ask score is a twenty-feature linear function refined by a one-ply
expectimax. Four of its features measure ownership progress and none of them is
multiplied by the probability that the ask succeeds: own-set progress
($+\numOwnSetWeight$), team control ($+\numTeamControlWeight$), lock completion
($+\numLockWeight$) and the completion bonus ($+\numCompletionWeight$).
Evaluated at the compiled weights they pay up to $\numOwnershipPay$ at $p = 0$,
against $\numHitWeight \cdot p$ for the hit-probability term. A card the actor's whole team already owns
therefore outscores a genuine chance at a card it does not.

Three further features push the same way. The score's only explicit information
term is the binary entropy of the ask outcome, at weight $\numEntropyWeight$: the
ask whose result would resolve a full bit is penalised most, and the ask whose
result is already known is penalised not at all. This is a negative
value-of-information term in the literal sense --- the policy does not merely
ignore information, it charges for it. A second feature rewards re-asking in the
half-suit last asked in ($+\numRepeatWeight$), which is a direct mechanism for the
\numRepeatAsk\% exact-repeat rate. A third scales the exposure penalty by $(1-p)$
and is therefore maximised exactly when the ask is certain to miss.

Evaluating the compiled weight vector at two representative feature vectors, a
provably dead ask inside a half-suit the team already owns outright scores
\numDeadScore{} while a maximally informative fresh ask scores \numLiveScore{}.
That is an illustration rather than a game measurement --- the shipped policy
adds an expectimax term, so the argmax is not the linear argmax --- but the
decomposition of real decisions agrees with it. At the traced deadlock onsets the
gap between the dead ask ranked first and the best live ask decomposes into the
same three contributions every time: entropy, exposure-on-miss and target hand
size outvote the hit-probability term.

Two facts establish that this is a choice and not a constraint. Starved turns,
at which no live ask exists, are \numStarved\% of asks; \numVoluntaryDead\% of
provably dead asks are made with a live alternative available. And the most
informative legal ask available at a deadlock state ranks, in \vfour{}'s own
score, \numInfoRank{} out of roughly forty on average. At the \numInfoFreeStates{}
deadlock states examined the ask actually played moved no teammate's exact
posterior in every case, while an ask that would have moved it existed in every
case. The ask rule is not slightly mispriced on information; it does not contain
the term.

\paragraph{Causal isolation.} Filtering provably dead asks out of the candidate
set, with nothing else changed, collapses the longest dead run from
\numFilterBaseDeadRun{} to \numFilterDeadRun{}, takes games containing a dead run
of six from \numFilterBaseGames\% to 0\%, takes declarations at or after the
horizon from \numFilterBasePostHorizon{} to zero, takes misdeclarations from \numFilterBaseDeclWrong\% to
\numFilterDeclWrong\%, and scores \numFilterWin\% against the shipped policy ---
neutral. The single change removes the pathology at no measured cost in strength,
which identifies it as the load-bearing defect.

Removing the repetition without pricing information is a different matter, and it
fails. Forbidding exact repeats alone terminates every game, taking events per
game from \numNoRepeatBaseEvents{} to \numNoRepeatEvents{} and the longest dead
run from \numNoRepeatBaseRun{} to \numNoRepeatRun{}, but it loses to the shipped
policy at \numNoRepeatWin\% over \numNoRepeatGames{} games. Breaking the cycle is
not sufficient; the candidate set has to be restricted to asks that can actually
do something.

\subsection{Terminating by force, and terminating by patience}
\label{sec:diag-dilemma}

\vfour{} terminates deadlocks with a graduated escalation keyed on the public
event count. Raising the horizon so that the rule never fires is worth
\numPatientWin\% against the shipped configuration (deal-clustered
\numPatientWinCI\%, $n = \numPatientN$), with the patient side declaring at
\numPatientDeclAcc\% accuracy against the forcing side's \numForcingDeclAcc\%.
But pure patience does not terminate: the unforced mirror reaches the action cap
in \numPatientCap\% of games, with \numPatientDeadAsk\% of asks provably dead and
a longest run of \numPatientDeadRun{}. Table~\ref{tab:dilemma} states the
dichotomy.

\begin{table}[t]
\centering
\caption{The two settings \vfour{} has, and what each costs. Sources:
\texttt{research/v05/results/P0-v04-pathology.md},
\texttt{research/v05/results/P0b-forcing-dilemma.md}.}
\label{tab:dilemma}
\begin{tabular}{lrr}
\toprule
& \textbf{forcing (shipped)} & \textbf{patient} \\
\midrule
declarations wrong                        & \numDeclWrong\%     & \numPatientDeclWrong\% \\
games ending by action-cap adjudication   & 0\%                 & \numPatientCap\% \\
provably dead asks                        & \numDeadAsk\%       & \numPatientDeadAsk\% \\
longest dead run                          & \numDeadRunLongest{} & \numPatientDeadRun{} \\
\bottomrule
\end{tabular}
\end{table}

Note the third row. In the patient mirror more than half of all asks are already
provably dead: the turns are being burned either way. They are simply not being
spent on asks that emit useful certificates. That is the sense in which the
project owner's judgement --- that a risky ask is better than a guaranteed-zero
declaration --- is a statement about the ask policy and not about the
declaration rule.

\paragraph{The second stage of the escalation is an unconditional cash.} Past a
second horizon the rule returns ``declare'' before inspecting the allocation
probability at all, and simultaneously zeroes the team-ownership floor and the
marginal pre-gate. The consequence is measurable in the policy's own stated
confidence: bucketing declarations by the probability the policy itself attached
to them, \numLowConfDecl{} declarations in \numLowConfGames{} games are made at a
self-assessed probability below \numLowConfThresh{} and every one of them is
wrong. Pre-horizon declarations are wrong
\numPreHorizonWrong\%; at or after the horizon, \numHorizonWrong\%. Disabling the
second stage is worth \numStageTwoCost{} points against a mirror opponent and
changes nothing against \vthree{}, the detective, the lockout policy, the hunter
or the bluffer --- differences of at most a third of a point, in both directions.
It is a pure strong-opponent tax.

One hypothesis about this rule did \emph{not} hold and is recorded so that it is
not repeated: the horizon does not fire on healthy long games. Every one of the
\numHorizonGames{} games (\numHorizonReach\% of \numLowConfGames{}) that reached it was genuinely deadlocked,
and raising the deadlock threshold to a run of 80 leaves the split unchanged. The
cost is not when the rule fires but what it does when it fires.

\subsection{The forced endgame: an allocation with posterior probability zero}
\label{sec:diag-forced}

When one team is cardless the other must declare every remaining half-suit. Over
\numForcedGames{} distinct mirror games, all \numForcedN{} forced declarations
are wrong. An independently instrumented replay that shares no code with the
diagnostic probe and recomputes each declaration's correctness from the opening
deal reproduces the result at five further mirror seeds, disagreeing with the
engine's own flag on none of the declarations it re-scored.
This is not a hard decision problem and not an accounting artefact; it is one
mechanical defect with a deterministic consequence. The routine that names the
allocation picks each card's owner by an independent per-card argmax over
ownership marginals, with no capacity constraint. In every observed forced
endgame the truth splits the unresolved cards across two teammates and the argmax
piles them onto one. The named allocation therefore violates the declarer's own
capacity constraint, has exact posterior probability zero, and loses the
half-suit with certainty. All \numForcedN{} violate capacity; all
\numForcedN{} have exact posterior probability zero; the declarer's own stated
confidence is exactly zero in all \numForcedN{}; and the best \emph{feasible}
allocation would have been correct in \numFeasibleRight\% of them. The policy
misses by the minimum possible margin every time --- one card in
\numForcedMissRate\% of them, two in the remainder, never more.

The same defect disables the willingness ladder that selects the declarer. The
ladder thresholds a statistic that the unconstrained allocation drives to a hard
zero, and that statistic is exactly zero in \numRungZeroPairs{} of
\numRungZeroPairs{} surveyed (player, live half-suit) pairs. Four different
ladder shapes, including one with a bottom rung at \numLadderBottomRung{} and one with no rungs at
all, produce bit-identical output. The allocator is the defect; the ladder is
merely inert.

\subsection{The opponent model is two scalars, and one of them is decorative}
\label{sec:diag-prior}

\vfour{} reweights its rules posterior by a policy prior with two coefficients:
a weight on ``this seat asked in this half-suit'' and a weight on ``this seat took
turns but never asked here''. They are identical for all five other players and
never updated during a game, so a deliberately silent opponent is misread by
construction --- which is what the project owner exploited at the table.

Three measured findings qualify that story, and two of them contradict the
hypothesis this study began with.

\begin{enumerate}
  \item \textbf{Deleting the prior is worse, not better.} The policy-agnostic
        configuration loses to the shipped one by \numAgnosticCost{} points
        [\numAgnosticCI] in a paired ablation over a panel of three deceptive
        archetypes and the honest mirror, and by \numPriorDeleteCost{} points
        [\numPriorDeleteCI] against the three deceptive archetypes alone, which
        is the figure a claim about deception should quote. Per cell the effect
        is not uniform at the resolution these banks have: over
        \numPriorCells{} (seed, archetype) cells the ablation is worse in
        \numPriorCellsWorse{}, tied in \numPriorCellsTie{} and reversed in
        \numPriorCellsReversed{} (sign test $p \approx \numPriorSignP$).
        The exposure is over-weighting, not weighting: doubling the pair costs
        \numPriorDoubleCost{} points.
  \item \textbf{The silence channel does not exist.} Deleting the second
        coefficient entirely costs nothing measurable anywhere ---
        \numPhiWin\% [\numPhiCI] against the honest mirror, a paired delta of
        \numPhiPaired{} points over a four-opponent panel, and a fifteen-fold
        sweep that is flat inside its intervals. The algebra explains it: the
        exponent rearranges so that the second term is card-independent, and the
        capacity normalisation then removes it. The channel the owner's manoeuvre
        attacks is carrying no information to begin with.
  \item \textbf{The damage deception does is a signed, cell-local
        miscalibration.} Against a table of feinting opponents, a seat that has
        asked once in a half-suit really holds a given unresolved card of it
        \numFeintTrue\% of the time and the policy believes \numFeintBelief\% ---
        an over-confidence of $+\numFeintBias$, rising at two asks. Against
        honest opponents the same cells are calibrated to within \numFeintCalib{} and
        conservative. The exploit runs through the hard certificate, which the
        soft prior only amplifies; global belief accuracy is essentially
        untouched, mean absolute marginal error moving from \numBeliefErrHonest{}
        under an honest opponent to within $\pm\numBeliefErrTol$ of it under every archetype
        tested.
\end{enumerate}

No archetype converted its corruption of the policy's beliefs into wins: shipped
\vfour{} beats all three (\numVsFeint\%, \numVsSilent\%, \numVsWithholder\%).
The silent figure here is against the hit-probability-capped variant of that
archetype, which is the competent one and the one this subsection's belief
measurements use; \S\ref{sec:results-deception} runs the uncapped variant, on
which both policies score far higher, so the two figures are not the same
measurement. The
owner's live result is therefore not reproduced by commitment deception alone,
and \S\ref{sec:limitations} records what distinguishes the two settings.

\subsection{The value function carries one feature, and the mirror was never fitted against}
\label{sec:diag-value}

The sixteen-feature value model is, to measurement precision, one feature. Its
held-out game-clustered $R^2$ is \numValueRSq{}; the $R^2$ of a bias term plus
the score differential alone is \numValueRSqScore{}, which is higher. Two of the
sixteen features are algebraically the same feature, and five more cancel
identically at every call site --- moving all five to $\pm\numDeadFeatureSwing$ changes none of
\numDeadFeatureDecisions{} declaration decisions. Capacity is not the binding
constraint: a depth-3 gradient-boosted tree on the same rows reaches
\numValueRSqTree{}. And the model is least informative exactly where it is needed
most, reaching only \numValueRSqTail{} in the tail of games past \numTailEvents{}
events.

Splitting the model by use is the actionable result. Deleting the ask-side
expectimax entirely is worth \numAskSideValue{} [\numAskSideCI] --- nothing ---
while reverting declarations to fixed thresholds costs \numDeclSideValue{}
[\numDeclSideCI]. The part of the value model that earns its keep is the
declaration rule; the ask-side use, which is the side the deadlock lives on,
does not.

Finally, the fitting objective could not have caught any of this. The mirror was
absent from the fitting panel in every recorded round, and the recorded
per-opponent win rates never fall below \numPanelFloor{}, so no roughly even opponent was
ever present. The objective is a soft minimum at $\beta = \numSoftMinBeta$, which
over a panel whose win rates span about six points gives a maximum-to-minimum
gradient weight ratio of \numSoftMinRatio{}: it is a weighted mean wearing the
name of a minimum. That is the direct explanation for why a mirror pathology of
this size survived fitting.

\subsection{Channels the policy cannot express}
\label{sec:diag-channels}

Two capabilities are absent rather than mispriced.

\paragraph{The target dimension.} Over \numAskDecisions{} mirror ask decisions,
at \numTargetTwo\% of decisions at least two opponents are hard-indistinguishable
holders of the card actually asked for, carrying a mean of \numTargetBits{} free
bits per ask at a mean hit-probability spread of \numTargetSpread{}; against
\vthree{} the figures are \numTargetTwoVthree\% and \numTargetBitsVthree{} bits.
The dominant refinement term in the ask score discards the target's identity
outright, and every target-dependent quantity anywhere in the score is material
--- hit probability, the target's reply threat, hand size and exposure. Nothing
prices what the choice of target tells a teammate. A separate verification
qualifies the ``free'' in that description: the indistinguishability classes are
defined by hard consistency, and choosing within one is not free once the
capacity-marginal probabilities are compared, so the channel must be priced.
Turn routing is available on the same dimension --- at \numRouteTwo\% of mirror
decisions the actor could hand the turn to any of two or more chosen opponents by
a provably dead ask --- and the policy has one soft penalty for donating the turn
to a dangerous seat and no concept of choosing whom to donate it to.

\paragraph{A knowledge model of anybody else.} The policy builds deduction state
only for itself; the only such objects it ever constructs are clones of its own.
Every technique that turns on what another seat knows --- teammate-directed
signalling, blackballing, choosing the best-informed declarer, baiting --- is out
of reach not because a feature is missing but because the state it would read
does not exist. Since the transcript is public, this is nearly free to add:
seat $j$'s knowledge is the public deduction state plus a posterior over $j$'s
hand.

\subsection{Summary of confirmed defects}
\label{sec:diag-summary}

\begin{table}[t]
\centering
\caption{Confirmed defects, each with the measurement that establishes it and the
mechanism of \S\ref{sec:mechanisms} that addresses it. Rows marked
``\emph{designed}'' are specified but not built.}
\label{tab:defects}
\begin{tabular}{p{0.30\textwidth}p{0.42\textwidth}l}
\toprule
\textbf{Defect} & \textbf{Measured cost} & \textbf{Fix} \\
\midrule
Ownership features not gated by hit probability &
\numDeadAsk\% dead asks, \numVoluntaryDead\% of them voluntary; longest dead run
\numDeadRunLongest{} & M1 \\
Unconstrained per-card allocation guess &
\numForcedWrongRate\% of \numForcedN{} distinct forced declarations wrong; willingness
ladder inert at \numRungZeroPairs{} of \numRungZeroPairs{} pairs & M2 \\
Escalation stage 2 cashes without inspecting probability &
\numStageTwoCost{} points against a mirror opponent, zero against every weaker
style & M8 \\
Negative value-of-information term &
weight $\numEntropyWeight$ on outcome entropy; most informative ask ranks
\numInfoRank{} of $\approx$40 & M3 (\emph{designed}) \\
Refinement term discards the target's identity &
\numTargetBits{} free bits per ask at \numTargetTwo\% of decisions & M5
(\emph{designed}) \\
No knowledge model of other seats &
every ``what does he know'' technique unreachable & M4 (\emph{designed}) \\
Deterministic argmax ask selection &
leakage pinned at the deterministic bound \cite{alvim-leakage};
\numRepeatAsk\% exact repeats & M6 (\emph{designed}) \\
Opponent model is two fixed global scalars &
silence channel measurably inert; feint miscalibration $+\numFeintBias$ & M7
(\emph{designed}) \\
Value model is effectively one feature &
held-out $R^2$ \numValueRSq{} against \numValueRSqScore{} for bias plus score
differential & M9 (\emph{designed}) \\
Mirror absent from the fitting panel &
soft-minimum gradient ratio \numSoftMinRatio{} at $\beta=\numSoftMinBeta$ & M10
(\emph{designed}) \\
\bottomrule
\end{tabular}
\end{table}
%% --- end sections_v05/04-diagnosis ---
%% --- begin sections_v05/05-mechanisms ---
\section{The v0.5 mechanisms}
\label{sec:mechanisms}

\vfive{} is a new policy registered alongside \vfour{}, which stays unchanged and
remains the reference opponent. Every mechanism is individually switchable, so
each can be ablated; the switch names are given with each mechanism. Three are
built and measured (\S\ref{sec:mech-built}); seven are specified and not built
(\S\ref{sec:mech-designed}); four plausible repairs were measured, rejected, and
are recorded here so that they are not rebuilt (\S\ref{sec:mech-rejected}).

\subsection{Built: M1, M2 and M8}
\label{sec:mech-built}

\paragraph{M1 --- live-ask gating and ownership scaling (\texttt{m1}, \texttt{m1p}).}
The candidate set is restricted to asks that are not provably dead in the sense of
Definition~\ref{def:dead}, computed from the actor's own deduction state. Because
the predicate is a hard deduction rather than a posterior quantity, the filter
cannot introduce a modelling error; it removes moves whose probability of success
is exactly zero and nothing else. When the filtered set is empty --- a genuinely
starved turn, \numStarved\% of asks --- the policy falls back to the full legal
set, because the rules still oblige it to move. A second half was designed
alongside it and is \emph{not} in the shipped configuration: under \texttt{m1p},
all four ownership features of \S\ref{sec:diag-features} --- own-set progress,
team control, lock completion and the completion bonus --- are multiplied by the
hit probability, so that a card nobody can hand over cannot outscore a live one
inside the surviving set. The two halves are separately switchable because they
are separate claims: the gate changes which asks are available, the scaling
changes how the survivors are ranked. The gate ships on
(\texttt{liveAskGate}); the scaling ships \emph{off}
(\texttt{ownershipByP}, false in \filepath{engine/src/v05.hpp}), because it was
measured and rejected --- once the gate has removed the $p = 0$ case there is no
dead ask left for the scaling to suppress, and what remains costs
\numMonePReject{} points on top of the other mechanisms at design time and
$\numAblOwnershipPDelta$ at the frozen configuration
(\S\ref{sec:results-ablations}). Every claim in this report about \vfive{}'s play
is a claim about the gate alone.

\paragraph{M2 --- capacity-feasible allocation (\texttt{m2}).} The per-card argmax
of \S\ref{sec:diag-forced} is replaced by a search over allocations that respect
the declarer's own capacity constraints. A half-suit has six cards and a team has
three seats, so the assignment space is at most $3^6 = 729$; the implementation
enumerates it, rejects any assignment that places a card at a seat excluded by a
certificate or that gives a teammate more cards than their remaining hand can
hold, and scores the survivors with the same joint estimator the declaration rule
already uses --- which runs on the deployed approximate belief, so the resulting
probability is an estimate and not an exact posterior value. The search is
exhaustive over the feasible set and is cheaper than the dynamic program it
replaces. The same routine feeds three call sites: the forced-endgame
declaration, the willingness ladder that selects the forced declarer --- which
recovers its function for free, since the statistic it thresholds is no longer
identically zero --- and, under a separate switch, the voluntary declaration
path.

\paragraph{M8 --- termination without misdeclaration (\texttt{stage2},
\texttt{norepeat}).} The second stage of the event-count escalation is off by
default; the first stage, which requires a better-than-even allocation
probability, is the only escalation that remains. With M1 in place the deadlock
that the second stage existed to break is gone. The design intended to put a
structural backstop where the event-count guillotine had been --- a repetition
guard preventing the same agent from asking the same (card, target) pair twice,
applied inside the candidate filter and relaxed before the starved-turn fallback
--- and that guard is implemented, switchable, and \emph{off} in the shipped
configuration (\texttt{repeatGuard}, false in \filepath{engine/src/v05.hpp}): it
costs \numGuardReject{} points against \vfour{} at design time and
$\numAblRepeatGuardDelta$ at the frozen configuration, for the reason
\S\ref{sec:results-ablations} gives. M8 as shipped is therefore a deletion and
nothing else. The design rule behind it is the project owner's, and it is a claim
about the game: a risky ask is better than a declaration the policy itself scores
at zero.

\paragraph{What each is worth alone.} Each mechanism was measured in isolation
against \vfour{} over \numMechDeals{} deals in six seat rotations at seed
\numMechSeed{}, with \vfive{} configured to reproduce \vfour{} exactly as the
control (\numControlWin\%). The results are in Table~\ref{tab:mechanisms}. Two
readings are worth stating plainly. First, the three mechanisms that remove the
pathology are close to strength-neutral in head-to-head play, and the one that is
not --- M8 --- gains by \emph{withholding} a bad declaration rather than by
playing better. Second, the repetition guard on its own is exactly neutral, which
is consistent with \S\ref{sec:diag-features}: with the dead asks already gone,
there is almost nothing left for it to catch.

\begin{table}[t]
\centering
\caption{Single-mechanism win rate against \vfour{}, \numMechDeals{} deals in six
rotations, seed \numMechSeed{}. The control is \vfive{} with every mechanism
disabled, which reproduces \vfour{}. These are single-opponent measurements and
are not a strength claim for the assembled policy; see \S\ref{sec:results}.}
\label{tab:mechanisms}
\begin{tabular}{llr}
\toprule
\textbf{Configuration} & \textbf{Mechanism} & \textbf{win rate} \\
\midrule
control (all off)       & ---                                & \numControlWin\% \\
\texttt{m1}             & live-ask gate only                 & \numMoneAlone\% \\
\texttt{m1p}            & ownership features scaled by $p$   & \numMonePAlone\% \\
\texttt{m2}             & capacity-feasible allocation       & \numMtwoAlone\% \\
\texttt{stage2=0}       & escalation stage 2 removed         & \numMeightAlone\% \\
\texttt{norepeat}       & repetition guard only              & \numGuardAlone\% \\
\bottomrule
\end{tabular}
\end{table}

\paragraph{The pathology gates.} The build order was gated on the three pathology
statistics rather than on win rate, following the standing preference that
robustness is the research question. \S\ref{sec:results} reports the full
comparison; the gates themselves were: after M1, no game may contain a dead run
of six; after M2, forced-endgame accuracy must leave zero; after M8, no game may
reach the action cap and no declaration may be made past the horizon.

\paragraph{The parameters behind this table are not \vfive{}'s.}
Table~\ref{tab:mechanisms} was measured before the refit, with the
\numFitDims{} fitted coordinates still \vfour{}'s, and on the design-time
assembly, which carried both of the mechanisms this subsection has just recorded
as rejected. They had been fitted with the escalation of
\S\ref{sec:diag-dilemma} present as a safety net, so with that net removed that
assembly under-claims: it makes \numVfourParamDecl{}
declarations per game against \vfour{}'s \numVfourDecl{}, and scores
\numVfourParamWin\% [\numVfourParamCI] head-to-head. Most of that shortfall
belongs to the two rejected mechanisms rather than to the coordinates, and the
distinction matters because the figure has been read as a parameter mismatch:
on the same bank, with the same \vfour{} coordinates, the assembly that actually
ships scores \numShipParamSameBank\% [\numShipParamSameCI], which is
\numRejectedCost{} points higher and level with \vfour{}
(\filepath{research/v05/runs/v04vector-in-v05.txt}).
\S\ref{sec:fitting} describes the refit that replaced those
coordinates in the frozen configuration, and every strength figure for \vfive{}
in this report is measured after it; the ablations of
\S\ref{sec:results-ablations} repeat this table's comparisons on the fitted
vector, and its control is correspondingly not \vfour{}.

\subsection{Designed and not built}
\label{sec:mech-designed}

None of the following is on \vfive{}'s decision path. Prototype code exists
outside it for the target-scoring and opponent-model mechanisms, and is used only
to check that the quantities they need can be computed; nothing below has been
measured in play, and each entry states the constraint the diagnosis places on
its eventual construction.

\paragraph{M3 --- a net-information term in the ask score.} Replace the negative
value-of-information term with a signed net term: the information the ask
delivers to teammates, minus the information it delivers to opponents, weighted
separately. This is the secrecy-rate trade-off of \S\ref{sec:related-info}, and
the two filters the engine already produces --- the policy-agnostic posterior and
the policy-weighted one --- are enough to compute both sides. The design
constraint established in \S\ref{sec:diag-information} is that the channel is
almost never \emph{provably} free: it must be priced against the probability that
the team owns the half-suit, which at the relevant states averages
\numOwnProbMean{} and peaks at \numOwnProbMax{}, not
gated on a certainty no observer has. A \vfive{} that signals only when signalling
is free will never signal. M3 requires M4.

\paragraph{M4 --- a knowledge model of the other five seats.} Maintain, for each
seat, the public deduction state plus a posterior over that seat's hand. Nearly
free because the transcript is public, and a precondition for M3, M5 and M7.

\paragraph{M5 --- target-dimension selection.} Score the target explicitly:
lockout value (which opponent must not receive the turn), void progress inside
the half-suit, and, where conventions are enabled, codebook value. The measured
headroom is in \S\ref{sec:diag-channels}, and the verification recorded there
requires that the choice be priced against the capacity-marginal hit probability
rather than treated as free.

\paragraph{M6 --- partner-aware stochastic action selection.}
\label{sec:mech-partner}
Select by softmax over the near-optimal set rather than by argmax. This is
required for any leak penalty to mean anything at all: leakage is a function of
the distribution over actions, so a deterministic policy sits at the
deterministic bound whatever its feature weights \cite{alvim-leakage}. The
temperature and the convention flag are both keyed on the partner regime. With
bot teammates --- the self-play condition, in which partners follow a known
blueprint --- lower temperature and convention-rich play are appropriate, because
the receiver decodes what the sender encodes. With human or unknown teammates the
blueprint assumption fails, and play must be grounded and less predictable,
because an opponent who has watched a deterministic argmax across many games
learns to read it. The two regimes are to be reported separately, and the
self-play configuration is never to be headlined as though it were the one a
human will meet.

\paragraph{M7 --- an online per-seat opponent model.} Replace the two global
scalars with a per-player posterior over a small type library, updated from that
player's public asks and from a \emph{time-local} silence statistic --- own turns
since last asking in this half-suit, and having been asked in it without
replying --- rather than from a whole-game ask count, which
\S\ref{sec:diag-prior} shows is inert. Shape it as a data-biased response
\cite{johanson-dbr}: confidence per information set, decaying to the population
prior where evidence is absent, so that a deliberately silent opponent produces
no evidence and receives the prior instead of being misread. Expose the
exploitation/robustness dial explicitly and report the frontier
\cite{johanson-rnash} rather than a single operating point.

\paragraph{M9 --- the value model, rebuilt or retired.} Two options, to be decided
by measurement: collect rows at declaration points as well as ask points, drop
the degenerate features, and refit jointly with the policy vector so that the
freeze step covers it; or retire it and score declarations directly on the
allocation probability against a fitted threshold. \S\ref{sec:diag-value} shows
that the ask-side use contributes nothing and the declaration-side use does, and
that a tree model finds real structure the linear form cannot express.

\paragraph{M10 --- fitting and evaluation.} Put the mirror in the panel. Replace
the soft minimum, which at the panel's spread is a weighted mean
(\S\ref{sec:diag-value}), with a worst-case or minimax-regret objective over the
style set, and report both the minimum and the regret --- a raw minimum over a
small style set stagnates on the hardest member \cite{beukman-remidi}. Promote
the three pathology statistics --- longest dead run, share of games with a dead
run of six, share of post-horizon declarations that are wrong --- to first-class
metrics gating every commit.

\subsection{Conventions, and which configuration is the headline}
\label{sec:mech-conventions}

The strategy corpus states that pre-agreed conventions are outside this game
\cite{develin}; the rules corpus documents a named convention in real use
\cite{pagat}; and a pair of policies trained together has an implicit
pre-agreement whether or not anyone writes it down. \vfive{} therefore builds the
signalling machinery behind an explicit switch, ships it \emph{off} as the
headline configuration --- which is legal under every reading of the rules we
found --- and publishes the with/without difference as a result in its own right
rather than folding it into a single number.

\subsection{Measured, rejected, and recorded so as not to be rebuilt}
\label{sec:mech-rejected}

\begin{enumerate}
  \item \textbf{A time-varying holding cost in the value model.} The literature
        prescription of \S\ref{sec:related-stopping} --- make waiting strictly
        costly and the stopping problem becomes well posed
        \cite{mguni-actioncost} --- was implemented as a linear penalty in
        elapsed events and swept over a forty-fold range of magnitudes. It is a
        mean-shifter. Events per game fall from \numHoldEventsBefore{} to
        \numHoldEventsAfter{} and the frequency of dead runs from
        \numHoldRunGamesBefore\% to \numHoldRunGamesAfter\%, but the tail does
        not move at all: the 90th percentile of events goes from
        \numHoldPninetyBefore{} to \numHoldPninetyAfter{} and the longest dead
        run from \numHoldLongestBefore{} to \numHoldLongestAfter{}. Declaration
        error rises from \numHoldDeclBefore\% to \numHoldDeclAfter\% and the
        configuration loses \numHoldLoss{} points head-to-head. A constant holding cost is
        also not missing --- the declaration margin already is one --- and in
        deadlocked positions the urgency flag bypasses the value-based
        declaration rule entirely, so the term binds in fewer than one in ten
        late blocked opportunities.
  \item \textbf{Deleting the policy prior.} Worse, not better: by
        \numAgnosticCost{} points [\numAgnosticCI] on a mixed panel of three
        deceptive archetypes and the honest mirror, and by
        \numPriorDeleteCost{} points [\numPriorDeleteCI] against the deceptive
        archetypes alone (\S\ref{sec:diag-prior}). The exposure is
        over-weighting.
  \item \textbf{A turn-transfer willingness ladder.} A genuine multi-candidate
        transfer decision arises \numTransferRate{} times per game and governs
        \numTransferAskShare\% of asks; the unilateral choice already matches an omniscient
        oracle at \numTransferOracle{} of \numTransferOracleN{} of them, and
        handing one team a ground-truth chooser is worth
        $\numTransferValue \pm \numTransferSE$ sets per game over \numTransferGames{}
        paired games --- indistinguishable from zero. The channel is also the leakier of
        the two: an eight-rung ladder transmits \numLadderLeak{} bits per
        candidate against the one bit the rules license, and it transmits them
        mid-game to live opponents. Budget nothing for it.
  \item \textbf{Confidence-ranked declaration arbitration.} Resolving
        simultaneous declarations by lowest seat rather than by confidence costs
        \numArbCost{} percentage points pooled over \numArbGames{} games. The
        rules-legal willingness ladder may be added if it is convenient, but
        nothing should be engineered for it.
\end{enumerate}
%% --- end sections_v05/05-mechanisms ---
% Slots 06, 07 and 09 are reserved, matching the v0.4 numbering: 06 for the
% knowledge model of the other five seats (M4), 07 for the net-information ask
% term (M3/M5), 09 for the evaluation protocol.  They are not yet written.
%% --- begin sections_v05/08-fitting ---
\section{Fitting}
\label{sec:fitting}

The mechanisms of \S\ref{sec:mech-built} change what the ask score has to
discriminate between, and they remove a forcing rule that the declaration
thresholds had been fitted underneath, so the compiled coordinates were
refitted. This section describes that refit, and is deliberately explicit about
how little of it there is: one round, one base seed, and a value model that was
not touched.

It should be read alongside a measurement of what it was worth, because the
refit is not what produces the result of this report, and the length of this
section is not evidence that it is. Carrying \vfour{}'s frozen
vector unchanged, the shipped mechanisms play level against \vfour{} ---
\numShipParamWin\% [\numShipParamCI] over \numShipParamGames{} games at
\numShipParamSeeds{} banks, range \numShipParamRange\%, and
\numShipParamDecl{} declarations per game against \numShipParamVfourDecl{} ---
which is the same wash \S\ref{sec:results} reports for the fitted vector. In the
mirror, on the instrument and seed of \S\ref{sec:results}'s pathology bank, that
same unfitted combination already takes the dead-ask rate to
\numShipParamDeadAsk\%, the longest dead run to \numShipParamDeadRun{}, games
containing a run of six or more to \numShipParamDeadRunGames\%, exact repeats
to \numShipParamRepeat\%, declarations at or after the old horizon to
\numShipParamPostHorizon{}, and declaration error to
\numShipParamDeclWrong\% (\filepath{research/v05/runs/v04vector-in-v05.txt}).
On every row of that instrument the repair is therefore mechanical rather than
parametric --- the one exception being the forced endgame, which fires
\numShipParamForcedN{} times on a bank this size and so cannot separate the two
at all. The pre-refit shortfall reported in \S\ref{sec:mech-built},
\numVfourParamWin\%, is not evidence against this: it was measured on the
design-time assembly, and \numRejectedCost{} of its points belong to two
mechanisms that were subsequently rejected and do not ship.

What the refit does, then, is re-seat \numFitKnobs{} decision knobs that were
fitted underneath a forcing rule which no longer fires. Every figure in
\S\ref{sec:results} is measured at the vector it produced and is conditional on
that vector; none of them is attributable to the refit rather than to the
mechanisms, and this section should not be read as though they were.

\subsection{The parameter vector, and why its offset is derived}
\label{sec:fit-vector}

The fitted vector has \numFitDims{} coordinates, in the same layout as
\vfour{}'s so that a \vfour{} vector can seed a \vfive{} fit: the
\numFitAskFeats{} ask-score weights $w_0,\dots,w_{19}$, followed by
\numFitKnobs{} decision knobs --- the two declaration thresholds, the ask floor,
the patience pool, the opponent-card floor, the two mixing weights, the minimum
team probability, the stopping margin, the two policy-prior coefficients $\theta$
and $\phi$, the lookahead width, and the chain and threat weights. Each knob is
clamped to its own interval as the vector is parsed, so that probabilities stay
in $[0,1]$ and counts stay non-negative integers.

The value model is \emph{not} in this vector, exactly as in \vfour{}. Mechanism
M9 of \S\ref{sec:mech-designed} left open whether to refit it jointly inside the
search or retire it in favour of scoring declarations directly on allocation
probability, and neither was done: the \numFitDims{} policy coordinates were
fitted conditional on sixteen value coefficients that were never re-estimated,
and whose measured behaviour \S\ref{sec:diag-value} describes.

The one structural change is the offset at which the decision knobs begin. It is
now \emph{derived} from the compiled ask-feature count in both places that touch
the flat vector --- the spec parser in \filepath{engine/src/factory.hpp} and the
freezer in \filepath{engine/freeze_config_v05.py} --- rather than written down in
either. The v0.4 pipeline hard-coded the offset at \numAliasOffset{}. When the
ask-feature count later grew to \numFitAskFeats{}, the last two ask weights were
read as the first two decision knobs, and every knob after them shifted: the
configuration the optimiser scored was not the configuration the freezer baked
in. Nothing in the pipeline could notice. Both halves ran without error, and both
produced a policy that plays legal Fish.

\vfive{} therefore carries a mechanical assertion in place of a convention. The
shipped configuration is played against the flat vector that is supposed to
encode it: \texttt{v05} on one side, and on the other \texttt{v05:allparams=}
followed by the frozen vector. Two identical policies meeting on deals each played in
both team orientations must split every deal, so the assertion is that the match
returns exactly one half, and any other value is a mismatch between the two
representations. It \numRoundTripPass{}: \numRoundTrip\% over
\numRoundTripGames{} games at \numRoundTripSeeds{} independent seeds, with a
deal-clustered interval that is a single point
(\filepath{research/v05/runs/roundtrip-assert.txt}).

An assertion that always passes is worth showing to have teeth. Re-encoding the
same vector so that a correct parser reconstructs exactly what a parser at offset
\numAliasOffset{} would have built produces a policy that loses to the shipped
one: it scores \numAliasWin\% [\numAliasCI] over \numAliasGames{} games, short of
even by \numAliasLoss{} points. That is the magnitude of the silent error the
v0.4 pipeline could not detect, and it is larger than most of the differences
this report goes on to measure.

One caveat belongs beside the assertion rather than in a footnote. The freezer
writes ask weights to four decimal places and knobs to five, so what ships is a
rounding of the optimiser's output and not the output itself. Played against the
baked configuration, the full-precision vector scores \numQuantWin\%
[\numQuantCI] over \numQuantGames{} games --- no measurable difference, but not
an identical policy --- so the round trip is asserted, and run, at the precision
that actually ships. The freezer's own convenience output does not currently do
that: it prints a suggested assertion command with every coordinate formatted at
five decimals, which therefore does not return one half. The vector it writes is
correct and the assertion above is run against that; the printed command is
recorded as a defect to fix, because a provenance check that fails silently is
worse than none, and that is the class of defect this whole subsection exists to
prevent.

\subsection{Optimiser and schedule}
\label{sec:fit-cem}

The optimiser is unchanged from \vfour{}: the cross-entropy method with a
diagonal Gaussian over the \numFitDims{} coordinates
(\filepath{engine/src/tuner.hpp}). Each generation samples \numFitPop{}
candidates around the incumbent mean, always evaluating the incumbent itself as
one of them, scores every candidate against every panel member, keeps the
\numFitElite{} best as the elite set, and smooths the mean and the
per-coordinate standard deviation towards the elite statistics. Every candidate
within a generation is evaluated on \emph{identical} seed banks, so
within-generation comparisons are paired and the deal draw cancels; the banks are
redrawn between generations, so the search cannot settle onto one set of deals.

The run was \numFitGens{} generations at \numFitDeals{} deals per panel member
per candidate, each deal played in both team orientations
(\numFitCellGames{} games per cell, \numFitTotalGames{} games in total), from
base seed $\numFitSeed$. The initial mean was \vfour{}'s frozen configuration, so
what follows is a local refinement of \vfour{} and not a search from scratch.
The base seed is recorded in \S9 of \filepath{docs/FISHBOT_V05.md} and in the
artifact manifest; that closes one v0.4 gap cheaply, since the round that
produced \vfour{}'s shipping configuration has no committed script carrying its
base seed and can be inspected but not reproduced.

\subsection{The objective, and the two corrections it makes}
\label{sec:fit-objective}

The quantity optimised is the same soft minimum over an opponent panel
$\mathcal{O}$ that \vfour{} used. For a parameter vector $w$, writing
$\mathrm{wr}_o(w)$ for its win rate against opponent $o$ on the current bank,
\begin{equation}
  \mathrm{score}(w) \;=\; -\frac{1}{\beta}\,
    \log \sum_{o \in \mathcal{O}} \exp\bigl(-\beta\,\mathrm{wr}_o(w)\bigr).
  \label{eq:softmin05}
\end{equation}
Two things about it changed, and \S\ref{sec:corrections} records both as
corrections to the v0.4 study rather than as improvements on it.

\paragraph{The panel now contains an opponent of the agent's own strength.} It is
\numFitPanel{} --- \numFitPanelN{} members, with \vfour{} at the head. No
recorded round of the v0.4 fit contained a self-play or mirror-strength opponent
at any point. \S\ref{sec:diag-value} gives the recorded per-opponent floor over
every committed v0.4 trace; a policy of a candidate's own strength sits near one
half by construction, which is well beneath it. A candidate that
deadlocked against a copy of itself while converting against everything weaker
therefore scored exactly as well as one that did not. The pathology of
\S\ref{sec:diagnosis} was not traded away during fitting; it was never priced.

\paragraph{$\beta$ is \numFitBeta{} rather than \numSoftMinBeta{}.} A soft
minimum is only a minimum if its gradient weights separate. The weight that
\eqref{eq:softmin05} places on opponent $o$ is proportional to
$\exp(-\beta\,\mathrm{wr}_o)$, so the ratio between the largest and the smallest
weight on a panel is exactly $\exp(\beta\Delta)$, where $\Delta$ is the spread of
win rates across that panel. \vfour{}'s selected profile spans
$\Delta = \numSoftMinSpan$ over its \numSoftMinPanelN{} opponents, and at
$\beta = \numSoftMinBeta$ that is a ratio of \numSoftMinRatio{}. Weights lying
within a factor of two of one another describe a weighted mean, not a minimum,
which is why the v0.4 description of the objective as a worst case across
playstyles is corrected rather than repeated. \vfive{}'s panel profile spans
$\Delta = \numFitSpan$, and at $\beta = \numFitBeta$ the ratio is
\numFitRatio{}.

Neither change alone accounts for that. The same $\beta = \numFitBeta$ applied to
\vfour{}'s narrow spread would give a ratio of only \numFitRatioCross{}, and
$\beta = \numSoftMinBeta$ applied to \vfive{}'s spread would give
\numFitRatioAtTen{}. The temperature sharpens the objective; putting an
even-strength opponent in the panel is what gives it something to sharpen.

What is still missing is the objective M10 actually asks for. This remains a soft
minimum with a raised temperature, not an explicit worst case, and not the
minimax-regret objective of \S\ref{sec:mech-designed}. Regret over the style set
is reported in \S\ref{sec:results}; it was never optimised, and \vfive{} does not
come out ahead of \vfour{} on it.

\subsection{The winner's-curse guard, and how to read a per-generation number}
\label{sec:fit-guard}

The score a cross-entropy trace reports for a generation is the best of
\numFitPop{} candidates evaluated on one shared seed bank. It is a maximum over a
noisy sample, so it estimates the winner's luck along with its strength. The
tuner guards against shipping that luck: after the final generation it
re-evaluates two candidates --- the distribution mean, and the single best
candidate seen in any generation --- at one common seed that neither was selected
on, and returns whichever of the two scores higher.

Here it returned the \numFitGuardWinner{}, which scored \numFitGuardMu{} against
the best generation's \numFitGuardBest{}, a gap of \numFitGuardGap{}. The
distribution mean is what \filepath{engine/freeze_config_v05.py} baked into
\texttt{V05Config}. The number to hold next to that is what the losing candidate
had scored in its own generation: \numFitScore{} at generation
\numFitBestGen{}, on a panel profile of \numFitProfile{}. The same vector,
re-evaluated on a bank it had not been selected on, scored \numFitGuardBest{}.

That single comparison is the reading rule for the whole trace, and it is worth
stating as a rule because per-generation numbers are the ones a reader is most
likely to quote. The per-generation scores in
\filepath{research/v05/runs/fit-round1.jsonl} are \emph{selection statistics}:
they rank candidates within a generation against a common bank, which is what
they exist for, and they are not measurements of a candidate's strength. The best
one in this trace overstated its own vector by \numFitCurse{} on the objective's
scale. Only the final common-seed re-evaluation, and the held-out banks of
\S\ref{sec:results}, are measurements.

\subsection{The held-out regression}
\label{sec:fit-heldout}

The best generation's own profile puts \vfive{} at \numFitWorst\% against the
mirror-strength member of the panel. Every seed bank used from
\S\ref{sec:results} onwards is disjoint from the fitting banks, and the
configuration was frozen before any of those runs. On them the same configuration
wins \numHeadWin\% of games against \vfour{} across \numHeadSeeds{} seeds (range
\numHeadRange\%, \numHeadGames{} games), and \numPanelVfiveVfour\%
[\numStyleCIVfour] on the per-style bank. In-panel \numFitWorst\% against
held-out \numHeadWin\% is a regression of \numFitRegression{} points, and the
held-out figure is the one this report uses everywhere the two could be confused.

The regression is what the guard above predicts and does not remove. The guard
protects the \emph{choice} between two candidates from selection bias; it does
nothing about the bias in the in-panel figure itself, which is a maximum taken
over \numFitPop{} candidates on a shared bank in each of \numFitGens{}
generations. An in-panel number from this pipeline should be read as an upper
estimate by construction, and the size of the gap here --- rather than any
argument about it --- is the calibration.

One further scope limit belongs here rather than in \S\ref{sec:limitations}.
Deals were held out; opponent \emph{identities} were not. All \numFitPanelN{} of
the fitting panel's members reappear in the evaluation panel of
\S\ref{sec:results}, so the held-out figures measure generalisation to new deals
rather than to new strategic populations. The three deceptive archetypes of the
deception panel are the exception: none of them was in the fitting panel, and
\vfive{}'s loss against one of them is reported there.

\subsection{What this fit is not}
\label{sec:fit-limits}

\begin{itemize}
  \item \textbf{It is one round.} \vfour{} ran \numVfourFitRounds{} rounds; this
        is \numFitRounds{}, from a single base seed, initialised at \vfour{}'s
        frozen vector. Nothing here distinguishes a converged optimum from the
        nearest local one, and no second round from a different seed was run to
        check that the search lands in the same place twice. Every strength
        figure in \S\ref{sec:results} is conditional on that.
  \item \textbf{The value model was left where it was.} It was fitted
        separately for \vfour{} by ridge regression on self-play decision points
        and then held fixed while the policy weights moved around it, and that
        arrangement is unchanged. \S\ref{sec:diag-value} measures what the model
        contributes and where; the refit did not act on it.
  \item \textbf{The panel is still unbalanced.} As the profile in
        \S\ref{sec:fit-guard} shows, most of its members are dominated by a wide
        margin, and a dominated opponent contributes almost no gradient at any
        temperature. The resolution the objective has is spent largely where the
        outcome was not in question.
  \item \textbf{One fitted coordinate moved into a known exposure.} The policy
        prior's $\theta$, the weight on ``this player asked in this half-suit,''
        rose from \numPriorThetaFour{} to \numPriorThetaFive{}. The diagnosis had
        already established that over-weighting that prior is precisely the
        deception exposure --- while deleting it outright is worse still against
        the deceptive archetypes themselves, by \numPriorDeleteCost{} points
        [\numPriorDeleteCI] --- and the deception panel of
        \S\ref{sec:results} measures the bill: \numDecFeintDelta{} points against
        the archetype that manufactures a false ask-legality certificate. The
        objective could not have seen this, because no deceptive archetype was in
        the fitting panel either.
\end{itemize}
%% --- end sections_v05/08-fitting ---
%% --- begin sections_v05/10-results ---
\section{Results}
\label{sec:results}

% Nine floats in a few pages.  LaTeX's default float fractions reserve so much
% of each page for text that the tables would be pushed several pages past the
% prose that discusses them; the values below let a page carry more table and
% fewer orphaned strips of text.  They are restored at the end of the section.
\renewcommand{\topfraction}{0.92}
\renewcommand{\bottomfraction}{0.75}
\renewcommand{\textfraction}{0.06}
\renewcommand{\floatpagefraction}{0.80}
\setcounter{topnumber}{3}
\setcounter{bottomnumber}{2}
\setcounter{totalnumber}{4}
\setlength{\textfloatsep}{12pt plus 2pt minus 4pt}
\setlength{\floatsep}{10pt plus 2pt minus 3pt}
\setlength{\intextsep}{12pt plus 2pt minus 3pt}

The result of this work is the removal of the failure mode of
\S\ref{sec:diag-baseline}, and it is not a strength jump. Over the
\numVfiveStyleN{} opponents of Table~\ref{tab:perstyle} the two policies' mean
win rates are \numVfiveMean\% and \numVfourMean\%, their worst cases are
\numVfiveWorst\% and \numVfourWorst\%, and \vfour{} is marginally the better of
the two on minimax regret over that style set. The evidence for \vfive{} is
Table~\ref{tab:pathology}, the forced-endgame measurement of
\S\ref{sec:results-forced}, and the gain of $\numDecWithholderDelta$ points
against the deception archetype that motivated the work, replicated at both of
the banks in Table~\ref{tab:deception}. The same table records a smaller and
equally replicated \emph{loss} of $\numDecFeintDelta$ points against a different
deceptive archetype, and the evidence against a strength claim is that loss
together with every row of Table~\ref{tab:h2h} and of Table~\ref{tab:perstyle}.
All of it is reported below at the same weight.

\subsection{The pathology, on the same instrument}
\label{sec:results-pathology}

\begin{table}[t]
\centering
\caption{\texttt{fish pathology} run on the \vfour{} mirror, the \vfive{} mirror
and the cross match. Population: \numPathologyDeals{} deals per column at both
team orientations; seed \numPathologySeed{} for the two mirrors and
\numPathologyCrossSeed{} for the cross match, full rules dialect, both policies
frozen. The cross match is \numCrossGames{} distinct games. The two mirror
columns are \numFourGames{}: in a mirror the two orientations run identical
policies and play byte-identical games, so the instrument's raw counts are each
game counted twice, and every count in those two columns is reported here per
distinct game --- the same halving C3 in \S\ref{sec:corrections} applies to the
forced-endgame figure. Re-running each bank at one orientation returns exactly
half of every count and every rate to the printed digit, so no rate in this
table is affected. The instrument also pools \emph{both} teams, which is why the
cross-match column is a property of the pair and not of either policy: the dead
asks it counts are \vfour{}'s, and \S\ref{sec:results-forced} returns to the
consequence of that pooling for the forced-endgame row. No intervals; these are
counts and rates over one bank per column. Source:
\texttt{research/v05/results/E2-pathology.txt}.}
\label{tab:pathology}
\footnotesize
\begin{tabular}{lrrr}
\toprule
& \textbf{\vfour{} mirror} & \textbf{\vfive{} mirror} & \textbf{cross} \\
\midrule
events per game                        & \numFourEventsMean{}  & \numFiveEventsMean{}  & \numCrossEventsMean{} \\
\quad median / p90 / p99               & \numFourEventsMedian{}\,/\,\numFourEventsPninety{}\,/\,\numFourEventsPninetynine{}
                                       & \numFiveEventsMedian{}\,/\,\numFiveEventsPninety{}\,/\,\numFiveEventsPninetynine{}
                                       & \numCrossEventsMedian{}\,/\,\numCrossEventsPninety{}\,/\,\numCrossEventsPninetynine{} \\
asks per game                          & \numFourAsks{}        & \numFiveAsks{}        & \numCrossAsks{} \\
ask hit rate                           & \numFourHit\%         & \numFiveHit\%         & \numCrossHit\% \\
\midrule
provably dead asks                     & \numFourDeadAsk\%     & \numFiveDeadAsk\%     & \numCrossDeadAsk\% \\
dead runs                              & \numFourDeadRuns{}    & \numFiveDeadRuns{}    & \numCrossDeadRuns{} \\
\quad longest                          & \numFourDeadRunLongest{} & \numFiveDeadRunLongest{} & \numCrossDeadRunLongest{} \\
games with a dead run $\geq 6$         & \numFourDeadRunGames\% & \numFiveDeadRunGames\% & \numCrossDeadRunGames\% \\
starved turns                          & \numFourStarved\%     & \numFiveStarved\%     & \numCrossStarved\% \\
exact repeat (actor, card, target)     & \numFourRepeatAsk\%   & \numFiveRepeatAsk\%   & \numCrossRepeatAsk\% \\
asks inside a half-suit the team owns & \numFourOwnLocked\%   & \numFiveOwnLocked\%   & \numCrossOwnLocked\% \\
\midrule
declarations                           & \numFourDeclN{}       & \numFiveDeclN{}       & \numCrossDeclN{} \\
\quad wrong                            & \numFourDeclWrong\%   & \numFiveDeclWrong\%   & \numCrossDeclWrong\% \\
at or after the forcing horizon        & \numFourPostHorizonDecl{} & \numFivePostHorizonDecl{} & \numCrossPostHorizonDecl{} \\
\quad of those, wrong                  & \numFourPostHorizonWrong\% & --- & --- \\
forced-endgame declarations (pooled)   & \numFourForced{}      & \numFiveForced{}      & \numCrossForced{} \\
\quad of those, wrong                  & \numFourForcedWrong\% & \numFiveForcedWrong\% & \numCrossForcedWrong\% \\
\midrule
games reaching the action cap          & \numFourCapRate\%     & \numFiveCapRate\%     & \numCrossCapRate\% \\
\bottomrule
\end{tabular}
\end{table}

Table~\ref{tab:pathology} is the paper's result. The three statistics promoted to
commit gates in \S\ref{sec:mech-built} --- longest dead run, share of games
containing a dead run of six, and declarations made at or after the forcing
horizon --- read \numFiveDeadRunLongest{}, \numFiveDeadRunGames\% and
\numFivePostHorizonDecl{} on the \vfive{} mirror bank, against
\numFourDeadRunLongest{}, \numFourDeadRunGames\% and \numFourPostHorizonDecl{}
for \vfour{} on the identical bank. The provably-dead ask rate falls from
\numFourDeadAsk\% to \numFiveDeadAsk\%, which is a property of the gate rather
than a measurement: the predicate of Definition~\ref{def:dead} is the same
deduction the filter runs, so a gated policy can record a dead ask only through
the starved-turn fallback, and on this bank the fallback never fired at all
(\numFiveStarved\% of asks).

What is a measurement, and what the gate did not guarantee, is the rest of the
table. Ask hit rate rises from \numFourHit\% to \numFiveHit\%, and ask volume
falls from \numFourAsks{} to \numFiveAsks{} per game: removing a class of moves
that could not succeed did not merely shorten the games, it changed what the
surviving asks are spent on. The event distribution moves at the tail and not
only at the mean, which is the property the swept holding cost of
\S\ref{sec:mech-rejected} failed to obtain --- the 99th percentile falls from
\numFourEventsPninetynine{} events to \numFiveEventsPninetynine{}, and the median
from \numFourEventsMedian{} to \numFiveEventsMedian{}. Declaration error falls
from \numFourDeclWrong\% to \numFiveDeclWrong\%, and the \numFourPostHorizonDecl{}
post-horizon declarations that were \numFourPostHorizonWrong\% wrong have no
counterpart because no game reaches the horizon.

Three rows are worth reading against the natural expectation rather than for it.
Exact repeats fall to \numFiveRepeatAsk\% but not to zero, and should not: cards
move, so re-asking a (card, target) pair can be the best move on the board once a
public transfer has put the card where it was not before. That is the same
observation that sank the repetition guard (\S\ref{sec:results-ablations}). Asks
inside a half-suit the actor's own team already owns fall from \numFourOwnLocked\%
to \numFiveOwnLocked\%, and again not to zero: those asks are legal, are not
provably dead when the actor cannot establish the ownership, and emit ask-legality
certificates. \vfive{} does not choose them for that reason --- the net-information
term M3 is unbuilt --- so the remaining rate is incidental, exactly as
\S\ref{sec:limitations} says of the \vfour{} rate. Finally, the cross-match column
records \numCrossDeadAsk\% dead asks with every dead run of length
\numCrossDeadRunLongest{}: the instrument pools both teams, these asks are
\vfour{}'s, and the run lengths are one because the seat that would have continued
the cycle is now gated.

\subsection{Engine and information-safety verification}
\label{sec:results-verify}

\texttt{fish verify} plays the ordered cross-product of the engine's
\numVerifyPool{}-policy pool --- \numVerifyPairs{} ordered pairs over a budget of
\numVerifyGames{} games --- with the knowledge audit active, and records
\numVerifyViolations{} violations in \numVerifyChecks{} soundness checks: no
derived knowledge state, capacity or certificate was ever contradicted by the
actual deal. Card conservation holds in every game, with \numVerifySetFail{}
set-conservation failures; \numVerifyLimitGames{} games reach the action limit;
and bit-identical replay is \numVerifyDeterminism{}.

The scope of that sweep needs stating, because it is narrower than it looks. The
pool it plays is \vfour{}, \vthree{}, \vtwo{} and the six scripted styles, and it
does \emph{not} contain \vfive{}; the determinism check is run on a
\vfour{}-against-\vthree{} bank. What E1 establishes is that the rules driver, the
public deduction state and the certificate machinery are sound --- which is the
machinery M1's gate is a pure function of, so a sound deduction state is exactly
the precondition the gate needs --- and not that \vfive{}'s own play was audited.
For \vfive{} the corresponding evidence is the action-limit rate, which is
\numFiveCapRate\% for the \vfive{} mirror and \numFourCapRate\% for \vfour{} on
the bank of Table~\ref{tab:pathology}, and is non-zero on \numLimitNonzero{} of
the \numLimitRows{} match rows behind the tables in this section: no result here
is an artefact of the cap ending a game. Sources:
\texttt{research/v05/results/E1-verify.txt}, \texttt{engine/src/main.cpp}.

\subsection{The forced endgame, at volume}
\label{sec:results-forced}

\begin{table}[t]
\centering
\caption{Forced-endgame declarations, per declaring team (E8). Population:
\numForcedEightDeals{} deals in six seat rotations, \numForcedEightGames{} games
per arm, seed \numForcedEightSeed{}, \vfive{} against \vfour{}, both frozen. The
two arms are the two sides of the same matches, scored separately; ``rate'' is
forced declarations per game by that team and ``correct'' is the share of them
that named the true allocation. Declaration totals are all declarations by that
team, forced and voluntary. No intervals: a forced declaration is a rare event
and these are pooled counts. Source:
\texttt{research/v05/results/E8-forced-endgame.txt}.}
\label{tab:forced}
\small
\begin{tabular}{lrrrr}
\toprule
& \multicolumn{2}{c}{\textbf{forced declarations}} & \multicolumn{2}{c}{\textbf{all declarations}} \\
\cmidrule(lr){2-3}\cmidrule(lr){4-5}
& \textbf{rate/game} & \textbf{correct} & \textbf{rate/game} & \textbf{correct} \\
\midrule
\vfour{} & \numVfourForcedRate{} & \numVfourForcedAcc\% & \numVfourDeclRateEight{} & \numVfourDeclAccEight\% \\
\vfive{} & \numVfiveForcedRate{} & \numVfiveForcedAcc\% & \numVfiveDeclRateEight{} & \numVfiveDeclAccEight\% \\
\bottomrule
\end{tabular}
\end{table}

The forced endgame is too rare for the pathology bank of
Table~\ref{tab:pathology} to resolve. Its \vfive{} mirror column counts
\numFiveForced{} of these declarations in \numFiveGames{} games, which is a
denominator that supports no rate. E8 therefore plays
\numForcedEightGames{} games per arm and scores each team separately.
Table~\ref{tab:forced} gives the result. \vfive{} enters the forced endgame
\numForcedRateRatio$\times$ less often --- \numVfiveForcedRate{} declarations per
game against \numVfourForcedRate{} --- and when it does, it names a
capacity-feasible allocation that is correct \numVfiveForcedAcc\% of the time
(\numVfiveForcedRightN{} of \numVfiveForcedN{} forced declarations), against
\numVfourForcedAcc\% for \vfour{} (\numVfourForcedRightN{} of
\numVfourForcedN{}). The denominators are worth carrying: the \vfive{} rate rests
on \numVfiveForcedN{} events, which is enough to separate it from \vfour{}'s and
not enough to place it to the nearest point. E8 does not decompose the two
effects by mechanism, so the attribution that follows is the design's and not
this experiment's: M1 is what removes the extended runs of guaranteed misses
during which a team is stripped of cards while half-suits remain undeclared,
which is how the cardless endgame was being reached, and M2 is what stops the
declarer, once there, naming an allocation that gives a teammate more cards than
that teammate can hold.

\numVfiveForcedAcc\% is not the ceiling. \S\ref{sec:diag-forced} measures the best
capacity-feasible allocation to be correct \numFeasibleRight\% of the time under
the reference posterior. That ceiling was measured on \vfour{}'s mirror
forced-endgame states, not on the states \vfive{} reaches here, so it bounds this
number only to the extent that the two populations resemble each other; read as
an indication rather than a bound, roughly two-fifths of the available accuracy
remains unclaimed. The gap is a live target rather than a residual: M2 ranks
feasible assignments by an estimate computed on the deployed approximate belief,
and it does not attempt the per-seat knowledge model (M4) that would sharpen the
ranking.

\paragraph{Why this is reported per declaring team.} \texttt{fish pathology}
pools both teams' forced declarations into a single count. In the cross match of
Table~\ref{tab:pathology} that pooled figure is \numCrossForced{} declarations,
\numCrossForcedWrong\% wrong, and it is a property of neither policy:
Table~\ref{tab:forced} scores the same kind of event separately by declaring team,
on a far larger bank, and finds \numVfiveForcedAcc\% correct for
\vfive{} against \numVfourForcedAcc\% for \vfour{}, so the pooled cross-match
figure is dominated by \vfour{}'s side of those games. This is the same class of
error corrected at C3 in \S\ref{sec:corrections}, where a mirror run's six seat
rotations were read as six independent samples although they are three deal
shifts times two team orientations and both orientations play byte-identical
games. Forced-endgame accuracy is attributable only when it is reported per
declaring team, which is why Table~\ref{tab:forced}, and not the pathology row,
is the measurement of record.

\subsection{Head-to-head against \vfour{}, on deals never fitted against}
\label{sec:results-h2h}

\begin{table}[t]
\centering
\caption{\vfive{} against \vfour{}, five independent evaluation banks (E3).
Population: \numHeadBankDeals{} deals in six seat rotations per bank,
\numHeadBankGames{} games per row and \numHeadGames{} in total, full rules dialect,
both policies frozen. Intervals are deal-clustered percentile bootstrap. The
deals are held out from the fit of \S\ref{sec:fitting}; \vfour{} itself was a
member of the fitting panel, so the opponent identity is not. Source:
\texttt{research/v05/results/E3-headtohead.jsonl}.}
\label{tab:h2h}
\small
\begin{tabular}{lrr}
\toprule
\textbf{Bank seed} & \textbf{\vfive{} win rate} & \textbf{95\% CI} \\
\midrule
\numHeadSeedOne{}   & \numHeadWinOne\%   & \numHeadCIOne\% \\
\numHeadSeedTwo{}   & \numHeadWinTwo\%   & \numHeadCITwo\% \\
\numHeadSeedThree{} & \numHeadWinThree\% & \numHeadCIThree\% \\
\numHeadSeedFour{}  & \numHeadWinFour\%  & \numHeadCIFour\% \\
\numHeadSeedFive{}  & \numHeadWinFive\%  & \numHeadCIFive\% \\
\midrule
mean of the five    & \numHeadWin\%      & range \numHeadRange\% \\
\bottomrule
\end{tabular}
\end{table}

Across the five banks of Table~\ref{tab:h2h}, \vfive{} wins \numHeadWin\% of
\numHeadGames{} games against \vfour{}, with individual banks spanning
\numHeadRange\% and every interval containing $50\%$. The mean margin is about a
point, which is smaller than the spread between banks, and no claim of a strength
advantage is made from it.

These are not the only banks in this report on which the frozen configuration
meets \vfour{}. Counting the per-style cell of Table~\ref{tab:perstyle}
(\numPanelVfiveVfour\%, the figure the abstract and \S\ref{sec:intro} quote
because it is the cell the worst case of that table is read from), the shipped
row of the ablation bank (\numAblFull\%) and the default dialect of
Table~\ref{tab:dialects-v05} (\numDialectDefault\%), \numHeadBanks{} independent
banks measure the same pairing and they span \numHeadBankSpan\%. That span, and
not any one of its members, is what this report knows about the head-to-head:
it straddles even money, its width exceeds every margin inside it, and its
lowest member, \numHeadBankLow\%, puts \vfive{} below \vfour{}.

\paragraph{The fitting regression, stated plainly.} The generation the fit of
\S\ref{sec:fitting} selected from scores \numFitWorst\% against \vfour{}
in-panel, on the banks it was being optimised on. Held out, after freezing, the
same configuration scores \numHeadWin\%. In-panel \numFitWorst\% versus held-out
\numHeadWin\% is the regression, and the held-out figure is the one reported
here and everywhere else in this paper. Nothing about the in-panel number should
be quoted as a result.

The composition of the margin is worth recording because it is not the
composition one would predict from the pathology table. \vfive{} takes
\numHeadSetsFive{} half-suits per game against \numHeadSetsFour{}, and declares
slightly more often --- \numHeadDeclRateFive{} declarations per game against
\numHeadDeclRateFour{} --- at a slightly \emph{lower} accuracy,
\numHeadDeclAccFive\% against \numHeadDeclAccFour\%. \vfour{} gets
\numHeadDeclErrFour\% of its declarations wrong in these games, against
\numFourDeclWrong\% in its own mirror, because against a \vfive{} opponent the
games end at \numHeadEvents{} events and \vfour{} never reaches the horizon that
forces its bad declarations.
The failure mode of \S\ref{sec:diag-baseline} is a property of the interaction,
not of one seat's rule in isolation, and repairing one side of the match removes
most of it from both.

\subsection{The per-style profile, and the worst case}
\label{sec:results-style}

\begin{table}[t]
\centering
\caption{Per-opponent win rate for both policies against the same
\numVfiveStyleN{}-member style set (E4). Population: \numStyleDeals{} deals in six
seat rotations, \numStyleGames{} games per cell, seed \numStyleSeed{}, full rules
dialect, both policies frozen. $\Delta$ is \vfive{} minus \vfour{} in percentage
points. The \vfour{}-against-\vfour{} cell is $50\%$ by construction, not by
measurement. Minimax regret is computed against the better of the two
\emph{measured} arms on each style, not against a per-style optimum, which was
not computed (\S\ref{sec:limitations}). Source:
\texttt{research/v05/results/E4-perstyle.jsonl}.}
\label{tab:perstyle}
\small
\begin{tabular}{lrrr}
\toprule
\textbf{Opponent} & \textbf{\vfive{}} & \textbf{\vfour{}} & \textbf{$\Delta$} \\
\midrule
\vfour{}      & \numPanelVfiveVfour\%       & \numPanelVfourVfour\%       & $\numDeltaVfour$ \\
\vthree{}     & \numPanelVfiveVthree\%      & \numPanelVfourVthree\%      & $\numDeltaVthree$ \\
\vtwo{}       & \numPanelVfiveVtwo\%        & \numPanelVfourVtwo\%        & $\numDeltaVtwo$ \\
lockout       & \numPanelVfiveLockout\%     & \numPanelVfourLockout\%     & $\numDeltaLockout$ \\
detective     & \numPanelVfiveDetective\%   & \numPanelVfourDetective\%   & $\numDeltaDetective$ \\
diversifier   & \numPanelVfiveDiversifier\% & \numPanelVfourDiversifier\% & $\numDeltaDiversifier$ \\
hunter        & \numPanelVfiveHunter\%      & \numPanelVfourHunter\%      & $\numDeltaHunter$ \\
bluffer       & \numPanelVfiveBluffer\%     & \numPanelVfourBluffer\%     & $\numDeltaBluffer$ \\
random        & \numPanelVfiveRandom\%      & \numPanelVfourRandom\%      & $\numDeltaRandom$ \\
\midrule
\textbf{worst case}          & \textbf{\numVfiveWorst\%} & \textbf{\numVfourWorst\%} & \\
\quad attained on            & \numVfiveWorstOpp{}       & \numVfourWorstOpp{}       & \\
mean over the style set      & \numVfiveMean\%           & \numVfourMean\%           & \\
minimax regret               & \numVfiveRegret{}         & \numVfourRegret{}         & \\
\quad attained on            & \numVfiveRegretOpp{}      & \numVfourRegretOpp{}      & \\
\bottomrule
\end{tabular}
\end{table}

Table~\ref{tab:perstyle} is the table this project's standing evaluation
preference asks for, and it says that on win rate the two policies are level.
\vfive{} is ahead on four styles, behind on four, and identical on one; the
largest gain is $\numDeltaLockout$ against the turn-starvation lockout and the
largest loss $\numDeltaVtwo$ against \vtwo{}. The two means differ by
\numStyleMeanGap{} points, computed from the unrounded cells. On minimax regret \vfour{} is marginally the
better policy, \numVfourRegret{} against \numVfiveRegret{}, its regret attained on
\numVfourRegretOpp{} and \vfive{}'s on \numVfiveRegretOpp{}. Neither figure
supports an ordering.

\paragraph{The worst case, and what \vfour{}'s published one excluded.} The v0.4
study reported that \vfour{}'s lowest win rate against any panel member was
\numVfourVsVthree\%. That panel contained no opponent of \vfour{}'s own strength.
Once one is included, the worst case for both policies falls to about $50\%$:
\numVfiveWorst\% for \vfive{} and \numVfourWorst\% for \vfour{}, both attained on
\numVfiveWorstOpp{}. This is not a defect that a better policy in this family
would repair. A policy's mirror match is a two-team game between identical
strategies, so \vfour{}'s own cell is exactly $50\%$ by construction, and any
policy strong enough that no member of a fixed panel beats it will have its worst
case set by a copy of itself. A worst-case figure quoted over a panel that
excludes the mirror is therefore a statement about the panel. C9 in
\S\ref{sec:corrections} records the scope correction; the operational consequence
is M10's, and it is that the mirror belongs in the panel that the worst case is
computed over as well as in the panel that is fitted against.

\subsection{The deception panel}
\label{sec:results-deception}

\begin{table}[t]
\centering
\caption{\vfive{} and \vfour{} against three deceptive archetypes (E10), at two
independent seed banks. Population: \numDecDeals{} deals in six seat rotations,
\numDecGames{} games per cell, seeds \numDecSeedA{} and \numDecSeedB{}, full
rules dialect, both policies frozen. All three are commitment deceivers: they
follow a fixed deceptive rule for the whole game. \emph{Silent} never asks in the
half-suit it holds most cards of until forced --- the uncapped variant, not the
hit-probability-capped one whose figure \S\ref{sec:diag-prior} quotes;
\emph{feint} preferentially asks in a half-suit it holds exactly one card of,
manufacturing a misleading ask-legality certificate when the material cost is
small; \emph{withholder} is the project owner's own manoeuvre, declining to ask in
a half-suit for several turns after being asked in it while holding other cards of
it. $\Delta$ is the mean over both banks of \vfive{} minus \vfour{}. Source:
\texttt{research/v05/results/E10-deception.md}.}
\label{tab:deception}
\small
\begin{tabular}{lrrrrr}
\toprule
& \multicolumn{2}{c}{\textbf{seed \numDecSeedA{}}} & \multicolumn{2}{c}{\textbf{seed \numDecSeedB{}}} & \\
\cmidrule(lr){2-3}\cmidrule(lr){4-5}
\textbf{Opponent} & \textbf{\vfive{}} & \textbf{\vfour{}} & \textbf{\vfive{}} & \textbf{\vfour{}} & \textbf{$\Delta$} \\
\midrule
silent      & \numDecSilentFiveA\%      & \numDecSilentFourA\%      & \numDecSilentFiveB\%      & \numDecSilentFourB\%      & $\numDecSilentDelta$ \\
feint       & \numDecFeintFiveA\%       & \numDecFeintFourA\%       & \numDecFeintFiveB\%       & \numDecFeintFourB\%       & $\numDecFeintDelta$ \\
withholder  & \numDecWithholderFiveA\%  & \numDecWithholderFourA\%  & \numDecWithholderFiveB\%  & \numDecWithholderFourB\%  & $\numDecWithholderDelta$ \\
\bottomrule
\end{tabular}
\end{table}

The withholder is a direct model of the manoeuvre reported from live play that
started this work: hold cards of a half-suit you have been asked for, then
decline to ask back in it. Against that archetype \vfive{} gains
$\numDecWithholderDelta$ points, and the two banks agree closely on the size of
it --- $\numDecWithholderDeltaA$ and $\numDecWithholderDeltaB$ --- which is the
strongest replication anywhere in this report. It also gains
$\numDecSilentDelta$ against the silent archetype, but there only the sign
replicates: the banks read $\numDecSilentDeltaA$ and $\numDecSilentDeltaB$, and
the first of those is a fraction of the width of a bank, so the silent result
should be read as a direction and not as a magnitude. The mechanism is not a better
opponent model --- M3 through M7 are unbuilt and the two policy-prior scalars are
the same construct in both arms. It is that \vfive{} no longer spends turns on
asks it can prove will miss, so a misleading \emph{absence} of asks has much less
leverage over the position. The manoeuvre works by making the deceiver's holdings
appear to lie elsewhere, and a policy already spending \numFourDeadAsk\% of its
asks on cards it could prove were unavailable was donating turns for that
misdirection to work in.

\paragraph{The loss, and its cause.} Against the feint \vfive{} is
$\numDecFeintDelta$ points worse, and that replicates at both banks too
($\numDecFeintDeltaA$ and $\numDecFeintDeltaB$). The
feint manufactures a false-positive ask-legality certificate rather than
withholding a true one, and the fit raised \texttt{priorTheta} --- the weight on
``this player asked in this half-suit'' --- from \numPriorThetaFour{} to
\numPriorThetaFive{}. \vfive{} therefore weights a manufactured certificate more
heavily than \vfour{} did. This is the exposure the diagnosis predicted:
\S\ref{sec:diag-prior} establishes that the vulnerability is over-weighting the
policy prior, not weighting it at all, and the corresponding removal is worse
still --- deleting both prior scalars costs \numPriorDeleteCost{} points
[\numPriorDeleteCI] on a panel containing no honest opponent, against the
\numAgnosticCost{} points \S\ref{sec:mech-rejected} records on the mixed panel.

\paragraph{The parameter is not identified, and was not tuned.} A sweep of
\texttt{priorTheta} over \numThetaSweepN{} values, $\{\numThetaSweep\}$, at
\numDecSeeds{} seeds does not identify it: the ordering of the settings is not
stable between the two banks and every pairwise difference sits inside the
cluster-bootstrap interval. The fitted value was therefore retained rather than
adjusted, because adjusting it on this evidence would be selection on noise. The
principled repair is M7 --- a per-seat online type posterior with data-biased
shrinkage, so that a certificate from a seat whose type posterior has drifted
toward ``deceptive'' is discounted --- and M7 is specified and not built
(\S\ref{sec:mech-designed}).

\paragraph{Consequence for the worst case.} Across the full
\numTwelveStyles{}-style set --- the \numVfiveStyleN{} of
Table~\ref{tab:perstyle} plus these three --- the lowest cell \vfive{} records
anywhere is \numDecWorst\% against the \numDecWorstOpp{}, on the weaker of the
two deception banks, marginally below its mirror cell of \numVfiveWorst\%.
Which of the two is \emph{the} worst case is not resolved by these banks:
averaged over both deception banks the \numDecWorstOpp{} reads
\numDecFeintFive\%, which is above the mirror rather than below it, and the two
candidates sit within half a point of each other either way. The honest
statement is that \vfive{}'s floor over the twelve styles is about
\numVfiveWorst\%, attained on the mirror or on the \numDecWorstOpp{} depending on
the bank. \vfour{}'s worst case remains its own mirror at \numVfourWorst\%.
Neither policy falls below a coin flip against anything in the set, and no
aggregate over it should be quoted without this table and
Table~\ref{tab:perstyle} beside it.

\subsection{Mechanism ablations}
\label{sec:results-ablations}

\begin{table}[t]
\centering
\caption{Frozen-parameter mechanism ablations (E5). Population:
\numAblFiveDeals{} deals in six seat rotations, \numAblFiveGames{} games per row,
seed \numAblFiveSeed{}, every row played against frozen \vfour{}, full rules
dialect. Every row runs \vfive{}'s fitted parameter vector; only the mechanism
switches change, so the control is \emph{not} \vfour{} --- it is the \vfive{}
vector with the mechanisms disabled, which is why it differs from the
\vfour{}-parameter control of Table~\ref{tab:mechanisms}. The last two rows
\emph{add} a mechanism to the shipped configuration rather than removing one.
``Decl.\ acc.'' is the ablated arm's own declaration accuracy and ``opp.'' is
\vfour{}'s in the same games. Intervals are deal-clustered percentile bootstrap.
Source: \texttt{research/v05/results/E5-ablations.jsonl}.}
\label{tab:ablations-v05}
\footnotesize
\begin{tabular}{llrrrrr}
\toprule
\textbf{Switches} & \textbf{Mechanisms} & \textbf{win rate} & \textbf{95\% CI}
  & \textbf{events} & \textbf{decl.\ acc.} & \textbf{opp.} \\
\midrule
\texttt{m1=0,m2=0,stage2=1} & control          & \numAblControl\%    & \numAblControlCI\%    & \numAblControlEvents{}    & \numAblControlDeclAcc\%    & \numAblControlOppDeclAcc\% \\
\texttt{m1=1}               & M1               & \numAblMone\%       & \numAblMoneCI\%       & \numAblMoneEvents{}       & \numAblMoneDeclAcc\%       & \numAblMoneOppDeclAcc\% \\
\texttt{m2=1}               & M2               & \numAblMtwo\%       & \numAblMtwoCI\%       & \numAblMtwoEvents{}       & \numAblMtwoDeclAcc\%       & \numAblMtwoOppDeclAcc\% \\
\texttt{stage2=0}           & M8               & \numAblMeight\%     & \numAblMeightCI\%     & \numAblMeightEvents{}     & \numAblMeightDeclAcc\%     & \numAblMeightOppDeclAcc\% \\
\texttt{m1=1,m2=1}          & M1 + M2          & \numAblMoneMtwo\%   & \numAblMoneMtwoCI\%   & \numAblMoneMtwoEvents{}   & \numAblMoneMtwoDeclAcc\%   & \numAblMoneMtwoOppDeclAcc\% \\
\texttt{m1=1,stage2=0}      & M1 + M8          & \numAblMoneMeight\% & \numAblMoneMeightCI\% & \numAblMoneMeightEvents{} & \numAblMoneMeightDeclAcc\% & \numAblMoneMeightOppDeclAcc\% \\
\texttt{m2=1,stage2=0}      & M2 + M8          & \numAblMtwoMeight\% & \numAblMtwoMeightCI\% & \numAblMtwoMeightEvents{} & \numAblMtwoMeightDeclAcc\% & \numAblMtwoMeightOppDeclAcc\% \\
shipped \vfive{}            & M1 + M2 + M8     & \numAblFull\%       & \numAblFullCI\%       & \numAblFullEvents{}       & \numAblFullDeclAcc\%       & \numAblFullOppDeclAcc\% \\
\midrule
\texttt{+m1p=1}             & \quad + ownership scaled by $p$ & \numAblOwnershipP\%  & \numAblOwnershipPCI\%  & \numAblOwnershipPEvents{}  & \numAblOwnershipPDeclAcc\%  & \numAblOwnershipPOppDeclAcc\% \\
\texttt{+norepeat=1}        & \quad + repetition guard        & \numAblRepeatGuard\% & \numAblRepeatGuardCI\% & \numAblRepeatGuardEvents{} & \numAblRepeatGuardDeclAcc\% & \numAblRepeatGuardOppDeclAcc\% \\
\bottomrule
\end{tabular}
\end{table}

Every row of Table~\ref{tab:ablations-v05} is a frozen-parameter ablation: the
fitted vector is \vfive{}'s in all of them, and no row was refitted around the
mechanism it removes. A cheap removal may still be doing work a refitted policy
would find another way to do, and an expensive removal may be expensive only
because the surrounding weights were fitted to rely on it. That qualification
applies to every row and is not repeated per row; a matched-budget refit is one of
the experiments \S\ref{sec:limitations} names as not run.

\paragraph{Stage 2 is dead code once M1 is on.} The \texttt{m1=1} and
\texttt{m1=1,stage2=0} rows are identical to the digit in win rate, events,
declaration accuracy and their opponent's declaration accuracy, and so are the
\texttt{m1=1,m2=1} and shipped rows. This is not a coincidence of the bootstrap:
with the live-ask gate on, games end at about \numAblFullEvents{} events and the
event-count escalation's second stage is never reached, so switching it off
changes nothing that happens. The measured cost of \numStageTwoCost{} points that
\S\ref{sec:diag-dilemma} attributes to that rule is a cost it can only levy on a
policy that deadlocks.

\paragraph{The rows with the best win rates are the rows that still deadlock.}
M8 alone is the largest number in the table, \numAblMeight\%, and M2 + M8 the
second largest, \numAblMtwoMeight\%; both play games of about
\numAblMeightEvents{} events against the control's \numAblControlEvents{}, so both
retain the failure mode in full. They win by declining to cash a declaration the
policy itself scores near zero: their own declaration accuracy rises from the
control's \numAblControlDeclAcc\% to \numAblMeightDeclAcc\%, while \vfour{}'s in
the same games falls from \numAblControlOppDeclAcc\% to
\numAblMeightOppDeclAcc\% --- an asymmetry this experiment records without
explaining, since the two arms play games of nearly the same length. This is the
forcing dilemma of
\S\ref{sec:diag-dilemma} measured again inside \vfive{}'s switch set: patience
without termination buys points and keeps the pathology, and the configurations
that terminate (\numAblFullEvents{} events, \numAblFullDeclAcc\% own declaration
accuracy, \numAblFullOppDeclAcc\% for the opponent) are the ones that give the
points back. The mechanism package is bought with win rate, not for it.

\paragraph{Two designed mechanisms were rejected on measurement.} Both are
retained as switches, both ship off, and both are recorded here so they are not
rebuilt. Scaling the four ownership features of
\S\ref{sec:mech-built} by hit probability on top of the gate --- \texttt{m1p},
which is a switch and not part of the shipped configuration --- costs $\numAblOwnershipPDelta$ points relative to the
shipped configuration, and was measured at $-\numMonePReject$ points on top of
M1, M2 and M8 at design time, which is why the switch ships off. Its own
diagnostic points at the reason: it \emph{raises} the hit rate, from
\numAblFullHit\% to \numAblOwnershipPHit\%, while losing games. Once the gate has
removed the $p = 0$ case there is no dead ask left for the scaling to suppress,
so all it does is rescale live candidates against one another, and it buys
immediate hits at the expense of the ownership progress the rest of the linear
score was fitted to trade against.

The repetition guard --- \texttt{norepeat}, forbidding a seat from asking the same
(card, target) pair twice --- costs $\numAblRepeatGuardDelta$ points, and was
measured at $-\numGuardReject$ points on top of M1, M2 and M8 at design time. It
also ships off, and it has an explanation worth stating, because the guard looks
like a structural backstop that could not hurt. The reason it does is that cards
move: forbidding a (card, target) repeat forbids asking for a card the actor has
just watched transfer to that opponent in a public event, which is frequently the
best ask on the board. The cost shows in the hit rate, which falls from
\numAblFullHit\% to \numAblRepeatGuardHit\%, and not in the event count, which is
unchanged at \numAblRepeatGuardEvents{} against \numAblFullEvents{}: the guard is
not preventing a deadlock --- there is none left to prevent --- it is deleting
asks the policy would otherwise land. And the
case it was designed to catch is already covered: a repeat of an ask that missed
is provably dead in the sense of Definition~\ref{def:dead}, so M1 removes it
without a separate rule, which is why the guard measured as exactly neutral when
tested alone in Table~\ref{tab:mechanisms} and negative once M1 was present.

\subsection{Calibration}
\label{sec:results-calib}

\begin{table}[t]
\centering
\caption{Aggregate forecast reliability for \vfive{} (E6). Population: the
calibration bank of \numCalibFiveDeals{} deals, seed \numCalibFiveSeed{},
\vfive{} against \vfour{}, weights frozen; the unit is one forecast and $n$
counts forecasts of that type. Voluntary and forced declarations are pooled.
All entries are point estimates: no clustered uncertainty is computed for Brier
score, log loss or expected calibration error, so small differences between them
are not resolved by these numbers. Source:
\texttt{research/v05/results/E6-calibration-v05.txt}.}
\label{tab:calib-v05}
\small
\begin{tabular}{lrrrrrr}
\toprule
\textbf{Forecast} & \textbf{$n$} & \textbf{Brier} & \textbf{log loss}
  & \textbf{ECE} & \textbf{mean pred.} & \textbf{mean obs.} \\
\midrule
ask succeeds        & \numCalibAskN{}  & \numCalibAskBrier{}  & \numCalibAskLogloss{}  & \numCalibAskECE{}  & \numCalibAskPred\%  & \numCalibAskObs\% \\
declaration correct & \numCalibDeclN{} & \numCalibDeclBrier{} & \numCalibDeclLogloss{} & \numCalibDeclECE{} & \numCalibDeclPred\% & \numCalibDeclObs\% \\
\bottomrule
\end{tabular}
\end{table}

Neither aggregate in Table~\ref{tab:calib-v05} characterises the estimator it
scores, for the reason the v0.4 study gave: both are dominated by one bin. Of
\numCalibDeclN{} declaration forecasts, \numCalibDeclTopShare\% fall in the top
decile at \numCalibDeclTopPred\% predicted against \numCalibDeclTopObs\%
realised, which is most of what the expected calibration error of
\numCalibDeclECE{} measures. The decision-relevant bins run the other way: in
$[0.7, 0.8)$ the declaration forecast is \emph{under}confident,
\numCalibDeclMidPred\% predicted against \numCalibDeclMidObs\% realised on
\numCalibDeclMidN{} forecasts: the estimate the stopping rule thresholds is
conservative in exactly the range where the threshold binds.

The ask forecast is the one this report has a reason to look at, because M1
gates on a hard deduction while the ask score prices everything above it. Its
aggregate is close to unbiased --- \numCalibAskPred\% mean predicted against
\numCalibAskObs\% realised over \numCalibAskN{} forecasts --- and it is
overconfident precisely at the bottom, where the lowest decile predicts
\numCalibAskLowPred\% against \numCalibAskLowObs\% realised on
\numCalibAskLowN{} forecasts, and underconfident in the middle, where $[0.4, 0.5)$
predicts \numCalibAskMidPred\% against \numCalibAskMidObs\% on
\numCalibAskMidN{}. The gate itself is unaffected, since it is a deduction and
not a forecast; what the bottom-decile bias means is that among the asks that
survive the gate, the near-dead ones are still priced somewhat too high, which is
a reason to expect the M3 net-information term to have something to price rather
than a defect this report repairs.

\subsection{Rule dialects}
\label{sec:results-dialects}

\begin{table}[t]
\centering
\caption{\vfive{} against \vfour{} under four rules dialects (E7). Population:
\numDialectDeals{} deals in six seat rotations, \numDialectGames{} games per row,
seed \numDialectSeed{}, both policies frozen. ``Forced rate'' is \vfive{}'s
forced declarations per game and the two accuracy columns are the share of each
team's forced declarations that were correct. The \vthree{} legacy dialect
implies declaration on turn only, so its rows track the no-out-of-turn dialect. In
the default and no-cardless dialects a forced declaration occurs a handful of
times per row, so those two accuracy cells rest on single-digit denominators and
should not be read as rates. Intervals are deal-clustered percentile bootstrap.
Source:
\texttt{research/v05/results/E7-rules.jsonl}.}
\label{tab:dialects-v05}
\footnotesize
\begin{tabular}{lrrrrr}
\toprule
& & & \textbf{forced} & \multicolumn{2}{c}{\textbf{forced correct}} \\
\cmidrule(lr){5-6}
\textbf{Dialect} & \textbf{\vfive{} win} & \textbf{95\% CI}
  & \textbf{per game} & \textbf{\vfive{}} & \textbf{\vfour{}} \\
\midrule
default (\S\ref{sec:rules}) & \numDialectDefault\%    & \numDialectDefaultCI\%    & \numDialectDefaultForcedRate{}    & \numDialectDefaultForcedFive\%    & \numDialectDefaultForcedFour\% \\
no out-of-turn declaration & \numDialectNoOOT\%      & \numDialectNoOOTCI\%      & \numDialectNoOOTForcedRate{}      & \numDialectNoOOTForcedFive\%      & \numDialectNoOOTForcedFour\% \\
no cardless declaration & \numDialectNoCardless\% & \numDialectNoCardlessCI\% & \numDialectNoCardlessForcedRate{} & \numDialectNoCardlessForcedFive\% & \numDialectNoCardlessForcedFour\% \\
\vthree{} legacy & \numDialectLegacy\%     & \numDialectLegacyCI\%     & \numDialectLegacyForcedRate{}     & \numDialectLegacyForcedFive\%     & \numDialectLegacyForcedFour\% \\
\bottomrule
\end{tabular}
\end{table}

The head-to-head margin is insensitive to the dialect: the four rows of
Table~\ref{tab:dialects-v05} span \numDialectDefault\% to \numDialectNoOOT\%,
a range narrower than the width of any one of their intervals, and every interval
overlaps every other. None of them separates from $50\%$; the closest is the
no-out-of-turn row, whose interval reaches down to exactly
$50\%$. The v0.4 study's dialect experiment moved by
several points, but it measured a different pairing on a different bank, so the
two are not comparable; all that is claimed here is that the \vfive{}-against-
\vfour{} margin does not depend on which of these four dialects is played.

What the dialect does change is the population of states the mechanisms act on. Disabling out-of-turn declarations raises the
forced-endgame rate by a factor of about \numDialectForcedRatio{}, from
\numDialectDefaultForcedRate{} declarations per game to
\numDialectNoOOTForcedRate{}, because a team that may only declare on its own
turn arrives at the cardless endgame far more often. On that population
\vfive{}'s forced declarations are correct \numDialectNoOOTForcedFive\% of the
time against \vfour{}'s \numDialectNoOOTForcedFour\%, and \vfour{}'s overall
declaration accuracy falls to \numDialectNoOOTDeclFour\% against \vfive{}'s
\numDialectNoOOTDeclFive\%. The absolute accuracies are much higher than the
\numVfiveForcedAcc\% of Table~\ref{tab:forced} because the states are not the
same states --- a forced endgame reached because declaration was illegal until
now is far more determined than one reached after a full game of asking --- so
the two measurements are not comparable in level. Nor are they comparable in the
size of the gap between the two policies: the gap here, \numDialectNoOOTForcedFive\%
against \numDialectNoOOTForcedFour\%, is less than half the one
Table~\ref{tab:forced} records between \numVfiveForcedAcc\% and
\numVfourForcedAcc\%. What transfers across them is only the sign of that gap,
which is what M2 predicts, and that sign is the only claim made from this row. Ask hit rate also falls sharply under
these dialects, to \numDialectNoOOTHit\% from \numDialectDefaultHit\%, which is a
property of the dialect and not of either policy.

\subsection{Throughput}
\label{sec:results-throughput}

\vfive{} plays \numThroughputFive{} self-play games per second over a timing run
of \numThroughputGamesFive{} games, a wall-clock rate for one run with no
interval. The more useful comparison is the matched one the ablation bank
supplies, because it holds the opponent, the deals and the machine fixed: on that
bank the mechanism-disabled control of Table~\ref{tab:ablations-v05} runs at
\numThroughputPatho{} games per second and the shipped configuration at
\numThroughputFixed{}, a factor of \numThroughputRatioFive$\times$. (The two
figures differ from the self-play number above because they are played against
\vfour{} rather than a mirror.) The failure mode of \S\ref{sec:diag-baseline} is
therefore a compute cost as well as a play defect: the deadlock is played out move
by move and every event in it runs a belief update. The factor exceeds the event
ratio --- \numAblControlEvents{} events per game against \numAblFullEvents{} ---
which is consistent with the certificate accumulation \S\ref{sec:limitations}
notes, but this timing is not instrumented finely enough to isolate that. Every
evaluation in this section completes in minutes at these rates; the fit of
\S\ref{sec:fitting} does not, and its cost is not reported here.

\renewcommand{\topfraction}{0.7}
\renewcommand{\bottomfraction}{0.3}
\renewcommand{\textfraction}{0.2}
\renewcommand{\floatpagefraction}{0.5}
\setcounter{topnumber}{2}
\setcounter{bottomnumber}{1}
\setcounter{totalnumber}{3}
\setlength{\textfloatsep}{20pt plus 2pt minus 4pt}
\setlength{\floatsep}{12pt plus 2pt minus 2pt}
\setlength{\intextsep}{12pt plus 2pt minus 2pt}
%% --- end sections_v05/10-results ---
%% --- begin sections_v05/11-corrections ---
\section{Corrections to the v0.4 study}
\label{sec:corrections}

The v0.4 study \cite{fishbot04} and its supporting artifacts contain claims that
the measurements of \S\ref{sec:diagnosis} contradict. They are stated here as
corrections, in one place, rather than quietly patched. Each entry gives what was
claimed, where, what is now measured, and what replaces it. Where the v0.4
\emph{paper} and a v0.4 \emph{artifact} disagree with each other, that is said
explicitly, because in the load-bearing case the paper is right and the artifact
is wrong.

\paragraph{C1. The information-freeze account of non-termination is withdrawn.}
The v0.4 termination artifact, \filepath{research/v04/results/E11-termination.md},
states that the locked-half-suit theorem implies that such a half-suit is frozen
``and, for the same reason, that no further information about its allocation can
ever arrive,'' and concludes that two policies which both decline to cash an
uncertain half-suit therefore deadlock. It repeats the claim as ``Theorem~1
freezes information as well as ownership,'' and closes by offering that
proposition as a contribution to the rules literature. That is false, and all
three statements are retracted.

It is false at the level of the rules before any measurement: an ask inside a
half-suit locked to the asker's own team satisfies every legality condition, and
the deduction machinery publishes the asker's own exclusion whether the ask
succeeds or not. It is false as measured: \numCertRate\% of such asks strictly
raise a teammate's exact probability of the correct allocation, mean
$+\numCertMean$, maximum $+\numCertMax$, and \numEntropyRate\% strictly lower the
exact marginal entropy, by a mean of \numEntropyMean{} nats
(\S\ref{sec:diag-information}). And it is not the mechanism of the observed
deadlock in any case: in \numNoLockPct\% of long games no half-suit is locked to
any team at any point in the dead run, and only \numCycleInLock\% of long games
have their entire cycle inside a lock. What replaces it is
\S\ref{sec:diag-cycle}: the deadlock is a deterministic two-question cycle at a
fixed point of the public information state, produced by an ask score in which
ungated ownership features outbid the hit-probability term at $p = 0$.

The v0.4 \emph{paper} is not implicated in this error and should not be read as
if it were. Its Observation on locked half-suits states the correct position ---
that ownership freezes while allocation information need not, that asks inside a
locked half-suit stay legal and carry certificates, and that ownership
monotonicity alone establishes neither informational freezing nor
non-termination. The error is confined to the artifact, was flagged in the
engine's own specification document, and was never corrected there.

\paragraph{C2. The reported termination property was purchased, and the price was
not reported.} The v0.4 study reports an action-limit rate of \numVfourLimitRate{}
across its experiments and describes the graduated event-count escalation as the
mechanism. Both statements are accurate. What was not reported is what the
escalation does in the population where it actually fires. In mirror play
\numPostHorizonDecl{} declarations are made at or after the forcing horizon and
\numPostHorizonWrong\% of them are wrong, against a baseline declaration error of
\numDeclWrong\%; bucketed by the confidence the policy itself attaches,
\numLowConfDecl{} declarations in \numLowConfGames{} games are made below a
self-assessed probability of \numLowConfThresh{} and every one is wrong. The escalation's second stage returns ``declare'' before
inspecting the allocation probability at all, and removing it is worth
\numStageTwoCost{} points against a mirror opponent while changing nothing
against any weaker style tested. A zero cap-hit rate against a panel of weaker
policies is therefore not evidence that termination was solved; it is evidence
that the horizon was rarely reached in that panel. The correct statement is that
\vfour{} converts non-termination into misdeclaration, and that the conversion
rate is a strong-opponent tax.

\paragraph{C3. Forced-endgame declarations were reported as inaccurate and
recorded as a known gap; they are a certainty of loss with a mechanical cause.}
The v0.4 evaluation tables carry a forced-declaration accuracy of zero, and the
engine specification records that the allocation routine takes a per-card argmax
with no capacity constraint. The two were never connected. \S\ref{sec:diag-forced}
connects them: over \numForcedGames{} distinct mirror games, all \numForcedN{}
forced declarations violate the declarer's own capacity constraint, all have
exact posterior probability zero, and the declarer's own stated confidence is
exactly zero in every one. The best feasible allocation would have been correct
in \numFeasibleRight\% of them. This is not a hard inference problem being lost;
it is a policy naming an allocation it has itself scored at zero. Two notes
attach to the same measurement, neither of which touches the rate. The count of
\numForcedN{} corrects, downward by half, the figure the diagnosis run first
reported: in a mirror match the six seat rotations are three deal shifts times
two team orientations, both orientations play byte-identical games, and every
such run splits its forced declarations exactly evenly between them. And the
result does not rest on that run alone --- an instrument sharing no code with it,
recomputing each declaration from the opening deal, reproduces it at five further
mirror seeds.

\paragraph{C4. The known gap about the value-function coefficients has the sign
backwards.} The v0.4 specification records that the sixteen compiled value
coefficients are not the ones in the fitting artifact, framed as a defect. The
compiled vector is the better of the two: substituting the fitted vector costs
\numEfourteenRange{} win-rate points, including \numEfourteenCost{} points
against \vthree{}, with the sign stable across nine opponents and three
independent seed banks. Two riders on that figure are corrections to the present
study rather than to the v0.4 one. The size is not uniform across the panel: at a
fresh seed the mirror is the \emph{smallest} effect and the only one whose
cluster-bootstrap interval includes zero, so a narrower per-opponent range stated
in an earlier draft of this analysis, and the reading that the compiled vector's
advantage is largest in mirror play, are both withdrawn. The real defect is
larger and different. The model is
effectively one feature: held-out $R^2$ is \numValueRSq{} for all sixteen and
\numValueRSqScore{} for a bias term plus the score differential alone, two
features are algebraically identical, and five more cancel at every call site ---
moving all five to $\pm\numDeadFeatureSwing$ changes none of \numDeadFeatureDecisions{} declaration
decisions. The v0.4 limitation that no goodness-of-fit figure was reported for
the deployed coefficients is upgraded here to a measurement, and it is not
flattering.

\paragraph{C5. The policy prior is one channel, not two.} The v0.4 study
describes the prior as two coefficients, a weight on having asked in a half-suit
and a weight on having taken turns without asking there, and reports a
frozen-policy ablation for the pair. The second is not an independent channel:
the exponent rearranges so that its term is card-independent, and the capacity
normalisation then removes it. Deleting it costs nothing measurable ---
\numPhiWin\% [\numPhiCI] against the honest mirror, a paired delta of
\numPhiPaired{} points across a four-opponent panel, and a fifteen-fold sweep flat
inside its intervals. The ablation result for the pair stands, and the
attribution of it to two channels does not. This also corrects a hypothesis of
the present study rather than of the v0.4 one: the silence channel that the
project owner's manoeuvre attacks was not misreading silence, it was not reading
it at all.

\paragraph{C6. Asks inside a locked half-suit: both halves of the v0.4 hedge are
now measured, and they point in opposite directions.} The v0.4 limitations
describe these asks as ``a channel the fitted configuration has no term for
\dots\ unexploited here, not shown to be worthless.'' They are not worthless:
\S\ref{sec:diag-information} prices them, and a greedy ladder of them reaches
certainty in \numLadderRate\% of attempts at an information exchange rate of
about $\numLeakRatio\!:\!1$ in the sender's favour. They are also not usable in
the form the hedge implies. Selecting them requires the actor to know its team
owns the half-suit, and that condition is provable from public information at
only \numProvableOwn\% of mirror decisions; at the deadlock states the exact
probability an owning team assigns to its own ownership averages
\numOwnProbMean{} and peaks at \numOwnProbMax{}, so not one of the
\numOwnProbPairs{} pairs measured comes near the \numProvableThresh{} bar. The channel has to be priced against ownership probability rather
than gated on ownership. The first pass of the present study wrote ``provably
owns'' where it meant ``in fact owns,'' and that wording is corrected here.

\paragraph{C7. The eight-level forced-endgame willingness ladder does not
function, and leaks more than the rules license.} The v0.4 dialect table presents
the ladder as a graded willingness mechanism and contrasts it with v0.3's
two-level version. It is inert: the statistic it thresholds is exactly zero in
\numRungZeroPairs{} of \numRungZeroPairs{} surveyed (player, half-suit) pairs, a
consequence of C3, and four ladder shapes including one with no rungs at all
produce bit-identical output. Separately, a ladder of this shape transmits
\numLadderLeak{} bits per candidate, against the roughly one bit of willingness
the rules make public. M2 restores the ladder's function for free; the leak
figure is the reason \S\ref{sec:mech-rejected} declines to extend the same device
to mid-game turn transfer, where the audience still holds cards.

\paragraph{C8. Waiting was free by construction, so observing that waiting
dominates proves nothing.} The v0.4 declaration rule compares the value of acting
against the value of waiting, and the waiting branch is evaluated by passing
all-zero deltas to a value model in which the public event count appears nowhere.
The value of waiting at time $t$ therefore equals the value of waiting at $t+1$
as an identity of the implementation, not as an approximation or a finding. Any
argument that patience weakly dominates in that model is self-fulfilling, and the
only thing that ever broke the tie was the hand-set horizon whose cost C2
records. \S\ref{sec:mech-rejected} reports that the literature's remedy for this
--- a strictly positive holding cost --- was implemented and swept, and that it
shifts the mean without touching the tail, so the correction here is not that
v0.4 chose the wrong cost but that the well-posedness of its stopping problem
came from a move cap rather than from anything in the model.

\paragraph{C9. Scope: the published head-to-head result stands; its
generalisation caveat now has an instance.} \vfour{}'s reported
\numVfourVsVthree\% against \vthree{} is not contradicted by anything here, and
the v0.4 study was explicit that it established nothing about generalisation to
unseen opponent populations. \S\ref{sec:diag-baseline} supplies a concrete
instance of that caveat biting: on the same instrument, against a population the
panel did not contain, the dead-ask rate differs by a factor of fourteen and a
third of games contain an extended run of guaranteed misses. The lesson is about
evaluation design and is acted on in M10 --- the mirror belongs in the fitting
panel and the pathology statistics belong in the commit gate.
%% --- end sections_v05/11-corrections ---
%% --- begin sections_v05/12-limitations ---
\section{Threats to validity}
\label{sec:limitations}

This section is in two parts. \S\ref{sec:lim-named} names the eight limitations
that bound what \S\ref{sec:results} may be read to show, each with the experiment
that would resolve it and none of which has been run. \S\ref{sec:lim-validity}
groups the remaining threats by the kind of validity each bears on. Caveats
stated in full at the point of measurement are cross-referenced rather than
restated.

\subsection{Named limitations, and the experiment each needs}
\label{sec:lim-named}

\paragraph{L1. The strength gain is small, and often within noise.}
Head to head \vfive{} scores \numPanelVfiveVfour\% against \vfour{}
[\numStyleCIVfour] --- $\numDeltaVfour$ points on \vfour{}'s
\numPanelVfourVfour\% mirror baseline, on an interval that contains even money.
The nine-style means are level at \numVfiveMean\% against
\numVfourMean\%, and on minimax regret over that style set \vfour{} is
marginally the better of the two, \numVfourRegret{} against \numVfiveRegret{}.
Most of the per-style differences are a point or
less --- hunter $\numDeltaHunter$, bluffer $\numDeltaBluffer$, random
$\numDeltaRandom$, detective $\numDeltaDetective$, diversifier
$\numDeltaDiversifier$, \vthree{} $\numDeltaVthree$ --- and the largest,
\vtwo{} $\numDeltaVtwo$ and lockout $\numDeltaLockout$, are of a size the
per-cell intervals do not separate. No style in the table is one on which the
two policies' intervals fail to overlap. Nothing in this report should be read as
\vfive{} playing better than \vfour{}; the case for it is the pathology table
and the deception panel. \emph{Resolving experiment:} a paired common-deal
evaluation at an order of magnitude more deals per cell, with a per-style
optimum computed for each style, so that regret is measured against a best
response and not against the better of two arms \cite{xu-mirror}.

\paragraph{L2. No exploitability probe was run against \vfive{}.} The v0.4 study
ran one; this one did not, and the battery of \S\ref{sec:results} contains no
best-response search of any kind. \vfive{} therefore inherits no robustness
property from \vfour{}'s probe and claims none of its own. This is a strict
scope limit rather than a hedge: the deception panel measures performance against
three \emph{fixed} deceptive rules, which is not the same quantity as
performance against a policy that is searching for this policy's worst case.
\emph{Resolving experiment:} a local best-response probe against \vfive{},
reported as an achieved score that is a lower bound within its search class
\cite{lisy-lbr}, run with and without M6, since a deterministic argmax sits at
the deterministic leakage bound whatever its feature weights
\cite{alvim-leakage}.

\paragraph{L3. Mechanisms M3--M7 are specified and unbuilt, and two patches for
them are unmeasured.} \S\ref{sec:mech-designed} specifies the net-information ask
term, the per-seat knowledge model, target-dimension selection, partner-aware
stochastic selection and the online opponent model. None is on \vfive{}'s
decision path. Patches for M4/M5 and for M7 exist in the repository
under \filepath{research/v05/patches/}; neither has been evaluated in play, and
no number anywhere in this report is attributable to either. They are recorded
because unmeasured code that looks finished is a hazard, not because they support
a claim. \emph{Resolving experiment:} apply each patch, measure it alone against
the frozen configuration and then jointly, in the build order of
\S\ref{sec:mechanisms}, refitting after each --- since
\S\ref{sec:mech-built} already shows frozen-parameter ablations understating a
mechanism whose parameters were fitted for a different policy.

\paragraph{L4. Forced-endgame accuracy is still far below the feasible
ceiling.} M2 takes forced-endgame accuracy from \numVfourForcedAcc\% to
\numVfiveForcedAcc\% over \numForcedEightGames{} games per arm, and the measured
ceiling for the best capacity-feasible allocation at those states is
\numFeasibleRight\%. The gap is not closed, and the reason is that feasibility
is exact while the score that ranks the survivors is not: it runs on the
deployed approximate belief. \emph{Resolving experiment:} a calibration study of
M2's allocation estimate against exhaustive enumeration on small reachable
states, to bound the gap between its exact feasibility test and its approximate
score, followed by re-ranking survivors with the exact belief where the state is
small enough to afford it.

\paragraph{L5. A single fitting round.} The configuration evaluated here comes
from \numFitRounds{} round of \numFitGens{} generations
(\S\ref{sec:fitting}); the v0.4 configuration came from
\numVfourFitRounds{} rounds.
The gap between in-panel and held-out performance is already visible in the one
round that was run: the best generation scored \numFitWorst\% against \vfour{}
inside the panel, while the frozen configuration --- the distribution mean the
winner's-curse guard preferred to it (\S\ref{sec:fit-guard}) --- scores
\numPanelVfiveVfour\% against \vfour{} on a bank the fit never saw.
The held-out figure is the one reported, and the difference is a regression to be
expected of a single round rather than an anomaly. \emph{Resolving experiment:}
further rounds on fresh deal banks, each evaluated on a bank the fit never saw,
with the round-by-round in-panel and held-out figures published together so the
regression can be read directly.

\paragraph{L6. The feint exposure is real, replicated, and not addressable by
retuning.} Against the feint archetype \vfive{} scores $\numDecFeintDelta$
points relative to \vfour{}, and does so at both seed banks
(\S\ref{sec:results}). The proximate
cause is identified: the fit raised the weight on ``this player asked here'' from
\numPriorThetaFour{} to \numPriorThetaFive{}, and \S\ref{sec:diag-prior}
establishes that the exposure is over-weighting the policy prior rather than
weighting it at all: against the deceptive archetypes themselves, deleting the
prior outright is worse still, by
\numPriorDeleteCost{} points [\numPriorDeleteCI]. A sweep of that weight over \numThetaSweepN{} values
$\{\numThetaSweep\}$ at two seeds does not identify it: the ordering is unstable
across seeds and every pairwise difference sits inside the cluster-bootstrap
interval, so the fitted value is retained rather than tuned on noise.
\emph{Resolving experiment:} M7, a per-seat online type posterior with
data-biased shrinkage \cite{johanson-dbr}, evaluated against the deception panel
at the seeds used here plus fresh ones, and reported as a frontier rather than an
operating point \cite{johanson-rnash}.

\paragraph{L7. Owner decision D2, partner-aware play, is specified and not
implemented.} The design requires two regimes --- convention-rich,
lower-temperature play with bot teammates, and grounded, less predictable play
with human or unknown teammates --- reported separately, with the self-play
configuration never headlined as the one a human will meet
(\S\ref{sec:mech-partner}). \vfive{} implements neither. It is a single
deterministic policy used in both regimes, and every number in this report is
from the self-play regime. There is no human-teammate evaluation of \vfive{} at
all, which matters most for the one result this study did not obtain by
measurement: the report that a well-trained human outperforms \vfour{} on some
reasoning. \emph{Resolving experiment:} build M6 with the temperature and the
convention flag keyed on the partner regime, then evaluate the two regimes as
separate arms --- bot-teammate self-play, and \vfive{} partnered with a policy it
was not trained with as a stand-in for an unknown teammate --- before any
human-teammate sessions are attempted.

\paragraph{L8. Owner decision D1, conventions behind a flag, has no flag.}
\S\ref{sec:mech-conventions} states the policy: build the signalling machinery
behind an explicit switch, ship it off, and publish the with/without difference
as a result in its own right. The switch does not exist. No codebook, no
convention machinery and no receiver-side decoder were built, so the difference
that was to be published is unmeasured, and the headline configuration is
``conventions off'' only in the trivial sense that there is nothing to turn on.
\emph{Resolving experiment:} implement M5's codebook term behind
\texttt{--conventions}, and report the paired difference between the two settings
on the same deal bank, separately for each partner regime of L7, since a
codebook is worth nothing to a teammate who does not share it.

\subsection{Threats grouped by kind of validity}
\label{sec:lim-validity}

\textbf{Internal validity.} Every ablation in Table~\ref{tab:mechanisms} is a
frozen-parameter ablation: no row was refitted, so no row measures a mechanism's
standalone value under parameters chosen for it. The causal isolation of
\S\ref{sec:diag-features} was performed with a prototype filter applied to the
\vfour{} policy, not with the shipped M1, and the two differ in the fallback rule
and in the ownership scaling. The diagnosis probes are copies of the shipped
policy rather than the shipped policy itself; equality with the original was
checked by reproducing its aggregate statistics to the reported digits after
every instrumentation change, which is a strong check and not a proof of
identity. On the other side of the ledger, the evaluated configuration is the
frozen one: the round-trip assertion that the compiled defaults and the fitted
vector name the same policy passes, and the rules audit reports
\numVerifyViolations{} violations in \numVerifyChecks{} checks with determinism
reproduced.

\textbf{Construct validity.} The provably-dead predicate of
Definition~\ref{def:dead} is a hard deduction and admits no false positives, but
several of the headline statistics that motivate M3 are ground-truth statistics:
the lock census, the ``asks inside a half-suit the team already owns'' rate and
the certificate-ladder measurements all classify half-suits by the true deal, not
by what any observer can establish. \S\ref{sec:diag-information} and
\S\ref{sec:corrections} state the consequence --- the provable version of the
condition holds at \numProvableOwn\% of decisions --- but the reader should not
read those rates as an achievable policy. M2's allocation probability is computed
with the deployed approximate belief, so it is an estimate; the feasibility
constraints it enforces are exact, the score it ranks survivors by is not.
More generally, the posterior throughout this work is exact with respect to the
rules and a uniform prior over consistent deals, which is the policy-agnostic
object; the true posterior weights deals by opponent policy, and the two-scalar
prior examined in \S\ref{sec:diag-prior} is a parametric stand-in for that
correction rather than the correction --- which is also why L6's exposure is a
modelling gap and not a tuning error. Finally, win rate does not capture
declaration reliability, calibration, or the pathology statistics that gate this
work, which is why they are reported separately.

\textbf{External validity.} Every opponent in this study was written inside this
project or ported from an earlier FishBot. The deceptive archetypes of
\S\ref{sec:diag-prior} are commitment deceivers --- they follow a fixed deceptive
rule for the whole game --- and while the withholder is a direct model of the
manoeuvre the project owner reported, it does not reproduce what he actually
did, which was to adapt within a game and across games. The gain of
$\numDecWithholderDelta$ points against it is therefore a result about a fixed
rule, not about a human. The report that a well-trained human outperforms
\vfour{} on some reasoning rests on one player's account of live sessions, is not
a controlled measurement, and is treated here as a hypothesis that generated
experiments rather than as a result. No expert human logs exist for this dialect
and no independently authored engine for it was located, so nothing here tests
the agent against a strategy source outside this project. One point runs the
other way and is worth stating: the three deceptive archetypes, together with
\vtwo{}, the bluffer and the random opponent, were absent from the fitting panel
(\S\ref{sec:fitting}), so the deception panel is a genuinely held-out
measurement of deal bank \emph{and} opponent style, which most of the nine-style
table is not.

\textbf{Statistical validity.} The per-style table is a single seed bank
(\numStyleDeals{} deals in six rotations per cell, seed \numStyleSeed{}), and the
deception panel is two; only the \vfour{} pairing has been run at more
(\numHeadSeeds{} banks, \numHeadDeals{} deals, mean \numHeadWin\%, range
\numHeadRange\%). Several of the diagnosis figures are stable in rate and
unstable in magnitude across seeds, and the pooled values are used above for that
reason. Re-running the information measurement at two fresh seeds reproduced the
rate (\numCertRate\% pooled against \numCertRateOrig\% at the original seed) but
moved the mean improvement from $+\numCertMeanOrig$ to $+\numCertMean$ pooled;
the certificate ladder reached certainty in \numLadderOrigN{} of
\numLadderOrigN{} attempts at the original seed and in \numLadderRate\% pooled,
at a median of \numLadderMedian{} asks rather than two; and the ratio at which a
teammate's best allocation moves toward the truth rather than away from it fell
from $\numMapRatioOrig\!:\!1$ to $\numMapRatioPooled\!:\!1$ pooled. All keep
their sign. The ``information-free at every state'' result of
\S\ref{sec:diag-features} is scoped to \numInfoFreeStates{} states already inside
a repeat cycle and is not a claim about \vfour{}'s asks in general --- indeed
\numOwnLocked\% of all its asks land inside a half-suit its team already owns and
do emit certificates, by accident rather than by choice. Intervals quoted are
deal-clustered percentile bootstrap intervals, which under-cover in evaluations
of this kind \cite{jordan-rleval}.

\textbf{Opponent-population dependence.} Nothing in this report bounds
exploitability, and no such bound is claimed; L2 is the scope statement and
\cite{lisy-lbr} the form the eventual measurement must take. Two of the
mechanisms that would bear on it are unbuilt: while action selection remains a
deterministic argmax, leakage sits at the deterministic bound whatever the
feature weights \cite{alvim-leakage}, so any leak penalty in the ask score is
acting on the wrong variable; and the opponent model remains a fixed global
prior rather than a per-seat one, so the robustness frontier that M7 is supposed
to expose does not yet exist to be reported.

\textbf{Rules-dialect dependence.} All results are for the dialect of
\S\ref{sec:rules}, which is one explicit selection among documented variants
\cite{pagat}. Arbitration by lowest seat is a modelling choice rather than a
rule, and the action cap is a backstop rather than part of the game. The
deadlock measurements in particular are sensitive to the cap only in the patient
configuration, where the cap is what ends the game.

\textbf{Computational and model-class limits.} The ask score remains a
twenty-feature linear function refined by a one-ply expectimax, which caps what
any reweighting can express; M3 and M5 add terms to that class rather than
replacing it. The value model inside it is effectively one feature --- held-out
$R^2$ of \numValueRSq{} for all sixteen against \numValueRSqScore{} for a bias
term plus the score differential alone, where a tree reaches \numValueRSqTree{}
--- so M9's decision to rebuild or retire it is still open, and until it is
taken the declaration rule is scoring on a model whose capacity has been
measured and found wanting. M2's search is exhaustive over at most $3^6$
assignments and returns no candidate at all when an opponent provably holds a
card of the half-suit or when no teammate may hold one, so the declaration path
still needs a defined behaviour at those states. The duplicate-certificate growth
noted in \S\ref{sec:diag-cycle} makes the exact belief progressively more
expensive inside a long cycle without changing what it computes; de-duplication
is an unbuilt engine improvement.

\textbf{Experiments not tied to a single limitation above}, none of which has
been run: a provability oracle over the transportation polytope, to establish how
far the \numProvableOwn\% provable-ownership rate can be widened, which decides
whether M3 can ever gate rather than price; and an adaptive human-like opponent,
or logged human play, to test whether the owner's live result survives a
controlled setting that commitment deceivers do not reproduce.

\textbf{Provenance of the figures in this report.} Of the \numProvUsed{}
distinct quantities quoted in the text, \numProvGenerated{} are emitted directly
from an experiment artifact by \filepath{engine/build_tables_v05.py} and cannot
drift from it; the remaining \numProvTranscribed{} are transcribed out of the
diagnosis reports of \S\ref{sec:diagnosis} rather than regenerated. Transcription
is the weaker guarantee, and it is the mechanism by which a number goes stale
when an artifact is regenerated and the prose is not. It is constrained rather
than eliminated: \filepath{paper/check_provenance.py} requires every transcribed
quantity to sit under a comment naming the artifact it came from and fails the
build if that artifact is absent, and it reports the split above so that this
paragraph cannot itself go stale. Extending the generator to parse the diagnosis
reports would close the gap; it has not been done, and until it is, the
\numProvTranscribed{} transcribed figures carry the same standard of evidence as
the v0.4 study's hand-entered numbers, which is precisely the standard
\S\ref{sec:corrections} criticises elsewhere.
%% --- end sections_v05/12-limitations ---
%% --- begin sections_v05/13-conclusion ---
\section{Conclusion}
\label{sec:conclusion}

\vfour{} was published on the strength of a panel of weaker opponents. Played
against a copy of itself it spends \numDeadAsk\% of its asks on cards it can
\emph{prove} the target does not hold, repeats \numRepeatAsk\% of its asks
exactly, and gets \numDeclWrong\% of its declarations wrong. This report
measures that failure, corrects the account the v0.4 study gave of it, and
reports the policy that removes it. What follows is what the evidence in
\S\ref{sec:diagnosis} and \S\ref{sec:results} supports, stated without
extension.

\paragraph{The deadlock is a fixed point of the ask policy, not a freeze of
information.} In a game where nothing compels a claim, the natural explanation
for two patient policies stalling is that they have run out of things to learn.
That explanation is wrong here twice over. It is wrong about the mechanism: the
observed stall is a deterministic two-question cycle, present in \numPureCycle{}
of \numLongGames{} long games at the forensic seed, and across the pooled seed
banks no half-suit is locked to any team at any point in the cycle in
\numNoLockPct\% of long games. And it is wrong about
information: where a half-suit \emph{is} locked, \numCertRate\% of the asks its
owning team may legally make inside it strictly raise a teammate's exact
probability of the correct allocation, by a mean of \numCertMean{} and up to
\numCertMax{}. Ownership monotonicity does not imply informational monotonicity,
and the two were conflated.

The practical consequence is that the repair is a subtraction, not a forcing
rule. Declining asks whose success probability is exactly zero --- a hard
deduction from the actor's own public knowledge, which cannot introduce a
modelling error --- takes the dead-ask rate from \numFourDeadAsk\% to
\numFiveDeadAsk\%, the share of games containing a run of six consecutive dead
asks from \numFourDeadRunGames\% to \numFiveDeadRunGames\%, and the longest such
run from \numFourDeadRunLongest{} to \numFiveDeadRunLongest{}. It does so with
the event-count escalation that used to force claims switched \emph{off}:
declarations made at or after the old horizon fall from
\numFourPostHorizonDecl{} to \numFivePostHorizonDecl{}, declaration error from
\numFourDeclWrong\% to \numFiveDeclWrong\%, and no game in either arm reaches
the action cap. Forcing claims was never the fix; it was the mechanism that
converted the stall into misdeclaration, at a measured cost of
\numStageTwoCost{} points against a mirror opponent.

\paragraph{A per-card argmax allocation is infeasible, not merely unlikely.}
The second defect is of a different kind from a low-probability guess, and the
difference is worth stating exactly. \vfour{} named its forced-endgame
allocation by taking, card by card, the teammate with the highest marginal ---
a procedure with no capacity constraint, which can therefore give a teammate
more cards than that teammate holds. All \numForcedN{} such declarations over
\numForcedGames{} distinct mirror games were wrong, and the reason is not that
the estimate was poor. The named allocation has posterior probability exactly
zero: it is excluded by the declarer's own arithmetic, not merely disfavoured by
it. A policy losing a hard inference problem and a policy naming an outcome it
has itself ruled out are different failures and need different repairs. Searching
the capacity-feasible set instead takes forced-endgame accuracy from
\numVfourForcedAcc\% to \numVfiveForcedAcc\% over \numForcedEightGames{} games
per arm --- against a measured ceiling of \numFeasibleRight\% for the best
feasible allocation at those states, so the gap is narrowed and not closed. The
state is also reached \numForcedRateRatio$\times$ less often, but that belongs to
M1 rather than to M2: what was stripping a team of cards while half-suits stood
undeclared was the runs of guaranteed misses, not the declaration rule
(\S\ref{sec:results-forced}).

\paragraph{What the v0.4 study said that is withdrawn.}
\S\ref{sec:corrections} is the register. The load-bearing entry is C1, the
information-freeze account of non-termination, which is retracted with the
measurement that refutes it; the v0.4 \emph{paper} states the correct position
and is not implicated, while the supporting artifact states the wrong one. The
remaining entries correct the reported termination property (purchased, and the
price unreported), the status of the forced-endgame accuracy figure (a
mechanical certainty, not a hard problem lost), the sign of the value-function
``known gap,'' the claim that the policy prior is two channels, the usability of
asks inside a locked half-suit, the willingness ladder, the well-posedness of the
stopping problem, and the scope of the published worst case. Three of those
entries also carry corrections to earlier drafts of the present study rather
than to the v0.4 one, and are recorded in the same place for the same reason.

\paragraph{What \vfive{} costs, and what it does not buy.} It does not buy win
rate. Over the nine-style panel of \S\ref{sec:results}
(\numStyleDeals{} deals in six rotations per cell, seed \numStyleSeed{}) the
means are level --- \numVfiveMean\% for \vfive{} against \numVfourMean\% for
\vfour{} --- and head to head \vfive{} scores \numPanelVfiveVfour\%
[\numStyleCIVfour], which is $\numDeltaVfour$ points on \vfour{}'s
\numPanelVfourVfour\% mirror baseline on an interval that contains even money.
Across \numHeadSeeds{} further seed banks the same pairing averages \numHeadWin\%
with a range of \numHeadRange\%. On minimax regret over the style set \vfour{} is
marginally the better of the two, \numVfourRegret{} against \vfive{}'s
\numVfiveRegret{}. Per style the differences run in both directions:
$\numDeltaLockout$ on the lockout opponent, $\numDeltaVfour$ on the mirror and
$\numDeltaDiversifier$ on the diversifier against $\numDeltaVtwo$ on \vtwo{} and
$\numDeltaVthree$ on \vthree{}, with the remainder at half a point or less and no
per-cell interval separating either policy from the other on any style.
\vfive{}'s worst case over the nine styles is \numVfiveWorst\%, against \vfour{},
and \vfour{}'s is \numVfourWorst\%, against a copy of itself. In both cases the
worst cell is the opponent of the policy's own strength, which is the worst case
for any policy in this family, and which is why the v0.4 study's floor of
\numVfourVsVthree\% describes a panel rather than a policy. The case for \vfive{} is the pathology table and the
deception panel below. It is not the win rate, and this report does not offer
one.

\paragraph{The result that speaks to the incident that started this.} The work
was prompted by a live report: the project owner held cards of a half-suit he
had been asked for and then declined to ask back in it, and \vfour{} was visibly
confused. Modelled as an archetype --- after being asked in half-suit $S$ while
holding other cards of $S$, avoid asking in $S$ for the next several turns ---
\vfour{} wins \numDecWithholderFour\% of its games against that archetype and
\vfive{} wins \numDecWithholderFive\%, a gain of $\numDecWithholderDelta$ points
replicated at \numDecSeeds{} independent seed banks of \numDecGames{} games
each. The silent archetype moves the same way, by $\numDecSilentDelta$ points.

The mechanism is worth being precise about, because it is not the one the design
document proposed. \vfive{} has no opponent model that \vfour{} lacks: M3--M7
are specified and unbuilt, the policy prior is still fixed and global rather than
per-seat, and nothing in \vfive{} watches a seat and updates a type. What changed is that
\vfive{} never spends a turn on an ask it can prove will miss. A misleading
\emph{absence} of asks works by making a policy's own inferences point at the
wrong half-suit; a policy whose candidate set is already pruned to asks with
strictly positive hard-consistent probability gives that misdirection much less
to move. Robustness here was bought by removing a defect, not by adding a
model --- which is the more encouraging of the two possibilities, and also the
one with a lower ceiling.

The same panel records a loss. Against the feint --- which manufactures a
\emph{false} ask-legality certificate rather than withholding a true one ---
\vfive{} scores \numDecFeintFive\% against \vfour{}'s \numDecFeintFour\%, a
difference of $\numDecFeintDelta$ points that replicates at both seeds. It is
not an accident of tuning noise: the fit raised the weight on ``this player asked here''
from \numPriorThetaFour{} to \numPriorThetaFive{}, and the diagnosis had already
established that over-weighting the policy prior is exactly this exposure ---
while deleting the prior outright is worse still against the deceptive
archetypes themselves, by \numPriorDeleteCost{} points
[\numPriorDeleteCI]. Across the full twelve-style set \vfive{}'s worst case is the
\numDecWorstOpp{} at \numDecWorst\%, marginally below its mirror. No aggregate
figure in this report should be quoted without that table beside it.

\paragraph{What comes next, in order.} The ranking follows from the exposure
above rather than from the design document's build order.

\begin{enumerate}[leftmargin=1.4em,itemsep=1pt,topsep=2pt,parsep=0pt]
\item \textbf{M7, the online per-seat type posterior.} This is the principled fix
      for the feint, and \S\ref{sec:limitations} shows why nothing cheaper will
      do: a sweep of the prior weight over \numThetaSweepN{} values at two seeds
      does not identify the parameter. Shaped as a data-biased response
      \cite{johanson-dbr}, a manufactured certificate from a seat whose type
      posterior has drifted toward deceptive is discounted rather than believed,
      and the exploitation/robustness dial becomes explicit and reportable
      \cite{johanson-rnash} instead of being frozen at a fitted constant.
\item \textbf{M4, M5 and M3, in that order}, because each is the precondition of
      the next. A per-seat knowledge model is nearly free from a public
      transcript; it makes the target dimension scoreable, which is currently
      discarded outright --- the ask score opens by voiding its target argument,
      throwing away a mean of \numTargetBits{} free bits per ask, with at least
      two hard-indistinguishable targets on offer at \numTargetTwo\% of
      decisions and roughly \numTargetBitsGame{} bits per team per game in
      total; and only
      then can the net-information term of M3 be priced rather than gated, which
      \S\ref{sec:diag-information} shows is the only form in which it can ever
      fire.
\item \textbf{An exploitability probe.} The v0.4 study ran one; this one
      did not, and until it does \vfive{} claims no robustness property that a best
      response could contradict. Such a probe reports an achieved score that is
      a lower bound within its search class \cite{lisy-lbr}, and it is worth
      running both with and without M6, since a deterministic argmax sits at
      the deterministic leakage bound whatever its weights
      \cite{alvim-leakage}.
\item \textbf{A value model that is more than one feature.} Held-out $R^2$ is
      \numValueRSq{} for all sixteen features and \numValueRSqScore{} for a bias
      term plus the score differential alone, while a tree reaches
      \numValueRSqTree{}. Either rebuild it on rows collected at declaration
      points as well as ask points and refit it jointly with the policy vector,
      or retire it and score declarations directly on allocation probability.
      The measurement decides; the present state --- sixteen coefficients
      carrying one feature's worth of signal --- should not survive either way.
\end{enumerate}

\paragraph{The narrow claim, and the broader one.} The narrow claim of this
report is that a specific, reproducible failure of a specific agent was
diagnosed to an arithmetic cause in its ask score, that the cause was removed,
that the removal cost nothing measurable in aggregate win rate and left \vfour{}
marginally ahead on minimax regret, and that it helped substantially against one
class of deceptive opponent while hurting measurably against another.
The broader observation is
about evaluation. Every statistic in \S\ref{sec:diagnosis} was available to the
v0.4 study; none of it was collected, because the panel it was measured against
contained no opponent of its own strength, and the failure is a strong-opponent
phenomenon. The mirror belongs in the fitting panel, the pathology counters
belong in the commit gate, and the per-opponent breakdown with an explicit worst
case belongs in the abstract --- not because they are more rigorous, but because
an aggregate win rate against weaker opponents would have reported all of this
as a success.
%% --- end sections_v05/13-conclusion ---

%% --- begin sections_v05/bibliography ---
\begin{thebibliography}{99}

%% ---------------------------------------------------------------
%% Fish and Literature
%% ---------------------------------------------------------------

\bibitem{fishbot03}
FishLab Research Project, ``FishBot v0.3: A Count-Conditioned, Empirically Tuned
Policy for Six-Player Canadian Fish,'' project technical report, 2026. Project
repository: \url{https://github.com/dylann29/fish-optimization}

\bibitem{fishbot04}
D. Nguyen, FishLab Research Project, ``FishBot v0.4: Observer-Conditioned Deal
Inference and Value-Based Declaration in Six-Player Canadian Fish,'' project
technical report, 2026. Project repository:
\url{https://github.com/dylann29/fish-optimization}

\bibitem{pagat}
J. McLeod, ``Literature,'' \emph{Pagat}, 2006--2025, last updated 1 July 2026.\\
\url{https://www.pagat.com/quartet/literature.html}

\bibitem{develin}
M. Develin, ``Canadian Fish,'' chapter 9 of \emph{Card Games}, author's online
manual.\\
\url{https://web.archive.org/web/20170720075756/http://bantha.org/~develin/cardgames.html\#ch9}

\bibitem{somani}
N. Somani, ``literature: card game implementation and learning bot,'' GitHub
repository, MIT licence, Python, 2018--.\\
\url{https://github.com/neelsomani/literature}

\bibitem{playstore}
``Literature: Fish Card Game'' (\texttt{com.cards.game.literature}), Google Play
store listing, closed source.\\
\url{https://play.google.com/store/apps/details?id=com.cards.game.literature}

\bibitem{rlcard}
D. Zha, K.-H. Lai, S. Huang, Y. Cao, K. Reddy, J. Vargas, R. Wei, J. Guo, and
X. Hu, ``RLCard: A Toolkit for Reinforcement Learning in Card Games,''
arXiv:1910.04376, 2019.

%% ---------------------------------------------------------------
%% Voluntary stopping, wars of attrition, strategic delay
%% ---------------------------------------------------------------

\bibitem{christensen-dynkin}
S. Christensen, K. Lindensj\"o, and B. A. Neumann, ``Markovian Randomized
Equilibria for General Markovian Dynkin Games in Discrete Time,''
arXiv:2307.13413v2, 12 August 2025. \url{https://arxiv.org/abs/2307.13413}

\bibitem{hamadene-dynkin}
S. Hamad\`ene, M. Hassani, and M.-A. Morlais, ``$\epsilon$-Nash Equilibria of a
Multi-player Nonzero-sum Dynkin Game in Discrete Time,'' arXiv:2201.03562, 2022;
\emph{Dynamic Games and Applications}, DOI 10.1007/s13235-023-00500-3.

\bibitem{hendricks-attrition}
K. Hendricks, A. Weiss, and C. A. Wilson, ``The War of Attrition in Continuous
Time with Complete Information,'' \emph{International Economic Review}
29(4):663--680, 1988.

\bibitem{mguni-actioncost}
D. Mguni, ``Stochastic Games with Minimally Bounded Action Costs,''
arXiv:2407.18010, 2024. \url{https://arxiv.org/abs/2407.18010}

\bibitem{admati-perry}
A. R. Admati and M. Perry, ``Strategic Delay in Bargaining,'' \emph{Review of
Economic Studies} 54(3):345--364, 1987.

%% ---------------------------------------------------------------
%% Information, leakage, and information-seeking actions
%% ---------------------------------------------------------------

\bibitem{wyner}
A. D. Wyner, ``The Wire-Tap Channel,'' \emph{Bell System Technical Journal}
54(8):1355--1387, 1975.

\bibitem{farrell-gibbons}
J. Farrell and R. Gibbons, ``Cheap Talk with Two Audiences,'' \emph{American
Economic Review} 79(5):1214--1223, 1989.

\bibitem{alvim-leakage}
M. S. Alvim, K. Chatzikokolakis, Y. Kawamoto, and C. Palamidessi, ``Information
Leakage Games,'' \emph{GameSec 2017}, LNCS 10575:437--457; journal version
\emph{ACM Transactions on Privacy and Security} 25(3):20, 2022;
arXiv:1705.05030.

\bibitem{lervold-infogain}
N. Lervold, G. L. Peterson, and D. W. King, ``Incentivizing Information Gain in
Hidden Information Multi-Action Games,'' \emph{Computers and Games (CG 2022)},
LNCS 13865, DOI 10.1007/978-3-031-34017-8\_6. Venue and finding verified from the
publisher abstract; the full text was not obtainable and no formula is taken from
it.

\bibitem{serrino-deeprole}
J. Serrino, M. Kleiman-Weiner, D. C. Parkes, and J. B. Tenenbaum, ``Finding
Friend and Foe in Multi-Agent Games,'' \emph{NeurIPS 32}, 2019.

%% ---------------------------------------------------------------
%% Online opponent modelling and robustness to deception
%% ---------------------------------------------------------------

\bibitem{ganzfried-multiplayer}
S. Ganzfried, K. A. Wang, and M. Chiswick, ``Opponent Modeling in Multiplayer
Imperfect-Information Games,'' arXiv:2212.06027, 2022 (rev.\ 2024).

\bibitem{rosman-bpr}
B. Rosman, M. Hawasly, and S. Ramamoorthy, ``Bayesian Policy Reuse,''
\emph{Machine Learning} 104(1):99--127, 2016; arXiv:1505.00284.

\bibitem{smith-qmixing}
M. O. Smith, T. Anthony, Y. Wang, and M. P. Wellman, ``Learning to Play against
Any Mixture of Opponents,'' arXiv:2009.14180, 2020; \emph{Frontiers in Artificial
Intelligence}, 2023.

\bibitem{johanson-rnash}
M. Johanson, M. Zinkevich, and M. Bowling, ``Computing Robust
Counter-Strategies,'' \emph{NIPS 20}, 2007.

\bibitem{johanson-dbr}
M. Johanson and M. Bowling, ``Data Biased Robust Counter Strategies,''
\emph{AISTATS 2009}, 2009.

\bibitem{ganzfried-safe}
S. Ganzfried and T. Sandholm, ``Safe Opponent Exploitation,'' \emph{ACM EC 2012};
\emph{ACM Transactions on Economics and Computation}, 2015.

\bibitem{muller-bestofboth}
A. M\"uller, J. Schneider, S. Skoulakis, L. Viano, and V. Cevher, ``Best of Both
Worlds: Regret Minimization versus Minimax Play,'' arXiv:2502.11673, 2025.

\bibitem{yang-bayestomop}
T. Yang, Z. Meng, J. Hao, C. Zhang, Y. Zheng, and Z. Zheng, ``Towards Efficient
Detection and Optimal Response against Sophisticated Opponents,'' \emph{IJCAI
2019}; arXiv:1809.04240.

\bibitem{milec-abd}
D. Milec, V. Kova\v{r}\'ik, and V. Lis\'y, ``Adapting Beyond the Depth Limit:
Counter Strategies in Large Imperfect Information Games,'' arXiv:2501.10464,
2025.

%% ---------------------------------------------------------------
%% Worst-case and minimax-regret objectives
%% ---------------------------------------------------------------

\bibitem{aggarwal-prmre}
N. Aggarwal and J. P. How, ``Beyond Bayesian Nash: Learning Minimax-Regret
Equilibria for Adversarial Team Games under Asymmetric Information,''
arXiv:2607.09993, 2026; \emph{Transactions on Machine Learning Research}.

\bibitem{xu-mirror}
L. Xu, A. Perrault, F. Fang, H. Chen, and M. Tambe, ``Robust Reinforcement
Learning Under Minimax Regret for Green Security,'' \emph{UAI 2021};
arXiv:2106.08413.

\bibitem{beukman-remidi}
M. Beukman, S. Coward, M. Matthews, M. Fellows, M. Jiang, M. Dennis, and
J. Foerster, ``Refining Minimax Regret for Unsupervised Environment Design,''
\emph{ICML 2024}; arXiv:2402.12284.

\bibitem{vinitsky-rap}
E. Vinitsky, Y. Du, K. Parvate, K. Jang, P. Abbeel, and A. Bayen, ``Robust
Reinforcement Learning using Adversarial Populations,'' arXiv:2008.01825, 2020.

%% ---------------------------------------------------------------
%% Team games, coordination, and search
%% ---------------------------------------------------------------

\bibitem{celli-gatti}
A. Celli and N. Gatti, ``Computational Results for Extensive-Form Adversarial
Team Games,'' \emph{AAAI 2018}, 2018; arXiv:1711.06930.

\bibitem{mcaleer-teampsro}
S. McAleer, G. Farina, G. Zhou, M. Wang, Y. Yang, and T. Sandholm, ``Team-PSRO
for Learning Approximate TMECor in Large Team Games via Cooperative
Reinforcement Learning,'' \emph{NeurIPS 2023}, 2023.

\bibitem{liu-hpsro}
N. Liu, M. Wang, X. Wang, W. Zhang, Y. Yang, Y. Zhang, B. An, and Y. Wen,
``Computing Ex Ante Equilibrium in Heterogeneous Zero-Sum Team Games,''
arXiv:2410.01575, 2024.

\bibitem{anagnostides-teamzero}
I. Anagnostides, I. Panageas, T. Sandholm, and J. Yan, ``The Computational
Complexity of Team Zero-Sum Games,'' arXiv:2606.16139, 2026.

\bibitem{bard-hanabi}
N. Bard, J. N. Foerster, S. Chandar, N. Burch, M. Lanctot, H. F. Song,
E. Parisotto, V. Dumoulin, S. Moitra, E. Hughes, I. Dunning, S. Mourad,
H. Larochelle, M. G. Bellemare, and M. Bowling, ``The Hanabi Challenge: A New
Frontier for AI Research,'' \emph{Artificial Intelligence} 280, 2020;
arXiv:1902.00506.

\bibitem{hu-otherplay}
H. Hu, A. Lerer, A. Peysakhovich, and J. N. Foerster, ```Other-Play' for
Zero-Shot Coordination,'' \emph{ICML 2020}, 2020; arXiv:2003.02979.

\bibitem{hu-obl}
H. Hu, A. Lerer, B. Cui, L. Pineda, N. Brown, and J. N. Foerster, ``Off-Belief
Learning,'' \emph{ICML 2021}, 2021; arXiv:2103.04000.

%% ---------------------------------------------------------------
%% Inference, counting, and evaluation
%% ---------------------------------------------------------------

\bibitem{ginsberg-gib}
M. L. Ginsberg, ``GIB: Imperfect Information in a Computationally Challenging
Game,'' \emph{Journal of Artificial Intelligence Research} 14:303--358, 2001.

\bibitem{frank-basin-ai}
I. Frank and D. Basin, ``Search in Games with Incomplete Information: A Case
Study Using Bridge Card Play,'' \emph{Artificial Intelligence}
100(1--2):87--123, 1998.

\bibitem{buro-kermit}
M. Buro, J. R. Long, T. Furtak, and N. R. Sturtevant, ``Improving State
Evaluation, Inference, and Search in Trick-Based Card Games,'' \emph{IJCAI-09},
pp.~1407--1413, 2009.

\bibitem{solinas-supervised}
C. Solinas, D. Rebstock, and M. Buro, ``Improving Search with Supervised
Learning in Trick-Based Card Games,'' \emph{AAAI-19}, 2019; arXiv:1903.09604.

\bibitem{rebstock-inference}
D. Rebstock, C. Solinas, M. Buro, and N. R. Sturtevant, ``Policy Based Inference
in Trick-Taking Card Games,'' \emph{IEEE Conference on Games 2019}, 2019;
arXiv:1905.10911.

\bibitem{lsw}
N. Linial, A. Samorodnitsky, and A. Wigderson, ``A Deterministic Strongly
Polynomial Algorithm for Matrix Scaling and Approximate Permanents,''
\emph{Combinatorica} 20(4):545--568, 2000.

\bibitem{cryan-dyer}
M. Cryan and M. Dyer, ``A Polynomial-Time Algorithm to Approximately Count
Contingency Tables When the Number of Rows Is Constant,'' \emph{Journal of
Computer and System Sciences} 67(2):291--310, 2003.

\bibitem{lisy-lbr}
V. Lis\'y and M. Bowling, ``Equilibrium Approximation Quality of Current
No-Limit Poker Bots,'' arXiv:1612.07547, 2016.

\bibitem{jordan-rleval}
S. M. Jordan, Y. Chandak, D. Cohen, M. Zhang, and P. S. Thomas, ``Evaluating the
Performance of Reinforcement Learning Algorithms,'' \emph{ICML 2020},
PMLR 119:4962--4973, 2020.

\end{thebibliography}
%% --- end sections_v05/bibliography ---

\end{document}
